\documentclass[../main.tex]{subfiles}
\begin{document}
%==========================================TIÊU ĐỀ TẬP========================================
\begin{table}[h]
  \begin{adjustwidth}{-1cm}{}
    \begin{tabular}{>{\centering\arraybackslash}p{6cm}|>{\raggedright\arraybackslash}p{12.5cm}}
      \multirow{5}{6cm}
      % Cột bên trái: ÉPISODE và số tập
      {\\[-40pt]
        \flaregothic\fontsize{20pt}{20pt}\selectfont \centering
        {\contour{errorfive}{\textcolor{black}{ÉPISODE}}}\color{black}\\
        \burbank\fontsize{72pt}{72pt}\selectfont
        \includegraphics[height=3.12359cm]{../../../../Image/Book?/number/num30.png}
        \\[18pt]
        % \flaregothic\fontsize{14pt}{14pt}\selectfont
        % \color{epattr}BIENVENUE, \uline{\color{Banpire} Mel} !
        % \flaregothic\fontsize{18pt}{18pt}\selectfont
        % \color{epattr}ÉPISODE PERDU
        \color{black}
      }
       &
      % Dòng thứ nhất: tên tập tiếng Nhật
      {\fotskip\fontsize{18pt}{18pt}\selectfont \ruby{青}{あお}\ruby{上}{がみ}\ruby{高}{こう}\ruby{校}{こう}の\ruby{運}{うん}\ruby{動}{どう}\ruby{会}{かい}}
      \\[6pt]
      % Dòng thứ hai: tên tập tiếng Anh
       & {\cafeteria\fontsize{18pt}{18pt}\selectfont Sport Festival}                                                                                                                                                                                          \\[4pt]
      % Dòng thứ ba: tên tập tiếng Việt
       & {\vnmsans\fontsize{18pt}{18pt}\selectfont Hội thao trường học}                                                                                                                                                                                       \\[8pt]
      % Dòng thứ tư: nhân vật xuất hiện
       & {\caviar{\fontsize{14pt}{14pt}\selectfont Nhân vật: \emp{sora} \emp{miko} \emp{roboco} \emp{haato} \emp{subaru} \emp{choco} \emp{ayame} \emp{aqua} \emp{shion} \emp{noel} \emp{pekora} \emp{kanata} \emp{towa} \emp{watame} \emp{lamy} \emp{botan}}}
      % \\[6pt]
      % Dòng cuối cùng: Thời gian diễn ra của tập (thường sẽ là ngày phát hành)
      % & {\continuum\fontsize{16pt}{16pt}\selectfont Ngày phát hành: 11/06/2022}\\[8pt]
    \end{tabular}
  \end{adjustwidth}
\end{table}
%===========================================NỘI DUNG==========================================
\color{black}

Tròn một tuần kể từ đêm nhóm bạn bước vào khu rừng sau thị trấn. Dòng chảy quen thuộc của thời gian lại mài mòn những ký ức lạ lùng, khiến chúng trở nên tròn trịa, dễ chấp nhận hơn.

Buổi sáng nào cũng thế, tiếng chuông trường Aogami vang lên đúng giờ, giảng đường lại ngập mùi phấn và bụi bảng, hành lang rộn rã tiếng cười đùa sau mỗi tiếng ra chơi, và tiếng bước chân học sinh tràn ngập hành lang mỗi khi tiết học kết thúc. Nhịp sống học đường cứ thế lặp đi lặp lại, như một guồng quay vô hình cuốn phăng mọi nghi ngờ.

Những bài kiểm tra đột xuất, những bản phân công nhóm, những câu chuyện vụn vặt trong giờ ra chơi\dots\ tất cả quay cuồng đều đặn, đủ để không ai còn thời gian nhớ đến bóng trắng đêm rừng.

Cuối tuần, nhà Hikawa bỗng trở thành điểm hẹn không cần lời mời. Azuki gõ cửa với ly trà đá mua dọc đường; Shion xách theo Aku như cái bóng; Wata thì mang bánh tới chia sẻ với mọi người; hay Ayame thi thoảng tới chơi vì ở nhà không có gì để làm. Các bạn học trong lớp 1-C, ai rảnh thì tới, cứ thế đầy phòng khách. Trời tắt nắng, họ ngồi thành vòng, kể chuyện không đầu không cuối, để bình trà nguội dần cũng không ai buồn dọn.

Ở trường, các câu lạc bộ đã khôi phục lịch tập. Sân bóng lại rộ tiếng giày da, phòng âm nhạc lại lất phất nốt cao vút, bảng tin đầy poster tuyển thành viên, mọi thứ như chưa từng có đêm gió rít bên ngoài khung cửa.

Riêng Kuon, cứ sau giờ học, cậu thay đồng phục bằng tạp dề nhỏ. Cửa hàng tiện lợi trong trung tâm thị trấn thuê cậu sắp hàng, thu ngân, dọn rác\dots\ một công việc có phần nhẹ nhàng hơn so với việc làm tổng quản ở thế giới gốc. Tiền sinh hoạt mà ``bố mẹ'' để lại không thiếu, nhưng cậu chưa từng được hứa thêm lần nào.

Chẳng mấy chốc, câu chuyện về ``những linh hồn trong rừng'' chỉ còn được nhắc đến như một giai thoại nửa đùa nửa thật, một mẩu tin đồn thú vị để giết thời gian. Đối với đa số học sinh, đó chỉ là phiên bản mới của truyện ma trường học: nghe cho vui, rồi quên đi như cơn gió thoảng.

Nhưng với những người đã trực tiếp đối mặt với đêm hôm đó, mọi thứ không hề phai. Kuon, Kaigou, Suzuki và Sakura vẫn cảm thấy hơi lạnh len lỏi mỗi khi gió thổi qua khung cửa lớp, vẫn thấy bóng đèn lồng rung nhẹ trong đầu mỗi lần nhắm mắt. Shino cũng đã được cung cấp thông tin về những con ma trong rừng, kể cả về một bóng hình nhỏ nhắn, tóc tết dài, đôi mắt xanh lục như ngọc bích ẩn dưới tán lá sương khuya.

Họ nhớ tiếng gió rít qua lá như lời hát ru, nhớ những bóng nhỏ co ro giữa màn sương mỏng, và đặc biệt là đôi mắt xanh lục của cô bé kia, ánh nhìn in sâu vào võng mạc, không thể xóa. Những hình ảnh ấy không biến mất, chúng chỉ trốn sau lớp áo bình thường, chờ ngày được phơi bày trở lại.

Công việc bán thời gian của cậu Tổng có được sau giờ học không chỉ đơn thuần là kiếm thêm thu nhập. Cậu đứng sau quầy thu ngân, mắt lặng lẽ quan sát từng gương mặt khách qua đường, học cách thị trấn Aogami hoạt động theo từng nhịp đập chậm rãi ẩn sau vẻ ồn ào thường nhật. Mỗi lần bước ra khỏi cửa hàng tiện lợi, cậu cố tình chọn con đường vòng qua những con hẻm nhỏ, để đôi chân không dừng lại trước ánh đèn mờ ảo vẫn treo lơ lửng trong tâm trí.

Con người, nhịp sống, và những điều chỉ người trong cuộc mới thấy. Dù phố phường có tĩnh lặng đến đâu, những linh hồn kia vẫn hiện hữu: chúng trôi dạt dọc phố nhỏ, nấp góc công viên, nép vào bờ tường rêu phong. Phần lớn dân thị trấn bước qua như bước qua cơn gió, nhưng những ai từng lặp lại thời gian như cậu có thể cảm được hơi lạnh thoảng qua, thấy bóng chúng lướt nhanh nơi khúc cua.

Những linh hồn ấy không làm gì cả. Chúng chỉ lặng lẽ trôi dạt quanh thị trấn. Không tấn công. Không gây rối. Hoặc ít nhất\dots\ chưa có dấu hiệu nào cho thấy chúng muốn làm như vậy, mà chỉ lặng lẽ tồn tại như lời thì thầm kể lại rằng đêm ba tuần trước ở khu rừng chưa hề khép lại.

\ct

\begin{center}
  \uline{Một tuần sau.\\Lớp 1-C - trường Trung học Aogami.\\Giờ chủ nhiệm.}
\end{center}

Những màn sương sớm cuối xuân còn đọng trên khung cửa sổ lớp 1-C khi tiếng chuông đầu tuần ngân dài. Học sinh lũ lượt kéo ghế, tiếng gỗ ken két hòa vào nhịp bước của cô chủ nhiệm. Yuko đặt chồng giấy xuống bàn, mắt nâng lên, đảo qua từng khuôn mặt như đếm lại ba tuần yên ả vừa trôi.

``Được rồi, cả lớp chú ý một chút.''

Giọng cô nhỏ nhưng đủ cắt tiếng xì xầm. Trong khoảnh khắc, lớp học như chiếc bình thủy tinh được lắp nắp; chỉ còn nghe tiếng hít thở và tiếng sách khép nhẹ.

``Trường vừa gửi thông báo,'' Yuko khoanh tay, khóe môi cong lên, nụ cười của người sắp thả một pháo sáng vào đêm. ``Ba tuần nữa, trường chúng ta sẽ tổ chức hội thao.''

Câu nói vừa rơi xuống, lớp đã nổ tung. Tiếng gõ bàn, tiếng reo hò, tiếng ghế kéo loạt xoạt như sóng thủy triều dâng. Cô giáo không cất lớn giọng, chỉ giơ tờ danh sách, để mặc cơn sóng ấy lan từ bàn đầu tới bàn cuối, rồi tự khắc lắng xuống khi học trò nhận ra ánh mắt cô vẫn bình thản.

``Từ tuần này, lớp 1-C chúng ta sẽ bắt đầu lên kế hoạch cho hội thao.''

Yuko khẽ gõ bút lên bảng, tiếng ``tách'' nhỏ như mở màn cho mùa hè sôi động. Ngoài cửa sổ, lá phong đã chuyển màu xanh đậm, gió đưa hơi nóng ẩm của sân bóng bay vào lớp. Cô đảo mắt một vòng, đo lường độ nhiệt trong không khí, vừa đủ để học trò bồn chồn, chưa đủ để bùng nổ.

``Trước tiên, mỗi người phải chọn cho mình một môn thi đấu.''

Cô xoay người, viết nhanh như người điều khiển đèn giao thông. Bảng xanh dần được lấp đầy bởi những mũi tên và dấu gạch đầu dòng trong bụi phấn trắng:

\begin{itemize}
  \item Tamaire - ném bóng vào rổ treo trên cột.
  \item Kéo co - nam nữ riêng biệt, thắng bại bằng sức nặng của đôi chân.
  \item Bóng né - dodgeball, nơi những quả bóng cao su có thể viết lại tình bạn.
  \item Đấu vật - tatami nóng hơn cả đường chạy.
  \item Chạy vượt chướng ngại vật - hàng rào gỗ, lốp xe, và bể nước cuối cùng.
  \item Chạy tiếp sức - thẻ gỗ nhỏ, tay ướt mồ hôi, tiếng reo vang của khán đài.
\end{itemize}

Ở góc phải, cô thêm một dòng nhỏ hơn, viết nghiêng: ``Bóng rổ - đăng ký riêng.''

``Quy định cố định: tối thiểu một môn, không giới hạn tối đa. Cứ việc nhiệt tình nếu các em tin bản thân mình đủ sức.'' Yuko bổ sung thêm, giọng mềm mại nhưng có chút gì đó cứng cáp.

Tiếng xì xào bắt đầu rộ lên, như sóng gợn trên mặt hồ khi có người ném viên sỏi đầu tiên, lan từ bàn cuối lên bàn đầu, từng cụm học sinh cúi sát vào nhau, tay che miệng, mắt sáng quắc, vai rung lên từng hồi vì những câu bông đùa chưa kịp thành lời; có đứa gõ bút xuống mặt bàn nhịp nhàng như trống trận, có đứa lại đạp nhẹ vào chân bạn để giành lấy sự chú ý, trong khi ghế sắt kéo loạt xoạt, quyển vồ bị vô tình đẩy rơi xuống nền, tạo nên một bản giao hưởng lộn xộn nhưng rộn rã của tuổi học trò đang bị dòng thác tin tức cuốn phăng khỏi bờ thời gian.

``Em ạ!''

Bàn tay Yae vụt lên trước khi cô lên tiếng, phá vỡ tiếng xì xào như cái chớp nhỏ giữa trưa hè. Cô dựa vào thành ghế, đôi mắt nửa khép, dáng vẻ như đang đọc thực đơn món tráng miệng của quán trà mùa hè chứ không phải danh sách môn thi đấu.

``Em tham gia bất cứ môn nào cũng được.''

Câu nói nhẹ nhàng như lời hứa, không khoe khoang nhưng đủ khiến bốn phía im lặng một nhịp. Không một ai ngạc nhiên, chỉ có tiếng gõ bút của Honoka còn vang lên như làm điểm nhấn.

Yae đã quá quen với việc được gọi là ``vận động viên'' của khối năm nhất với khả năng thể thao có thể gọi là toàn diện. Chạy cự ly, nhảy xa, ném bóng rổ, leo dây—cô tung tẩy như thể thể thao là ngôn ngữ mẹ đẻ. Có tin đồn cô từng bắt kịp xe đạp của anh giao báo chỉ để\dots\ nhặt lại tờ báo bị gió cuốn.

``Thậm chí nếu cần người `chơi tất' thì em cũng tham gia ạ,'' cô tiếp tục, khóe môi nhếch lên, nửa đùa nửa thách. Tiếng cười của cô Yuko vang lên trước, rồi cả lớp mới dám cười theo, như sóng xô bờ.

``Được rồi, đội chạy tiếp sức nhận Shiromi-san trước tiên.'' Cô giáo quay lại bảng, phấn trắng kẻ dòng thẳng tắp bên cạnh tên Yae, một chữ ký khởi đầu cho mùa hè sắp bùng nổ.

``Cả em nữa ạ.''

Tiếng nói thứ hai vang lên thuộc về Shiina.

Cô gái đứng dậy, tay khoanh trước ngực, ánh mắt lạnh lùng như băng sương sớm, một vẻ bình thản quen thuộc đến mức khiến người ta không khỏi rùng mình. Giọng cô vang lên, ngắn gọn, rành mạch: ``Em chọn đấu vật.''

Không cần chờ phản hồi, cô tiếp tục, như thể chỉ đang đọc thêm một mục trong danh sách:

``\dots và bóng né.''

Không ai bất ngờ. Ở Aogami, Shiina là cái tên gắn liền với những câu chuyện ``một mình một ngựa''. Cô không cần lời giới thiệu, chỉ cần đứng lên, cả lớp đã hiểu: đó là người có thể dùng nắm đấm để kết thúc mọi tranh cãi. Số lần cô ``nói chuyện'' bằng tay chân nhiều đến mức hầu như không một ai có gan dám tiếp xúc gần với cô.

Vì thế, khi cô đăng ký hai môn thiên về tiếp xúc và va chạm, không ai dám cười nhạo, cũng chẳng ai hoài nghi. Họ chỉ lặng lẽ tưởng tượng cảnh đối thủ bị cô quật ngã trên tatami, hoặc bóng cao su bay vèo qua đầu với tốc độ của viên đạn.

Yuko không nói gì thêm, chỉ khẽ gật, rồi hướng về bảng xanh, ghi thêm hai dòng mới bên cạnh tên Shiina, như thể ghi nhận một điều hiển nhiên đã được viết từ trước.

``Em cũng muốn tham gia ạ.''

Giọng nói tiếp theo đến từ Kaigou.

Cô vẫn ngồi thẳng, hai tay đặt phẳng trên mặt bàn, giọng không cao nhưng đủ cắt qua tiếng xì xào: ``Em chọn hai môn kéo co và đấu vật ạ.''

Cả lớp im bặt một nhịp. Không ai quên buổi sáng thứ hai của cô - khi tiếng động như bom nổ tầng dưới vang lên, lúc mà tên trùm của bọn côn đồ bị nốc-ao với vết tường lún sâu chỉ với một đòn đấm. Hỏi ra mới biết, đó là do chính Kaigou làm ra, khi cô vội chạy từ nhà tới lớp cùng cô em gái vì ngủ dậy muộn. Từ đó, họ biết sức mạnh của cô nằm ở đâu, không ở bắp tay lộ rõ, mà ở cú đấm có thể dồn cả trọng lực xuống một điểm.

Chỉ riêng Kuon và Suzuki hiểu hơn thế. Họ đã nhìn thấy cô nhấc một chiếc hộp kim loại nặng nửa tấn lên, ắt cũng chỉ để tìm cách phá đám Game trong thế giới của ông, và theo lời của Kaigou, cô cũng chỉ mới sử dụng chưa tới một phần trăm sức mạnh của mình. Những ai may mắn - hay xui xẻo - chứng kiến cảnh ấy đều hiểu: trong các môn cần sức mạnh và sức bền, Kaigou gần như là lựa chọn hoàn hảo.

Yuko gật nhẹ, viết tên cô lên bảng, dòng phấn trắng nằm song song với Yae và Shiina như khẳng định bộ ba không thể trượt.

Cô giáo quay lại, phấn kêu ``tách'' lần nữa: ``Ba bạn Shromi-san, Fukami-san và Kaigou-san sẽ chạy tiếp sức ở phần cuối hội thao nhé.''

Cô nàng tóc ngắn cười toe, vỗ tay vào nhau như thể vừa bắt được bóng: ``Được thôi ạ, em sẽ cầm cột về!''

Cô nàng bờm sư tử gật đầu lạnh lùng: ``Tôi chạy đoạn giữa, giữ nhịp.''

Cô chị lớn nhà Hikawa chỉ khẽ ``ừm'', nhưng khóe môi nhếch lên một milimet, đủ để cả lớp biết cô đã sẵn sàng đặt cả khối năm nhất lên vai mình nếu cần.

Sau một hồi lặng thinh, Kaigou ngẩng lên, giọng đều đều: ``À, Yuko-sesei, em cũng muốn thử sức với môn chạy vượt chướng ngại vật ạ.''

Yuko không ngạc nhiên, chỉ khẽ cười: ``Được thôi, Kaigou-san. Cô sẽ ghi thêm cho em.'' Phấn trắng lướt nhẹ, kẻ thêm một dòng bên cạnh tên cô, xác nhận Kaigou sẽ tham gia cả bốn môn: kéo co, đấu vật, chạy tiếp sức và chạy vượt chướng ngại vật.

Cô chủ nhiệm quay lại, mắt lướt qua những hàng ghế của lớp, tìm những nhân tố tiếp theo. ``Còn ai muốn xung phong nữa không?''

Không khí vừa kịp lắng xuống thì một cánh tay thon dài vút lên, nhanh như cái chớp giữa trưa hè.

Cô nàng luôn đeo găng tay Yuki.

Cô ngả lưng vào thành ghế, đôi mắt vàng hơi nheo, giọng nhẹ nhưng đủ khiến bốn bề im lặng: ``Em muốn thử sức với chạy vượt chướng ngại vật ạ.''

Một thoáng gió luồn qua khe cửa, làm tóc cô bay lên vài sợi. Cô khẽ nhún vai, tiếp lời: ``\dots và chắc luôn cả chạy tiếp sức ạ.''

Yuko không hỏi lại, chỉ gật đầu, phấn trắng kẻ một đường thẳng tắp bên cạnh tên Yuki, một dấu chấm nhỏ khép lại bản giao hưởng của buổi sáng. ``Được rồi, Kaizuka-san. Em sẽ chạy cùng với ba bạn trước đó nhé?''

``Vâng ạ\~{}''

Trong khi Yuko tiếp tục tìm thêm những bạn khác, Kana đã nghiêng người qua cửa sổ. Cô khẽ huých nhẹ vào cánh tay Kanade, giọng thì thầm như sợ gió tan.

``Kanade, thử đi.''

Nàng thiên sứ giật mình, đôi mắt xanh-tím nhìn xuống bàn, ngón tay luồn vào lọn tóc trắng muốt. ``Ơ\dots\ tôi à?''

Nàng ác ma mỉm cười, nụ cười ấm như nắng cuối chiều len qua kẽ lá. ``Ừ. Bà hợp với tamaire đó. Tôi thấy bà thích chơi những trò nhẹ nhàng mà.''

Cô hạ giọng thêm một nấc, gần như thì thầm vào tai bạn: ``Với cả ấy, đó cũng là một cách để giúp cả lớp đấy.''

Kanade im lặng. Mắt cô lướt qua bảng phấn, dừng lại ở dòng chữ ``Tamaire'' in hàng trên cùng. Một hồi chuông xa vang nhẹ trong đầu, có lẽ là tiếng phong linh thuở nhỏ cô từng treo trước hiên nhà.

Cô gật đầu, tay trái nắm lấy cổ tay phải như tự trấn an. ``\dots vậy Yuko-sensei, em đăng ký tamaire ạ.''

Tiếng nói nhỏ như hạt bụi phấn rơi, nhưng đủ để Yuko quay lại. ``Amano-san nữa.''  Cô giáo khẽ mỉm cười, phấn trắng kẻ thêm một dòng song song bên dưới tên Yuki, một vì sao mới trong chòm sao tiếp sức sắp thắp sáng.

Cuối lớp, một tiếng ``hừ'' khẽ cất lên như mũi tên vô hình bay thẳng vào không khí đang nóng dần. Koishin ngả lưng, hai tay khoanh trước ngực, ánh mắt nheo lại đủ để cả hàng ghế sau lưng tự động khép miệng. Cô không đứng dậy, chỉ nhấc cằm, thả giọng vừa đủ cả lớp nghe:

``Em chơi bóng né.''

Bốn chữ ngắn gọn nhưng đầy gai, như thể cô vừa tuyên bố chiến trường riêng chứ không phải môn thể thao. Ai cũng biết Koishin ghét thua đến mức sẽ ăn vạ hoặc tìm đủ cách đến khi nào giành chiến thắng mới thôi. Môn bóng né, nơi những quả bóng cao su lao vun vút, chính là đấu trường mà phản xạ và ý chí của cô được phép đốt cháy không kiểm soát.

Cô khẽ nhếch môi, tiếp lời như thể đang ký vào bản hợp đồng máu: ``Ghi tên em vào, Yuko-sensei.''

Yuko không quay lại ngay. Cô chủ nhiệm khẽ nghiêng đầu, phấn trắng dừng giữa không trung một giây, rồi cô mới cất giọng đều đều: ``Được, Akaba-san. Cô biết là em không thích hơn thua ai, nhưng nhớ nhẹ tay thôi nhé.''

Những tiếng cười nhẹ vang lên, nhưng Koishin không buông tay. Cô đã sẵn sàng với trận chiến sống còn trước mắt.

Tiếp theo là Ayame. ``Em cũng muốn tham gia ạ, yo!''

``Được rồi, Onidoro-san. Em muốn môn thi nào?''

Ayame ngồi thư giãn, lọn tóc bạc-đỏ lấp lánh dưới đèn lớp. Cô chống cằm, mắt lim dim, như thể đang xem một bộ phim chứ không phải đang chọn môn thi đấu, rồi khóe môi bật lên nụ cười ranh ma: ``Em thử chạy vượt chướng ngại vật!''

Tiếng cười lẻ loi vang lên vài chỗ, ai cũng hiểu cô nàng này đang nghĩ gì. Cứ nhìn vẻ mặt giả ngây thơ ấy là biết: cô nàng đã tưởng tượng mình nhảy qua hàng rào, vấp vào lốp xe, rồi bật ngửa xuống bể nước cuối cùng; cảnh hỗn loạn ấy đối với cô chính là món quà vui nhất hội thao, mặc kệ luôn vị trí cao nhất trong hội học sinh của trường, nơi thể diện là thứ phải có.

Yuko ghi tên, nhưng chưa dứt lời đã nhướng mày: ``Mà, Onidoro-san này, em còn việc quan trọng hơn.''

``Sao vậy cô?''

Giọng cô chủ nhiệm nhẹ, nhưng cả lớp lập tức im phăng phắc. ``Em đang là chủ tịch hội học sinh đúng chứ? Vậy nên em sẽ đồng thời làm MC toàn trường cho hội thao năm nay.''

Ayame hếch vai, mắt đỏ tròn xoe: ``Em biết mà, sensei. He he.'' Cô cười khúc khích, tựa như một đứa trẻ được tặng thêm một vé vào khu trò chơi cảm giác mạnh: vừa được chạy, vừa được hét, vừa được cầm míc-rô ra lệnh cho cả trường nhảy theo nhịp của mình.

Yuko khẽ vỗ tay, phấn trắng còn vương trên ngón tay. Cô vừa ghi xong tên Ayame, bảng xanh đã bắt đầu dày đặc phấn trắng hơn. Nhưng ánh mắt cô không dừng ở đó; cô quay xuống lớp, như thể muốn đếm lại xem còn ai chưa góp mặt. Và rồi, tựa như có dây vô hình giật nhẹ, mọi đầu lông mày đều ngoái về phía trước mặt cô, nơi nam sinh duy nhất của lớp vẫn đang dựa lưng, mắt ngước lên trần nhà như đang suy nghĩ một điều gì đó xa xăm.

Kuon cảm nhận được làn sóng tĩnh lặng đổ dồn về phía mình. Cậu không giật mình, chỉ khẽ nhíu mày, đôi mắt đỏ hỏn lên như vừa phát hiện được mối nguy hiểm.

``\dots Mọi người nhìn gì tôi vậy?''

Hanabi từ bàn bên tròn mắt hỏi: ``Kuon-kun không có ý định tham gia môn nào à?'' Mitsuki ở bên dưới hỏi tiếp: ``Cậu định né tránh hội thao ư?''

Cậu Tổng thở dài, ngón tay gõ nhịp lên mặt bàn. ``Tôi chẳng né tránh gì đâu, nhưng mà\dots\'' cậu khựng lại, cười méo, ``thành thật mà nói, tôi thực sự rất tệ môn thể chất. Chắc cả lớp cũng biết rồi đấy.''

Tiếng cười khúc khích nổi lên vài chỗ. Có đứa còn huýt sáo, như thể vừa nghe một câu chuyện khôi hài trước giờ ra chơi. Đã ba tuần kể từ lúc cậu học chung với các cô gái của lớp, nhưng đến lúc này hầu như không một ai tin rằng một người con trai lại thực sự kém khoản thể dục thể thao, mặc dù số lượng người như Kuon không phải là hiếm.

Koishin khoanh tay, ánh mắt nheo lại, giọng có gì đó hơi chảnh chọe: ``Con trai mà đầu hàng từ vòng gửi xe à? Thua cả Kanade-san kìa.''

Kuon nhún vai, chẳng buồn đáp lại. Cậu đã quen với những lời trêu như thế từ thuở cấp hai: những buổi chạy bền về đích cuối cùng, những lần bóng rổ ném trúng vành\dots\ tai mình. Thậm chí, bóng ma học lại (dưới danh nghĩa học cải thiện) môn thể chất hồi đại học năm ngoái vẫn còn nguyên đó, ám ảnh cậu mãi. Nhưng lần này, cậu chưa kịp lên tiếng thanh minh, vì Shino đã đứng dậy trước.

``Koishin-san.'' Giọng cô mềm như khăn lụa, nhưng đủ khiến cả góc lớp im lặng. ``Không ai bắt buộc phải giỏi thể thao để chứng minh mình là con trai hay con gái cả.'' Cậu Tổng cảm thấy một luồng ấm chạy dọc sống lưng: có người hiểu cậu, ít ra là hiểu phần nào. Không khác nào những cô gái đã cùng sát cánh bên cậu khi là một tổng quản ở thế giới thực.

Sakura chống cằm, mắt xanh lá mạ liếc qua Koishin rồi dừng lại ở Kuon. ``Tôi cũng không thích thể dục,'' cô thì thầm đủ nghe, ``đã thế còn trốn tiết thể dục để ra sau dãy phòng thí nghiệm. Chuyện bình thường mà.'' Cô cười khẽ, như chia ngọt cho cậu bạn cùng khổ.

Yae định bật dậy, bắp tay săn chắc lấp lánh dưới ống tay áo sơ mi: ``Hay để tớ huấn luyện cậu ấy--''

Cô nàng tóc hồng giơ tay lên, ngăn lại. ``Không cần đâu, Shiromi-san. Như Shino-san nói, không phải ai cũng cần chạy nhanh để tồn tại.'' Cô quay sang Kuon, nháy mắt, ``Cậu có thể lo phần cổ vũ tinh thần. Đội lớp mình sẽ cần một người giữ lửa khán đài.''

Kuon bật cười, lần đầu tiên kể từ khi tiếng chuông reo. Cậu gật đầu, nhẹ như trút được gánh. Yuko khẽ thở dài, ánh mắt dừng lại nơi mà cậu con trai đang ngồi. Cô gọi thẳng tên cậu, và mọi ánh mắt lập tức  ' Cô hiểu mà.''

Cô dừng một nhịp, đủ để tiếng cười nhẹ lướt qua vài khuôn mặt rồi tắt ngúm.

``Nhưng đây là hoạt động bắt buộc. Em phải tham gia ít nhất một môn. Nếu không thì hạnh kiểm của em sẽ bị ảnh hưởng đấy.'' Lời yêu cầu của cô giáo đanh thép kia dường như không làm dịu được với tông giọng ngọt lịm đó.

Cậu con trai cứng đờ. Cậu quay lại bảng xanh, đôi mắt đảo nhanh như tìm lối thoát trên một bức tường trơn.

Kéo co? Không. Sức kéo của cậu không thể so bì nổi với cả Ryuuhou.

Đấu vật? Tuyệt đối không. Thân hình của cậu quá mỏng manh để chịu đựng nổi sức nặng cơ thể.

Chạy? Càng không nốt. Chưa nói đến việc nhảy vượt rào, chỉ chạy thôi là đã quá sức với cậu con trai. Những lần cậu cố hết sức chạy, nhưng đó cũng là từng đó lần cậu suýt ngất lịm vì quá mỏi mệt.

Cậu thực sự \emph{không có thế mạnh} nào cả.

Không khí lớp trở nên lúng túng, như chiếc áo may sai chỗ khuyết một cúc. Kaigou khẽ nghiêng người, nhìn cậu con trai với sự lo lắng: ``\textit{Mình không nghĩ là Kuon-kun sẽ thực sự tệ môn thể chất đến mức đó đấy\dots}'' Suzuki thì hai tay nắm chặt váy, đôi mắt hai màu nhìn cậu với một sự hy vọng đầy nhỏ nhoi từ người anh tự nhận của mình.

Tiếng ghế kêu ken két, vài người cúi xuống giấu nụ cười, vài người khác lại nhìn Kuon như nhìn một con vật bị mắc kẹt giữa đường ray. Một sự thương hại, lo lắng, nhưng không ai thổ lộ - vì hoặc là không dám, hoặc là giấu nhẹm đi.

``\dots Đâu, Kuon-kun có phản xạ tốt mà.''

Cậu con trai ngoái đầu. Hotaru cất tiếng, ánh mắt nâu dí vào cậu như đang nhắc nhở một kẻ lãng quên vũ khí của mình.

``Tuần trước, lũ con trai lớp bên chặn cậu trước cổng trường đấy thôi. Cậu lùi một bước, nghiêng người, né được ba cú đấm liên hoàn như chơi đàn đó.''

Tiếng xì xào lan ra. Vài người vỗ trán ``A, đúng rồi!'' Yuki gật nhẹ xác nhận, Uzuki tiếp lời: ``Bọn mình trên tầng hai nhìn thấy rõ mồn một, cậu ấy thậm chí còn tránh được một đòn bẩn từ chúng cơ mà.''

Koishin không lên tiếng, chỉ khoanh tay, mắt hơi nheo như thể đang nói: ``Tôi cũng thấy, đừng có mà chối''.

Kuon khẽ nhắm mắt. Trong đầu, cảnh tượng ấy tái hiện: buổi sáng đầu giờ, gạch sân rêu phong, bốn thằng to xác vây quanh; cú đấm đầu tiên lao tới, cậu lùi gót, cú thứ hai cắt ngang tai, cậu cúi đầu, cú thứ ba nhắm mặt, cậu xoay vai\dots\ tất cả chậm như phim quay chiếu, không một mảnh áo rách, không một giọt máu.

Cậu khoanh tay lại, suy nghĩ một hồi lâu, ngẫm lại những gì mà các bạn học đã chia sẻ. Rồi, cậu thở dài, tay đưa lên như người chịu trận. ``Được rồi. Vậy\dots\ em chọn tamaire.''

Lớp khựng nửa giây, rồi nổ tung: tiếng vỗ bàn, tiếng huýt sáo, tiếng reo ``Cuối cùng Kuon cũng tham gia!'' vang lên như trống khai hội.

``\dots và bóng né nữa ạ.''

Hai môn đỡ tốn sức nhất - cậu nghĩ vậy - mà vẫn đủ để không bị điểm kém hạnh kiểm.

Yuko mỉm cười, phấn trắng kẻ dòng cuối: ``Hikawa Kuon - tamaire và dodgeball''.

Shino thở phào, khóe môi nhẹ nhõm. Koishin hất cằm, mắt như bảo ``Đấy, có gì đâu mà lằng nhằng ghê thế.''. Hotaru giơ ngón tay cái lên, Kaigou vỗ vai cậu động viên, còn Suzuki thì chỉ cười tít, hai màu mắt lấp lánh như vừa được tặng kẹo.

\ct

Khi bảng phấn đã kín tên, Yuko gõ nhẹ hai lần lên bảng, kéo mọi ánh mắt trở lại.

``Ngoài thi đấu, hội thao còn có phần thi cổ vũ giữa các lớp. Điểm phần này sẽ ảnh hưởng trực tiếp đến thứ hạng chung cuộc.''

Lời của cô vừa dứt, tiếng xì xào lập tức rộ lên. Những ai chưa kịp ghi tên vào môn thi đấu nay đã có thêm lối thoát; những người đã ``chốt đơn'' cũng bắt đầu tính toán xem mình có thể góp sức thêm chỗ nào.

``Các em không thi đấu, hoặc thi đấu ít, sẽ tự động vào đội cổ vũ. Nhiệm vụ của các em là: khẩu hiệu, cờ, trống, pháo sáng và cả giọng hát nếu cần. Làm bất cứ điều gì để cổ vũ cho lớp mình một cách nhiệt huyết nhất.''

Yuko mỉm cười như mèo vừa nhìn thấy cá, chờ đợi phản ứng của cả lớp.

Suzuki giơ tay trước tiên. ``Em xin vào đội ạ.''
Sakura gật theo. ``Tôi cũng thế.''
Shino nhẹ nhàng đưa tay. Hotaru và vài cái tên khác lập tức được ghi vào danh sách mới bên cạnh bảng thi đấu.

Khi dòng tên cuối cùng khép lại, Yuko đảo mắt một vòng, giọng bỗng trầm xuống đầy uyển chuyển. ``Một đội cổ vũ không thể thiếu người cầm lái. Ai tự tin làm trưởng đoàn nào?''

Im lặng. Mọi người nhìn nhau, rồi đồng loạt quay về phía Kuon, người duy nhất vừa được ``tha bổng'' từ việc chọn những môn thể thao vừa nãy. Cậu giật mình, hai tay giơ lên như đầu hàng.

``Đừng nhìn tôi. Tôi chỉ hứa cổ vũ chứ không nói là sẽ chỉ huy đội đâu.''

Yuko nhếch mày. ``Kuon-san, em sẽ là phó trưởng đoàn kiêm mũi chỉ huy cho lớp nhé. Em là người đàn ông duy nhất trong lớp đấy, đừng để cả lớp thất vọng.''

Cậu Tổng định phản đối thì đã thấy Suzuki ngoái lại, mắt lấp lánh: ``Khi Onii-sama không thi đấu, em sẽ ở bên cạnh anh, cùng anh chỉ huy đội cổ vũ!''

Kaigou bật cười khẽ, lắc đầu như người chị nhìn đứa em út mới tập nói. ``Cứ như thể cả thế giới phải quay quanh cậu vậy, Kuon-kun ạ.''

Cô em gái không giận, chỉ cười tủm tỉm. ``Đúng vậy đó! Ta cùng nỗ lực phấn đấu nhé, Onii-sama!''

Yuko gõ bút, kết luận. ``Vậy là, Hikawa Kuon-san sẽ phụ trách sáng tạo và điều phối; Hikawa Suzuki-san hỗ trợ trực tiếp; còn Ukai-san\dots\ em sẽ làm phụ đoàn, chuyên lo hậu cần, bánh kẹo, nước suối, thuốc xịt họng. Em đồng ý không?''

Kaoru, vốn đang ngáp dài, giật mình ngồi thẳng khi bị gọi tên. Cô ngơ ngác nhìn cô, tự chỉ tay vào mình: ``Dạ? Em đấy ạ?''

``Cậu chứ ai,'' Mari nhắc. ``Nãy giờ cậu có nghe cô nói gì không đấy?''

``C-Có chứ! E-Em sẽ cố gắng hết sức!'' nàng tóc ngắn quay về phía Yuko, quýnh quáng đáp lại.

Yuko khẽ gật, rồi quay sang một bên. ``Hitsujikado-san, em sẽ cùng Ukai-san lo hậu cần nhé. Một mình em ấy chắc chắn là không đủ đâu.''

Wata giật mình, mái tóc xoăn tung bay như đám mây nhỏ. Cô vội đứng dậy, hai tay nắm lấy váy đồng phục, cúi đầu một cái nhẹ nhàng. ``Dạ! Em sẽ chuẩn bị nước uống và đồ ăn nhẹ đầy đủ ạ. Có cả bánh quy và trái cây sấy nữa!''

Cả lớp bật cười. Trong tiếng cười ấy, bản danh sách cổ vũ hoàn tất, và mùa hội thao, với những trận cầu máu lửa lẫn những khúc trống rộn rã, chính thức bước vào hồi khởi động.

Cuối cùng, còn lại phần ghi hình cho hội thao.

Uzuki giơ tay phía bên cửa sổ, đôi mắt cam long lanh sau mái tóc bờm thỏ. ``Em có máy quay ạ.''

Yuko gật đầu, không cần hỏi thêm. Cả trường đều biết cô nàng này là ``con mắt sống'' của câu lạc bộ Nhiếp ảnh.

``Mizuta-san nhận phần quay, được chứ?''

``Dạ, được ạ!''

Kuon nghe vậy cũng nhẹ nhàng bật dậy, giọng trầm ấm: ``Tôi có thể phụ một tay.''

Uzuki quay sang, khẽ nhíu mày. Cô nhớ buổi đầu tiên khi chụp ảnh lớp, chính đích thân cậu đã mượn máy ảnh của cô và căn chính như một thợ ảnh chuyên nghiệp và cho ra lò những bức ảnh đẹp mắt (ngoại trừ một vài bức ảnh cậu chụp lỗi do yếu tố ngoại cảnh.)

Nhưng đó là việc chụp ảnh. Quay phim lại là một lẽ khác. Kuon có khả năng ăn ảnh rất tốt, nhưng việc lấy nét để quay một thước phim cũng là một kỹ thuật không hề dễ dàng gì.

``Cậu chắc chứ?'' cô nàng hạ giọng, đôi tai cụp xuống tỏ vẻ sự thận trọng. ``Mình không muốn cuốn phim rung lắc đâu. Vả lại, chẳng phải một mình cậu đang ôm đồm quá nhiều việc hay sao?''

Kuon cười trừ, nhún vai điềm tĩnh đáp: ``Tôi biết vài kỹ thuật quay phim cơ bản đấy. Trước từng dựng cả phim mà.'' Cậu không hề nói dối, vì sự thật là cậu đã cùng A-chan quay những thước phim \textit{hologra}, nên việc này không quá khó với cậu. ``Với cả đừng lo cho tôi. Tôi quen với việc làm nhiều việc cùng lúc rồi.''

``Ừm\dots\ vậy được.'' Uzuki gật đầu, má ửng hồng nhẹ. ``Mình tin tưởng vào cậu đấy.''

Hai người quay lại chỗ ngồi, mở điện thoại ghi nhanh kế hoạch: Kuon phụ trách ghi toàn cảnh sân khấu, khán đài, hậu trường. Cô nàng bờm thỏ thì ``bám'' các môn Kuon thi đấu để lưu lại khoảnh khắc cậu ném bóng hoặc né đòn. Sau hội thao, cả hai sẽ cùng dựng clip ngắn, thêm nhạc nền và chữ chạy, gửi Yuko-sensei trước buổi tổng kết.

Cô bạn bật cười khúc khích, vỗ nhẹ vào vai Kuon. ``Vậy hội thao này, chúng ta có cả phim điện ảnh rồi nhé!''

Nhìn hai người họ bàn bạc với nhau, Shino đưa tay lên miệng, khẽ cười trừ. ``Kuon-kun, cậu định cáng đáng cả hội thao luôn sao? Hai môn thi đấu, chỉ huy cổ động, lại còn cầm máy quay\dots\ Có khi nào cậu tự biến thành đội quân một người không?'' Một vài người cạnh bên nghe vậy bật cười khúc khích.

Kaigou gật gù, mắt híp lại đầy thú vị. ``Làm ba vai cùng lúc\dots\ Quả không hổ danh `cậu Tổng' có khác!'' Cô vỗ vai Kuon, thêm một cái thở dài giả vờ. ``Nhưng mà nhớ giữ sức, đừng để đến lúc máy quay rung lắc vì cậu kiệt sức giữa chừng đấy.''

Ở bục giảng, Yuko đang khoanh tay trước ngực, đôi mắt nửa khép như mèo ngắm chuột. Khóe môi cô nhếch lên một chút: vừa tán thưởng, vừa cảm thấy buồn cười trước `đội quân một người' kia. Phấn trắng lăn nhẹ trên đầu ngón tay, cô khẽ lắc đầu, thầm nghĩ: ``Hikawa Kuon\dots\ Em làm cô bât ngờ hơn cô nghĩ đấy.''

\ct

Không khí lớp học sau đó chuyển sang một giai điệu khác, nhộn nhịp như chợ hè. Những cái tên khác cũng đã an bài cho sân cỏ nay lại được điều đi khắp phòng: băng rôn, loa kẹo, khăn ướp lạnh, chai nước, thậm chí cả màu mực dùng để vẽ cờ. Mỗi người tự nguyện hoặc bị ``hiến'' vào một nhiệm vụ, tiếng cười xen lẫn tiếng van xin, tiếng gõ bút, tiếng ghế kéo loạt soạt tạo thành một bản giao hưởng lộn xộn nhưng đầy sức sống.

Yuko đứng trước bảng, tay khoanh lại, mắt lướt qua từng dòng phấn trắng như người thợ đang kiểm tra lại bản vẽ cuối cùng. Dưới ánh đèn huỳnh quang, bảng tên giờ đã dày đặc: thi đấu, cổ vũ, quay phim, hậu cần, băng rôn, âm thanh, nước uống, khăn lạnh\dots\ Một mạng lưới được dệt kín, không chỉ bằng kế hoạch mà còn bằng tinh thần ``tập thể lớp 1-C chúng ta là một''. Cô khẽ gật, như thể vừa ghim chiếc cúc cuối cùng lên chiếc áo đồng phục rộng thùng thình của mùa hè. Hội thao chưa bắt đầu, nhưng trên bảng xanh kia, mùa hè đã hiện hình rõ nét với những cái tên, những nét phấn, và cả những ánh mắt đang rực lên chờ trận cầu đầu tiên của năm học.

Trong tiếng râm ran của lớp, Ayame đứng phắt dậy, ghế kêu ken két. Cô đập bàn một cái *RẦM*, phấn trắng bật tung như pháo.

``Lớp 1-C chúng ta sẽ không thua bất kỳ lớp nào!''

Một chất giọng lớn vang lên như loa phát thanh buổi sáng, đủ khiến cửa kính rung nhẹ.

``Không chỉ không thua,'' cô nàng nheo đôi mắt đỏ lại, khóe môi nhếch lên như mèo vừa nhìn thấy cá, ``chúng ta sẽ đánh bại cả khối trên! \textbf{Mọi người có đồng ý với tôi không nào!?}''

Một thoáng im lặng, rồi cả lớp nổ tung. Cô nàng tóc trắng giơ hai ngón tay thành hình chữ V trước ngực, quay nửa vòng như ca sĩ trên sân khấu.

``\textbf{Nghe đây!}'' cô hét to hơn, ``\textbf{chiều nay ai rảnh ra sân trường! Tập khẩu hiệu, tập trống, làm tất cả mọi thứ để khủng bố tinh thần đội đối thủ nào, yo!}''

Tiếng cười, tiếng huýt sáo, tiếng vỗ bàn hòa thành một bản giao hưởng ồn ào.

Yuko đứng bục giảng, khẽ nhếch mày, không ngăn lại.

Ayame nhảy lên ghế, tay chỉ thẳng lên trần nhà như chỉ huy trận hải chiến.

``Từ giờ, mỗi người chúng ta là một ngọn đuốc! Khi reo hò, lửa sẽ cháy đến tận khán đài!''

Cô nàng phóng xuống, ôm ngang vai Kanade và Kuon, kéo cả hai quay vòng. ``Tamaire, dodgeball, chạy tiếp sức\dots\ dù thi đấu hay cổ vũ, chúng ta đều phải ghi điểm tuyệt đối!''

``\textbf{O---le}!'' cả lớp đồng thanh, tiếng vang dội cả hành lang.

Ayame hất tóc bạc-đỏ, mắt long lanh như vừa uống cạn ly năng lượng. Hòa chung vói đó, những tiếng reo hưởng ứng ầm ĩ, như thể trống khai hội đã đánh ngay trong lớp học. Yuko khẽ lắc đầu với cô học sinh nghịch ngợm, nhưng không giấu được nụ cười: ``Onidoro-san, em thực sự sinh ra để đứng trước đám đông đấy.''

Wata đứng đằng sau, mắt dán vào Ayame đang hét vang giữa lớp, lòng thầm nghĩ: ``\textit{Cậu ấy lại bùng cháy rồi.}'' Mái tóc xoăn của cô khẽ rung theo nhịp tim, vừa tự hào, vừa nhói nhẹ: ``\textit{Không biết liệu mình có thể bắt kịp được cậu ấy không nhỉ?}''

\il

Ngoài hành lang, bóng in dài trên nền gạch. Một linh hồn bé gái hiện ra, không tiếng bước, chỉ như tấm màn thời gian vén lên. Đôi mắt xanh lục của cô bé nhìn thẳng qua khung kính, rồi dừng lại ở Hikawa Kuon.

``\dots''

Cô bé không nói gì, chỉ có làn gió khe cửa lay nhẹ hai bím tết từng nhịp rung rinh như đang thở.

Ánh nhìn ấy mỏng như lá thu, mà nặng như lời hẹn cũ. Và rồi, khi cô bé chớp mắt một cái, cô tan biến, để lại hành lang trống và tiếng chuông xa vọng như hồi chuông báo đêm.

Cậu Tổng vẫn ngồi yên, nhưng trong lòng cậu chợt vang lên tiếng lá rơi, một linh hồn vừa đi qua, và một lời hứa nào đó vẫn còn đang dở dang.

\dots Một cảm giác thân quen không hề nhẹ.

\il

\ct

Ngay sau tiết học định mệnh ấy, làn sóng tin tức về hội thao trường Aogami tràn ra khắp ngóc ngách như thể ai vừa mở van một đường ống áp lực. Các bức poster A3 in màu rực rỡ được dán chồng chéo: từ bảng tin chính giữa hành lang đến mặt kính trước thư viện, từ vách gạch bên nhà ăn tới cột đèn gần sân tập nằng nặng mùi cỏ mới cắt.

Dòng chữ in đậm nổi khối, bóng đỏ cam chói lòa:

\begin{center}
  \Large{「第84回青上高校の運動会」\\2021-2022年度\\(HỘI THAO TOÀN TRƯỜNG TRUNG HỌC AOGAMI LẦN THỨ 84\\Năm học 2021 - 2022)}
\end{center}

Phía dưới, bảng lịch thi đấu in chữ nhỏ, Một vài tiếng reo hò, tiếng trầm trồ và tiếng bàn tán xôn xao vang lên, bởi lẽ cuối tuần này, sân trường sẽ biến thành một võ đài khổng lồ với đủ thứ môn từ tamaire mềm mại đến đấu vật bụi bặm.

Giờ ra chơi, dòng người qua lại như bị dán chặt vào tường. Nhóm con trai chỉ trỏ, nhóm con gái che miệng cười, thầy cô đi ngang cũng khựng lại đọc lướt. Tiếng xì xào vang lên không dứt:

\begin{itemize}
  \item ``Năm nay có thêm môn leo tường kìa, nghe đồn thưởng gấp đôi điểm.''
  \item ``Bọn mình đăng kí chạy tiếp sức rồi, phải giữ vững ngôi vô địch khối!''
  \item ``Nè, lớp bên kia nghe đồn có Hikawa Kaigou cũng ra gì đấy, lớp mình cũng không thể thua được.''
\end{itemize}

Chỉ trong vòng buổi sáng, hội thao đã vượt khỏi tầm một sự kiện thể thao, trở thành thứ ngôn ngữ chung của cả trường: trong lớp học, trong căn-tin, trên tin nhắn điện thoại và cả ánh mắt long lanh chờ đợi của những học sinh đang tận hưởng thời điểm đẹp nhất của tuổi học trò.

Những ngày sau đó, sân trường như được thay máu. Tiếng còi thể dục vừa vang, các lớp trong khối ùa ra nắng, không còn vẻ lê thê ngày thường. Ai cũng mang theo một mục tiêu nhỏ: tamaire thì tung bóng cho đúng tâm, đua tiếp sức thì căn đúng vạch, bóng né thì tính quỹ đạo để không ăn đòn. Không khí tập thể bỗng chốc biến thành mùi cỏ non đầu hè, mồ hôi của sự nỗ lực và quan trọng hơn hết: tinh thần thắng thua.

Yae là điểm nóng đầu tiên và cũng là trung tâm của mọi tiết thể dục. Cô gái tóc xám ấy chạy vòng sân như con thoi không biết mỏi, từng bước chân đều đặn đến mức có thể đặt đồng hồ. Thỉnh thoảng cô dừng lại, hai chân trước sau sát vạch trắng, miệng đếm từng nhịp rồi bật vọt về phía trước. Giáo viên thể dục đứng khựng, hai tay khoanh, gật gù như vừa tìm thấy viên ngọc thô ráp chưa mài giũa.

Bên kia sân, Shiina hóa thành đốm trắng lướt trên nền xanh. Quả bóng cao su nào tới gần đều bị cô tóm bằng cú quay người vừa nhanh vừa mềm, rồi phóng ngược lại khiến đối phương chật vật. Tiếng bóng *BỐP* vào lòng bàn tay cô đều nhịp như trống jazz, khán giả tình cờ qua đường không kìm được vỗ tay khen ngẫu hứng.

Gần đó, Kaigou cùng bốn bạn nữ cầm dây thừng to bằng cổ tay. Cô chỉ nhếch mép cười khi tiếng còi vang, rồi bất ngờ dồn một lực ``nhẹ'', hai chân trụ xuống như cọc tiêu. Đối thủ lớp bên loạng choạng, cả hàng bị kéo trượt một bước trên cỏ. Một nam sinh khối trên nhìn cảnh tượng đó, không khỏi lẩm bẩm: ``Cô ấy có phải người không vậy?'' Không ai đáp, bởi họ cũng đang mở to mắt như thể vừa chứng kiến robot biết\dots\ giận dữ.

Cứ như thế, từng tập thể, từng cá nhân, tất cả đều tự thắt vào mình một dây đếm ngược: bốn ngày nữa, sân khấu sẽ mở màn. Gió thu xoáy qua từng đường pitch, mang theo mùi sắp thắng, sắp thua, và nhất là mùi của thanh xuân đang rực cháy.

Và, không chỉ có các vận động viên, đội cổ vũ cũng lập tức vào guồng.

Vào những buổi cuối giờ học, trên mái nhà thể chất hay góc sân sau trường trở thành sân khấu thu nhỏ cho các lớp. Suzuki vừa đếm ``một, hai'' vừa vỗ tay, giọng ngân dài như người dẫn dắt dàn đồng ca, trái ngược hòan toàn với tính cách thu gọn của cô.

``Saya-san, giơ cao cờ ngay khi mình hét `chiến', nhớ giữ thẳng góc 45° độ nha!''

Cô gái tóc vàng gật đầu, đôi mắt hổ phách lấp lánh quyết tâm: ``Ừm, mình biết rồi!''

``Nanase-san, trống nhỏ đánh ba nhịp nhanh rồi dừng một, tạo điểm nhấn cho khẩu hiệu.''

Cô gái tóc xanh vỗ nhẹ mặt trống căng, khẽ huýt sáo ra hiệu đã sẵn sàng. Dường như cô ấy không ngại bị Suzuki sai bảo, ``\textit{ắt bởi hai thứ: hướng tới chiến thắng của lớp và\dots\ bản thân cô em gái của Kaigou cũng dễ thương}'' (theo suy nghĩ của cô.)

``Aku-san, míc-rô cứ để ở chế độ chờ; chờ tín hiệu của mình, cậu bật echo một giây rồi hét to `1-C' ngay sau đó.''

Aku đẩy kính râm (chẳng biết có ăn nhập hay không), tay cầm míc-rô rung nhẹ nhưng giọng vẫn đều: ``\textbf{R-Rõ, đội trưởng!}''

Suzuki thở sâu, giơ hai tay lên cao, mắt nhắm nửa như thể cả dàn nhạc đang chờ cây gậy của cô. ``Một\dots\ hai\dots\ \textbf{chiến!}''

Trong khi đó, Kaoru và một số bạn học khác cắm cúi trang trí tấm bạt, biến những dòng chữ ``1-C chiến thắng'' thành bức tranh màu nổi ba chiều, với sự hỗ trợ của Mari.

Wata lặng lẽ đặt chiếc giỏ cói xuống góc sân, lần lượt xếp từng hộp cơm tam giác, chai nước ion và khăn ướp lạnh thành từng dãy thẳng tắp. Cô khẽ đếm thầm: ``Mười chai lớn, mười chai nhỏ, thêm bốn chai dự phòng cho đội bóng né.'' Xong, cô rút từ túi áo một miếng giấy nhỏ, ghi chú bằng bút chì từ mỗi chai nước, mặt mỉm cười như thể đang trang trí bữa tiệc cho người thân.

Mỗi ngày, họ thêm một nhịp mới: băng rôn được kéo căng rồi thả lỏng để tập ``bật tung'' đúng khoảnh khắc reo hò, hay việc cờ lớp phải xoay tròn ba vòng trước khi dừng lại giữa không trung. Tiếng cười vỡ òa khi ai đó vấp chân, tiếng vỗ tay dồn dập khi khẩu hiệu vang đúng nhịp, và mùi cỏ non hoà mồ hôi làm nên mùi ``mùa hội'' đặc quánh.

Kaoru thậm chí còn góp ý: ``Khẩu hiệu không chỉ to, phải có điểm nhấn để khối trên phải giật mình!''; Suzuki đáp: ``Vậy thì chỗ `chiến thắng' mình kéo dài âm, khuấy động tinh thần cho lớp lên!'' và thế là cả nhóm lại hét thử, chim sẻ trên mái nhà giật mình bay tán loạn.

Khi hoàng hôn buông xuống, bóng họ in trên tường gạch như những nhân vật bước ra từ phim học đường: người cầm cờ, kẻ giật băng, người ôm trống nhỏ, kẻ vẫn còn khen tay mình vỗ đỏ ửng. Họ không chỉ tập cổ vũ, mà còn đang tập cho một mùa hè sắp nổ tung, để ngày thi đấu đến, tiếng reo của 1-C vang xa nhất trường, và để sau này nhớ lại, mỗi người đều biết mình từng là một phần không thể thiếu của bản giao hưởng ấy.

Giữa muôn vàn hoạt động hướng đến ngày hội thao, Kuon như người đánh trống hai đầu: vừa phải tập luyện, vừa phải lo hậu cần.

Buổi sáng nắng gắt, cậu khoác đồng phục thể dục ra sân, cùng Shiina và Koishin cùng các bạn học trong đội tập môn bóng né. Cú ném bất ngờ của Koishin tới tấp, cậu chỉ khẽ nghiêng vai, bóng vụt qua tai, tiếng cười trời cho vang lên như nhạc nền. Hotaru đứng ngoài gật gù: ``Phản xạ này đúng là vũ khí bẩm sinh.''

Hết giờ tập, trong khi bạn bè ngồi bệt uống nước, cậu lại tất bật chạy lên khán đài, giúp Suzuki điểm danh đội cổ vũ. Cô bé hổ phách đếm nhịp, cậu gõ bút ghi chú, hai người như dàn nhạc không cần dàn dựng.

Giờ tan học, khi trường học đóng cổng, cậu còn vội vã đạp xe xuống phố, kịp giờ ca tại tiệm tiện lợi, nơi cậu đã xin làm việc ba tuần trước đó. Thấy cậu thở hổn hển đầy mệt mỏi, ông chủ quán tỏ vẻ cảm thông, vỗ vai: ``Xem chừng căng đấy nhể? Chắc là hội thao hả?'' Kuon cười xin lỗi, xoay ca muộn, dậy sớm hôm sau để bù giờ.

Ngày nào cũng vậy, sáng hôm sau, cậu vẫn có mặt ở sân tập, mắt hơi thâm vì thiếu ngủ nhưng nụ cười chưa tắt. Có người hỏi: ``Cậu Hikawa-san không mệt sao?'', chỉ để nhận cái lắc đầu: ``Tôi còn phải giữ lửa cho cả lớp mà.'' Dưới ánh đèn huỳnh quang của cửa hàng đêm, trên sân cỏ ban mai, bóng cậu in lên hai thế giới, hai vai trò, hai lời hứa, và cậu không muốn đánh mất bất kỳ bên nào.

Song song với việc luyện tập của học sinh, công tác chuẩn bị hội thao của nhà trường cũng bắt đầu. Cả thấy và trò kéo nhau ra sân, dựng khung sắt, căng dây, vặn ốc, tiếng ken két vang lên như nhạc đệm cho bài khởi động của mùa hội.

Cột cờ nhỏ xíu được cắm thành hàng, mỗi màu một lớp, rung rinh trước gió như đang khoe sắc. Băng rôn lớp 1-C được mấy bạn hội học sinh giật lên giữa khán đài tạm, phất phơ như lá cờ hiệu ngầm báo: ``Khu vực này đã có chủ''.

Ở trung tâm sân, loa to như thùng rượu được gài dây, thử tiếng *XẸT XẸT* rồi bật ngay bản nhạc disco cũ lạc quẻ. Vạch trắng chạy tiếp sức được kẻ lại, thẳng tắp như ông thầy toán cầm thước kẻ bài.

Góc sân bên phải, mấy anh khối trên đang thử đường hầm vải, đứa nào cũng lóp ngóp chui ra với mái tóc rối như tổ quạ, cười ha hả. Lưới leo và thanh nhảy thấp được xếp thành đống như đồ chơi xấp hình khổng lồ, chờ ngày ``hạ màn''.

Gần nhà kho, cô giáo thể dục ngồi xổm, đếm bóng cao su từng quả, vừa đếm vừa lẩm nhẩm: ``Bóng cho tamaire còn thiếu ba quả, mà bóng né lại thừa năm\dots\ Thế là thế nào?'' rồi điều động các thầy cô gọi các bạn tập lại. Thỉnh thoảng, Ayame - với vai trò là hội trưởng hội học sinh -  lướt qua, cầm bảng ghi chép, mắt lia như chỉ huy trưởng, gật gù rồi phóng đi nơi khác, để lại phía sau mùi phấn và mồ hôi đầu hè.

\ct

Ba ngày trước giờ G, khắp hành lang trường Aogami bỗng vang lên tiếng loa phát thanh học đường rè rè: ``Các lớp chú ý! Sau giờ học hôm nay, toàn trường tập trung sân sau để phân đội và tập cổ vũ chung. Vắng mặt phạt lao động công ích!''

Tiếng ken két của cửa sổ đẩy mạnh, gió hè đầu mùa ùa vào lớp 1-C như thúc giục. Mọi người đồng loạt ngẩng khỏi bản vẽ cờ, khỏi bản nhạc trống chưa thuộc, khỏi mớ dây ruy-băng vừa tháo ra. Ayame đập bàn, mắt đỏ rực lên: ``Nghe chưa! Giờ không chỉ là tập lớp mình nữa, mà là tập khối! Đội Đỏ - đội Trắng, ai cũng phải biết mình đứng về phe nào!''

Tiếng ghế kéo loạt soạt. Kuon vừa cài máy quay vào balo, vừa kịp ném về phía Suzuki một ánh mắt nửa lo nửa hào hứng. Cô em gái hai màu mắt nắm chặt cờ lớp, gật nhẹ: ``Onii-sama, em đã chuẩn bị khẩu hiệu cho đội mình và lớp mình rồi.''

Cuối giờ, dòng người như thác đổ xuống sân sau. Nắng xiên qua mái tôn tạm, in hai vệt sáng đỏ - trắng rực rỡ. Giữa sân, thầy Hiệu phó cầm loa, giọng vang như còi tàu: ``Nghe hiệu lệnh của thầy! Các lớp sẽ tự chọn đại diện màu! Các lớp xếp thành hình quả chuông, chỗ đứng càng ngoài càng tốt!''

Tiếng xì xào bùng lên. Mấy cô cậu khối trên loay hoay\dots\ rốt cuộc lớp 1-C đứng chỗ nào? Chưa ai kịp định hình, Ayame đã đẩy vai các bạn học từ lớp khác, sắp xếp hàng ngũ một cách chuyên nghiệp: ``Cậu cao, đứng ngoài cùng! Còn mấy cậu kia,'' cô chỉ về phía đám đông lớp bên, ``đứng sau lưng cậu, hát vang lên!''

Chưa đầy năm phút, sân sau biến thành bãi chiến trận sắc màu. Đội Đỏ mặc áo thể thao đỏ rực, đội Trắng quàng khăn trắng loe loét. Tiếng trống nhỏ *tùng tùng* vang lên, một thầy giáo thể dục giơ tay làm nhịp: ``Một - hai - ĐỎ!'' Cả khối hét lớn như sóng vỗ: ``\textbf{Đỏ! Đỏ! Đỏ!}'' Liền đó, khối bên cạnh đáp trả: ``\textbf{Trắng! Trắng! Trắng!}'' Tiếng vang dội vào nhau, gió thổi căng cờ, mặt cỏ rung rinh.

Ayame nhảy lên bậc gỗ, hai tay vẫn giữ míc-rô tự chế, hét vang: ``1-C! Giọng ca của chúng ta đâu?'' Suzuki phất cờ, Kaoru đánh trống thùng *bình bịch*, cả lớp hò reo: ``Vô - địch - 1-C!'' Tiếng reo chưa dứt, khối bên cạnh đã nhảy vào nhịp, hô to: ``\textbf{ĐỎ bay xa! ĐỎ thắng lớn!}'' - khiến lớp đối thủ giật mình, rồi cũng cất cao: ``\textbf{TRẮNG cóng lên! TRẮNG về đích!}'' Không ai biết mình đang hò reo cho màu nào nữa, chỉ biết phải hét to hơn đối phương.

Sau mười phút, thầy Hiệu phó lại gõ loa: ``Giờ chuyển sang tập đồng ca. Đội nào hát được cả khối, điểm cổ vũ nhân đôi!'' Một tiếng ``Á?'' đồng thanh vang lên. Cả trường như bị bật ngửa, rồi lập tức xì xào: ``Nhạc gì? Lời đâu?'' hay ``Không có bản phát! Tự chế!''

Ayame nhanh như cắt, quay sang lớp: ``Beatbox! Beatbox nền đi!'' Kuon bật cười, đập tay vào đùi *phách phách*. Shiina huýt sáo *vút*, Kaigou gõ thùng rỗng *cộp cộp*. Chỉ ba giây sau, cả khối Đỏ bắt nhịp, hát vang bài ca tự chế: ``\textbf{Đỏ! Đỏ! Đỏ! Chúng tôi là hỏa xa!}'' Khối Trắng không chịu thua, bật ngay bản ``\textbf{Trắng! Trắng! Trắng! Bão tuyết về đây!}'' Hai dàn hợp xướng học đường đối đầu, tiếng vỗ tay, tiếng gõ chai, tiếng cười như nổ tung trong không trung.

Mồ hôi, nắng, cỏ, cờ, khẩu hiệu, beatbox - tất cả hòa vào nhau thành một bản giao hưởng không tưởng. Không còn kẻ thù, không còn phe phái, chỉ còn nhịp đập chung của tuổi trẻ đang rực lửa. Và khi tiếng loa phát ra hiệu lệnh cuối để kết thúc buổi tập trung, cả sân bỗng im phăng phắc, rồi vỡ òa trong những tràng pháo tay không dứt.

Trong khoảnh khắc ấy, dù chưa ai tiết lộ, nhưng mỗi người đều biết: dù đứng bên nào, họ đã cùng nhóm lại thành một trái tim khổng lồ, chỉ chờ ngày để bùng nổ.

\ct

Hai tuần trôi qua một cách nhanh chóng, trường Trung học Aogami như được lột xác. Sân trường trước chỉ có tiếng giày thầm thì nay đã thành bản giao hưởng vang dội: tiếng bóng cao su nảy đều, tiếng reo hò thử nghiệm, tiếng cười nắng gắt. Các khối học tách thành từng đội quân nhỏ, mỗi nhóm ôm một mục tiêu riêng: tung bóng tamaire cho trúng rổ, ném dodgeball sao cho đối phương ``cháy áo'', hay chỉ đơn giản là chạy tiếp sức mà không để đồng đội vỡ mặt.

Trong hành lang, giấy A4 dán chồng chéo: bản vẽ cờ, bảng màu, câu khẩu hiệu được gạch xóa đến nát đầu bút. Có lớp còn giấu như bảo vật: khẩu hiệu phải đến phút xuất quân mới chịu lộ diện, sợ đối thủ đánh cắp mất chiêu. Trên bảng tin mỗi lớp học, những lời động viên, những bảng xanh thông tin được trang trí được viết bằng bút dạ quang, băng keo, màu sáp, loé lên dưới ánh đèn huỳnh quang như mời gọi mọi ánh nhìn. Cả ngôi trường lúc này không khác nào như đang ở mội hội chợ.

Cảm giác háo hức không còn là làn gió nhẹ, nó đã thành cơn sóng thủy triều, lấp đầy từng ngóc ngách: căn-tin vang lên những tiếng ``trống'' bằng thìa nhựa, thư viện nghe thấy tiếng lục lọi tìm bản nhạc cổ vũ, những lớp học thì ánh đèn huỳnh quang bị thay bằng ánh đèn flash chụp ảnh thử cho khăn rằn.

Ngày hội thao còn chưa gõ cửa, nhưng trường Aogami đã nghiêng mình lao vào nhịp đập của một trận đấu lớn, nơi mỗi học sinh tự thấy mình là một mắt xích không thể thiếu của chiến thắng chung.

\begin{center}
  \uline{Hai tuần sau.\\Năm giờ sáng - trường Trung học Aogami.}
\end{center}

Ngày hội thao - ngày mà toàn trường Aogami mong chờ suốt nhiều tuần - cuối cùng đã gõ cửa.

Sương sớm còn vương ở khu vực trung tâm diễn ra các trận đấu, nhưng khắp nơi đã rộn rã tiếng bước chân. Kẻ loay hoay dây kéo co, người thử một lần cuối khẩu hiệu; có đứa chỉ muốn đi dạo cho vào guồng, nhưng tất cả đều mang theo mùi nắng mới tinh.

Kuon, Kaigou và Suzuki cũng đến sớm, với những toan tính riêng.

Kaigou xoay cổ tay, giật dây thừng ``bộp bộp'' như kiểm tra lại lực căng của chính mình. Thỉnh thoảng, cô còn chạy vài vòng ngắn để làm nóng cơ thể, từng bước chân chắc nịch và đều đặn, như thể đang tự điều chỉnh nhịp độ cho chính mình. Ở cuối khán đài tạm, bên dưới bóng cây bàng râm mát, một bé gái tóc tết đen nhánh lặng lẽ hiện ra, viền kimono nâu sẫn phất phơ theo làn gió sớm. Đôi mắt xanh lục trong veo hướng thẳng về cô nàng, hai nắm đấm nhỏ giơ ngang ngực, khẽ rung như cổ vũ thầm lặng.

Suzuki thì ở góc sân sau, giơ tay chỉnh đội hình cổ vũ. Cô vừa đếm ``một - hai'' vừa nghiêng đầu, mái tóc cột bím đung đưa theo nhịp. Kaoru ôm băng rôn chạy lòng vòng, thỉnh thoảng dừng lại kêu ``Cao thêm năm phân!'', giọng nghe như mẹ gà kêu con. Dù chỉ là buổi tập cuối, nhưng sự nghiêm túc hiện rõ trong từng ánh mắt. Và khi cô nàng mắt hai màu giơ tay ra hiệu, cô bé tóc tết lại lặng lẽ xuất hiện sau lưng đội cổ vũ, chỉ khẽ gật đầu như gửi lời động viên tới từng thành viên, rồi lại hóa vào làn sương ban mai trước khi ai kịp ngoảnh mặt quay lại.

Còn Kuon, sau khi khoác vội tạp dề máy quay, bắt đầu màn ``mở máy'' quen thuộc: lắp chân, căn góc, kiểm thử ánh sáng, thử âm. Cậu làm mọi thứ chậm đến mức một con rùa cũng phải ngáp, nhưng đó là kiểu cẩn thận của người từng bị lỗi cướp mất cảnh quay đẹp nhất đời.

``\dots?''

Đang xoay nhẹ ống kính, cậu bỗng chốc khựng lại.

Giữa sân, chỗ chỉ có vạch trắng và mặt cỏ còn đọng sương, có một bóng tròn nhỏ lơ lửng. Cô bé ấy lại đây: mái tóc tết sam đen nhánh, viền kimono nâu sẫn phất phơ, đôi mắt xanh lục như hồ kín đáy rừng. Chân cô bé không chạm cỏ, chỉ nghiêng người về trước, hai nắm đấm giơ ngang ngực, đúng điệu cổ vũ quen thuộc mà học sinh thường làm.

``Em ấy\dots\ đang động viên mình sao?'' tiếng cậu tuôn ra nhỏ xíu, như sợ gió nghe được. ``Nhưng\dots\ tại sao chứ?''

Từ đồng hồ đeo tay của cậu, một giọng nói cất lên: ``\eng{Có thể cô bé chỉ đơn giản thích cậu. Hoặc thích cảm giác được cổ vũ. Hoặc - lý do thứ ba - cậu đang nằm trong `route' của em ấy, và đây là flag auto-save.}''

Kuon khẽ ``hả?'' một tiếng, và rồi khi liếc lại\dots

\dots chỉ còn lại một bãi sân trống. Chỉ còn làn sương mỏng tan nhanh dưới nắng bình minh, và tiếng lá rơi khô khốc như ai đó vừa rút lui đầy chóng vánh.

Cậu nuốt khan, gắt gao nhìn vào khoảng không nơi đôi mắt xanh vừa lắng lại. ``\eng{\dots Ông nói như thể tôi đang ở trong visual novel vậy\dots}'' cậu thì thầo, rồi quay về máy quay, tay run nhẹ khi vặn ống kính lần cuối.

\ct

Không bao lâu, tiếng chuông báo hiệu giờ khai mạc vang lên. Cổng trường như bị cuốn phăng, học sinh từ mọi khối tuôn vào sân, lấp kín từng khoảng cỏ. Các lớp xếp thành hai khối vuông vức, các lớp đội Đỏ đối diện với đội Trắng giữa khoảng trống trung tâm, lá cờ hai màu rung nhịp nhàng trên cột cao như thách thức.

Lớp 1-C hòa vào màu đỏ, Kuon đứng thứ ba từ dưới lên, tay khoanh trước ngực, mắt lim dim ước lượng hướng gió. Kaigou và Suzuki đứng ở vị trí đầu, với cô em gái cầm cờ.

Không khí chợt nghẹt lại khi hiệu trưởng bước lên bục. Ông không cần míc-rô to; giọng trầm ấm vẫn đủ bao trùm cả sân. Những điều ông nói - như tinh thần thể thao, đoàn kết, ký ức thanh xuân - đều đã cũ kỹ tới mức phát chán với những khóa trên, nhưng vẫn mang nặng như lớp sương đọng trên cỏ, khiến nhiều người bất giác thẳng lưng.

Tiếp theo, Ayame trong bộ đồ thể dục thường thấy tiến lên. Cô không vội; từng bước giày in nhịp trên nền gỗ tạp, mái tóc bạc-đỏ khẽ lắc. Cô cầm míc-rô bằng hai ngón, như cầm roi phép.

``Thưa thầy cô và các bạn,'' giọng cô trong veo, ``với tư cách là hội trưởng hội học sinh, mình mong hai đội sẽ thi đấu hết mình, công bằng\dots\ và đừng ai gian lận, nhé?''

Chữ ``nhé'' được thả xuống nhẹ tênh, nhưng khóe môi nhếch lên đủ khiến khối Trắng ồ lên, khối Đỏ cười khúc khích. Một vài thầy cô lắc đầu, không giấu nụ cười.

Giọng điệu của cô nghe qua thì rất ``chuẩn mực'' và phù hợp của đại diện Hội học sinh. Nhưng, nhìn cách khóe môi hơi nhếch lên, cùng ánh mắt lấp lánh tinh nghịch, khiến không ít người bật cười khẽ.

Kuon đứng trong hàng, khoanh tay. Cậu biết rẩt rõ. Dù lời nói có vẻ trung lập, nhưng cô nàng chắc chắn đang nghiêng về đội Đỏ. Cụ thể hơn, tập thể lớp 1-C. Vì dù sao thì\dots\ cô cũng cùng lớp với cậu.

\ct

Khi tiếng vỗ tay vừa tắt, một khoảng lặng ngắn bao trùm khán đài. Ánh mắt mọi người đổ dồn về lối đi giữa sân, nơi hai đội trưởng bước lên. Họ không hẹn mà cùng giơ cao bàn tay phải, ngón tay dang rộng như thể đang cầm lấy bầu trời xanh trong buổi bình minh.

Lời thề thấm thoát vang lên, không dài dòng, chỉ vỏn vẹn ba câu. Giọng nam trầm, giọng nữ trong, hòa vào nhau thành một âm sắc duy nhất. Cuối mỗi câu, cả khối học sinh đồng thanh nhắc lại, tiếng vang dội xuống nền cỏ, lên mái hiên, rồi lan tận hàng cây bàng cuối sân.

Không có động tác rườm rà, không có cờ phướn hay nhạc nền.
Chỉ có tiếng gió nhẹ lùa qua lá, như thể chính thiên nhiên cũng đang chứng giám.
Trong chớp mắt ấy, hàng ngàn trái tim khác nhau bỗng đập cùng một nhịp:

\begin{center}
  ``\textbf{Chúng tôi sẽ thi đấu trung thực, hết mình, và tôn trọng đối thủ.}''
\end{center}

Nghi thức khép lại khi bàn tay hai đội trưởng chạm nhẹ vào nhau. Kèm theo đó là những tràng pháo tay và những tiếng reo hò từ khán giả - những người hậu phương đứng đằng sau để ủng hộ đội của mình. Một lời hứa vừa được gieo xuống đất, chờ cho tới thời điểm được bùng nổ.

\ct

Tiếng loa phát nhạc thể dục vang lên, toàn bộ học sinh của trường lập tức điều chỉnh hàng ngũ, lớp nào lớp ấy đều xếp thẳng tắp. Các khối học như những dải lụa đỏ-trắng được kéo căng, học sinh giãn cách đúng một bước chân, vừa đủ để cái bóng của người trước đè lên gót người sau.

Bài tập khởi động quen thuộc đến mức ai cũng có thể làm trong mơ: xoay cổ tay ba vòng, gập duỗi đầu gối, vươn vai nghe khớp răng rắc như hạt ngô nổ. Dưới ánh nắng bảy giờ, mồ hôi chưa kịp thành giọt đã bốc hơi, chỉ để lại một lớp mỏng muối trên mép áo.

Tiếng nhạc, tiếng hít thở, tiếng giày cao su  ma sát với cỏ - tất cả trộn thành một bản giao hưởng màu xanh lá, màu đỏ khối, màu trắng sân. Có đứa len lén quay đầu tìm bạn, có cô giáo nhịp chân theo phách, có anh lớp trên giả vờ kéo cơ nhưng thực ra đang nhìn lớp dưới để tìm xem có ai hợp gu không.

Khi khúc nhạc dạt lên nốt cao cuối, cả trường như một chiếc đồng hồ lớn vừa được lên dây cót xong: lò xo căng, kim chỉ đúng, chỉ chờ tiếng còi khai mạc để bùng nổ.
Và thế là mùa hội Aogami chính thức mở màn - không cần thêm một lời tuyên bố, chỉ cần ba nốt nhạc cuối cùng rơi xuống cỏ.

\ct

\begin{center}
  {\color{vol} Môn thi đầu tiên: \uline{Tamaire.}}
\end{center}

Phần thi đầu tiên của hội thao nhanh chóng được công bố.

Ayame bước ra giữa sân, míc-rô nhỏ gọn nằm gọn trong lòng bàn tay, dáng đi thong thả nhưng đủ khiến mọi ánh mắt dồn về phía trước. Nụ cười của cô lúc này không hoàn toàn thân thiện, nó vẫn dáng dấp gì đó như là thách thức, như thể cô vừa giấu một lá bài tẩy trong áo.

Giọng cô vang lên, rõ ràng, gọn ghẽ, không cần gõ loa: ``Luật thì trắng đen như phấn trên bảng. Mỗi đội có một đống bóng nhỏ xíu, đủ màu để khỏi nhầm. Trong vòng một phút, cứ tha hồ ném lên rổ. Mấy cái rổ này cao hơn đầu người một chút, nhưng không đến mức phải với tay giật cơ. Khi còi dứt, trọng tài sẽ đếm bóng bên trong lưới, không tính bóng nằm vạch hay lăn tung tóe. Đội nào nhiều hơn thì ăn điểm, đơn giản thế thôi.''

Cô khựng lại nửa nhịp, mắt liếc qua hai hàng học sinh đang xếp thành lối đi hẹp, rồi nhún vai như thể vừa kể một câu đố cũ: ``Nghe dễ như ăn kẹo, đúng không?''

Ở cuối hàng lớp 1-C, Kanade khẽ nghiêng về phía Kuon, môi gần chạm vành tai cậu, giọng thì thào đủ để hai người nghe: ``\dots Chẳng phải đây là trò dành cho trẻ con tiểu học sao?''

Kuon gật nhẹ, mắt vẫn dán vào chiếc rổ trống hoác phía trước: ``Tôi tưởng học cấp ba thì bỏ môn này rồi chứ.''

Yae đứng ngay phía trước hai người, nghe được câu bình luận, quay lại cười khẽ. Cô không phủ nhận, chỉ nhấc vai như thể đang khoác lên một chiếc áo mỏng: ``Đúng vậy, tamaire vốn là món quà dành cho lũ nhóc mẫu giáo. Nhưng đừng xem thường; nó là bài khởi động hoàn hảo để tạo không khí. Không cần sức mạnh, chỉ cần đồng điệu: ai ném, ai chuyền, ai hứng, ai la hét. Nếu cả đội là một chiếc máy đồng hồ, từng quả bóng sẽ là chiếc răng cưa ăn khớp. Và khi máy chạy trơn, thì ngay cả lớp khác giỏi cỡ đâu cũng phải hít khói.''

Kanade chớp mắt, như thể vừa thấy khía cạnh mới của trò chơi tưởng chừng vô hại: ``\dots Ra vậy. Nghe cũng\dots\ có lý đấy chứ.''

\ct

*\textbf{TOÉÉÉÉÉT!!}*

Tiếng còi từ trọng tài vang lên. Ngay khi đồng hồ một phút bắt đầu, lập tức cả hai đội lao vào nhặt bóng và ném về phía chiếc rổ ở giữa sân. Ai cũng cúi, vớ, ném; tiếng bóng cao su va vào lòng bàn tay rộn ràng như trống trận.

Yae bưgần như nhập cuộc ngay lập tức. Cô không chạy, không vội, chỉ lướt qua từng quả bóng như người thợ đan chọn sợi chỉ. Từng quả bóng được tung lên với góc độ ổn định, lực vừa đủ, tạo thành những đường cong gọn gàng trước khi rơi chính xác vào rổ. Cú ném của cô đều đặn như máy: cánh tay vung lên, cổ tay khẽ xoay, bóng vẽ cung mượt rồi lọt thỏm vào rổ. Không tiếng la hét, chỉ có tiếng *BỘP* của bóng bay vào rổ.

Thỉnh thoảng, cô khẽ nhắc nhở đồng đội: ``Đừng ném vội! Giữ nhịp!''

Kanade thì như con chim non vừa rời tổ. Cô vừa nhặt vừa loạng choạng, bóng văng đi tứ tung, có quả tạt sang trái, có quả nảy về sau, có quả lăn lông lốc dưới chân thầy trọng tài. Mỗi lần bóng trượt rổ, cô giật mình kêu ``á'' như người vừa giẫm phải ốc.

``Ơ-- Khoan-- từ từ! Mình phải làm thế nào đây!?''

Yae quay sang, mắt nhường lại một chút ánh sáng an ủi: ``Nhìn miệng rổ, không phải nhìn vào vành.'' Cô nàng thiên sứ gật, hít căng lồng ngực, lần sau bóng chạm vành rổ rồi rơi vào trong.

``Được rồi!'' cô reo như vừa bắt được cầu vồng.

Kuon đứng ngoài vòng xoáy, tay khoanh lại, mắt đong đưa theo từng quỹ đạo. Cậu quan sát đường gió bằng cách nhìn mái tóc Ayame khẽ bay, đoán độ nảy bằng cách nghe tiếng bóng chạm sân. Khi đã thấm, cậu mới cúi xuống nhặt một quả, xoay cổ tay, ném. Bóng đi chếch, hơi lệch trái, nhưng vẫn lọt lưới. Cậu nheo mắt, môi động không thành tiếng: ``Gió từ hướng Đông.''

Quả thứ hai cậu chỉnh thêm một gam nhỏ, bóng rơi thẳng tắp. Quả thứ ba, quả thứ tư\dots\ và cứ thế, mỗi lần ném là một dấu chấm hoàn hảo trên phông trắng sân cỏ. Ánh mắt cậu sáng lên như người vừa tìm ra đáp án cuối cùng của bài toán khó.

Yae vừa ném xong quả bóng, khẽ quay đầu. Ánh mắt cô lướt qua khoảng trống giữa hai hàng, dừng lại ở người con trai đang đứng nghiêng, môi lẩm nhẩm đếm góc, tay xoay nhẹ cổ tay rồi thả bóng.

``\dots Cậu tính cả hướng gió à?'' Tiếng cô như thở, chỉ đủ hai người nghe.

Kuon không ngẩng, vai khẽ lên xuống. ``Cho chắc thôi.''

Yae mím môi, khóe miệng cong lên thành một cung rất nhỏ. ``Không tệ. Nhưng mà nhanh lên nhé, còn 30 giây thôi đấy.''

``Biết rồi.''

Đúng lúc gió đổi chiều, một lực đẩy bất ngờ ập tới từ phía sau khiến Kuon loạng choạng bước lùi. Quả bóng cao su tuột khỏi lòng bàn tay, lăn tròn vài vòng rồi nằm lặng trên nền cỏ còn vương sương. Cậu xoay người, ánh mắt lạnh lùng đập vào khuôn mặt đã từng chặn mình ngay cổng trường buổi sáng thứ hai, là kẻ có bộ mặt cáo già và nụ cười luôn leo lên mép như thể thế giới nợ hắn ta.

``Ối, xin lỗi nhé\~{}'' tiếng nói kéo dài, đầy mỉa mai, vang lên như mảnh thủy tinh vỡ trong không khí. Kuon không đáp. Cậu chỉ nhìn hắn vài giây, đủ để nhớ lại cú đấm bất ngờ ngày ấy và lời đe dọa chưa kịp thực hiện.

Rồi, cậu gật nhẹ, như thể vừa lục lại một trang sách cũ: ``À, cái cậu định chơi bài móc cua sau lưng tôi này.''

Nụ cười của đối phương cứng đờ, chỉ kịp co rút trước khi biến thành tiếng cười khẩy. ``Lần này tao sẽ khiến mày bẽ mặt trước cả trường, giống như mày đã làm với thằng đại ca của bọn tao.''

Kuon nhún vai, không chút bận tâm. ``Ừ, chúc may mắn,'' giọng cậu đều đều, như thể vừa chúc một người xa lạ đạp xe đến trường thuận gió. Kẻ kia tặc lưỡi, ánh mắt vẫn nung nấu điều gì đó, nhưng bước chân đã quay đi, để lại phía sau mùi thuốc lá phảng phất và một khoảng lặng đầy mùi đồng cỏ nghiền nát.

Kuon cúi xuống, nhặt quả bóng, phủi nhẹ lớp bụi mỏng. Cậu trở lại vị trí, tiếp tục ném, mỗi lần thả tay là một dấu chấm tròn hoàn hảo trên nền xanh. Không ai kịp nhận ra cơn gió vừa xoáy qua, chỉ có bóng cậu vẫn thẳng tắp giữa sân, như thể chưa hề có va chạm.

Một lúc sau, Shiina lướt tới, mái tóc bạc phơ bay nhẹ. ``Gì thế?''

Kuon không quay lại, mắt vẫn dán vào rổ. ``Không có gì.'' Cậu ném thêm một quả nữa, bóng vẽ cung mềm rồi lọt thỏm. ``Chỉ là có kẻ muốn phá đám đội mình thôi.''

Cô nàng bờm Sư tử khẽ nhíu mày, nhưng im lặng. Cô đứng đó, cùng cậu nhìn theo quỹ đạo những quả bóng, để lại sau lưng kẻ đã khuất dần, một mối đe dọa còn lơ lửng, chưa kịp rơi xuống.

Mười giây cuối của đồng hồ bỗng nén lại thành một khoảnh khắc giòn tan. Nhịp tim hai đội đập loạn nhịp, bóng cao su bay lên như những vì sao lạc đường, vội vã tìm lấy đường parabol cuối cùng. Kanade ném bằng cả nỗi hối hận, Yae ném bằng cả sự tĩnh tại, Kuon ném bằng cả tính toán còn sót lại trong kẽ tay. Không ai còn nghe thấy tiếng gió, chỉ còn tiếng bóng chạm lưới như mưa sa trên mái tôn mỏng.

Tiếng còi dứt, mọi chuyển động bị cắt ngang như ai vừa giật dây diều. Không khí nặng trĩu, mắt mọi người dán vào chiếc rổ như thể đang đọc một bản trắc nghiệm số phận. Giáo viên cúi xuống đếm, từng quả bóng lăn ra như hạt giọt máu chậm rãi.

Ayame giơ míc-rô, khựng lại một nhịp, đủ để cả sân nghe thấy tiếng thở của chính mình.

``Và kết quả là\dots\ \textbf{đội Trắng chiến thắng!}''

Tiếng reo vỡ òa bên kia, tiếng thở dài bên này. Tên côn đồ ngoái lại, nụ cười như mảnh dao cũ kéo qua mặt cậ: ``Thấy chưa!? Đội Đỏ của chúng mày kiểu gì cũng sẽ thua thôi!''

Kuon chỉ đáp bằng ánh mắt mỏng: ``Chỉ mới thua một trận thôi. Mừng gì sớm thế.'' Yae gật nhẹ, như thể cả hai vừa ký vào một hiệp ước ngầm: đây mới chỉ là phần mở đầu của bản giao hưởng.

Kanade cúi đầu, mái tóc che kín nỗi hổ thẹn. Yae vỗ vai cô, nhẹ như gió thu thoáng qua: ``Cậu đã cố gắng hết sức rồi, không sao đâu.''

Trên bục, Ayame mỉm cười như người giữ lửa, phất tay như phất đi một mùa hè còn dài ở phía trước. ``Cảm ơn hai đội chơi! Tiếp theo sẽ là phần thi của khối tiếp theo, xin các lớp hãy chuẩn bị!''

Ở một nơi gần đó, Uzuki khép máy quay lại như khép một cuốn nhật ký bằng hình. Trong ống kính, từng khung hình đã được bao trọn: những quả bóng vừa rơi, tiếng reo vỡ, và cả khoảnh khắc một bàn tay đen đẩy nhẹ, đủ làm đổi hướng gió, đủ để một cậu bé loạng choạng giữa sân.

Cô không gọi tên điều đó. Ống kính đã thay lời: bằng chứng lặng lẽ, im như bụi trên băng ghế gỗ.

Ayame, từ bục chỉ huy, liếc về phía ấy chỉ bằng một góc mi mắt, nhanh như lưỡi cắt giấy. Ánh nhìn cô gái bờm thỏ và nữ chủ tịch hội học sinh gặp nhau giữa không trung, chạm nhau, rồi lướt qua, không tiếng động.

Cơn gió thổi qua, cuốn mùi cỏ non lên cao; hai người đứng đó, im lặng, đã thành đồng lõa bằng im lặng.

Chẳng cần nói gì thêm, họ cũng biết rất rõ ràng: họ sẽ không để vụ này chìm, nhất là từ một kẻ dám động tới cậu con trai đáng quý của lớp. Chỉ là, họ sẽ tìm một thời điểm thích hợp hơn để nói.

\ct

Ba khối vừa kết thúc cuộc thi, bảng tổng sắp liền được dán lên tấm bạt lớn. Mực đỏ chói in hàng chữ: đội Đỏ thắng. Tiếng reo vỡ òa như bong bóng xà phòng bị gió thổi bể, lan từ khán đài xuống thềm cỏ. Những cánh tay đỏ phất cờ, những bờ vai đỏ va vào nhau, còn khối Trắng lặng lẽ như bức tường vừa bị rút ổn. Họ không buồn, chỉ hít sâu rồi mỉm cười gượng, tinh thần thể thao được dệt bằng lớp sương mỏng trên khóe mắt.

Kuon lùi lại phía sau, dựa vào cột cờ nhỏ. Ánh mắt cậu lướt qua bảng điểm, không vui, không buồn, chỉ như kẻ đang đọc một câu thơ đã biết trước vần. ``Ít nhất thì các khối khác cũng gánh cho ta rồi,'' cậu thì thầm, giọng khàn vì cỏ và gió.

Yae đứng cạnh, khép cánh tay đẫm mồ hôi sau lưng. Cô cười khẽ, tiếng cười như hạt mưa rơi trên lá chuối. ``Đừng dựa dẫm vào mọi người quá chứ.'' Cô vừa nói vừa xoay cổ tay, khớp xương lách cách như hạt đậu rang, chuẩn bị cho môn kế tiếp.

Kanade thì gục cằm xuống, ngực phập phồng như con chim non vừa thoát khỏi bão. ``\dots Mình đi nghỉ chút\dots'' Giọng cô mỏng, dính vào không khí như lớp sương sớm. Yae gật, đưa tay vỗ nhẹ lên lưng cô, một cái chạm thay lời an ủi.

Bên khán đài tạm, Kana vén tà váy xuống, bước nhanh về phía Kanade. Mắt cô sáng lên niềm kiêu hãnh rõ mồn một. ``Kanade-chan, bà đã làm tốt rồi.'' Cô nắm lấy tay bạn, siết nhẹ, truyền qua lòng bàn tay một dòng nhiệt như muốn nói: dù bóng có lăn lông lốc, thì người cầm bóng vẫn đáng được vỗ tay. Kanade ngẩng lên, khóe mắt ươn ướt mồ hôi, nhưng nụ cười đã kịp nở, một bông hoa mùa hè vừa vượt qua đêm giông.

Kuon lững thững rời cột cờ, bóng dài trên nền cỏ như bóng cây thông đang bước. Cậu vừa đi, vừa cởi chiếc băng đô ướt sũng, vắt lên vai. Khi ngước lên, cậu thấy Suzuki đứng trước khán đài, tay giơ cao tấm khăn mùi xoa trắng tinh. Cô không nói gì, chỉ mỉm cười, nụ cười của người biết rõ cậu vừa trải qua những phút giây bị bóng đá vào lòng. Cô vẫy nhẹ, đưa khăn ra. Kuon tiến gần, đón lấy, lau vội vệt mồ hôi trên trán.

Sau đó, cậu quay về phía đội cổ vũ, nơi những chiếc trống nhỏ đã được dọn ra, nơi những dải ruy-băng đỏ còn vương gió. Suzuki bước sát bên, hai bóng dài song hành như hai dòng thơ chưa kịp xuống dòng. Họ không nói thêm, bởi tiếng reo của lớp 1-C đã vang lên, gọi họ về với guồng quay còn dài, chờ những người trẻ nắm lấy phách.

\begin{center}
  {\color{vol} Môn thi thứ hai: \uline{Kéo co.}}
\end{center}

Phần thi thứ hai nhanh chóng được triển khai ngay sau tamaire.

Ayame lại bước ra, míc-rô gọn trong lòng bàn tay, cô đứng giữa sân, nơi cỏ còn đọng sương, nơi hai đội đang nén hơi thở thành những dải băng căng cứng. Ánh mắt cô lướt qua họ, vừa đủ để mỗi người tưởng mình vừa được nhìn thấu, rồi lập tức bị thả lại vào khoảng lặng.

``Ba hiệp,'' cô nói, giơ ba ngón tay lên, giọng vẫn nhẹ tênh như lụa, ``đội nào thắng hai hiệp trước sẽ giành chiến thắng.'' Ngón tay cô chỉ về phía dây thừng, đợi người ta truyền lửa vào thân nó. Cô tiếp tục:``Chỉ việc kéo đối phương qua vạch. Không buông tay, không đá gót, nhớ giữ an toàn cho nhau.''

Cô dừng lại một nhịp, đủ cho gió lách qua kẽ áo, đủ cho tiếng tim ai đó lỡ nhịp. Nụ cười nhẹ như hơi rượu trên môi kẻ vừa thức dậy:  ``Chúc các bạn kéo vui vẻ\~{}''

Tiếng cuối cùng tan trong không, như mũi tên đã rời cung, chỉ còn dây rung nhẹ giữa hai hàng học sinh, những kẻ sắp biến thành hai bờ sóng, đâm vào nhau bằng cả trọng lượng tuổi trẻ.

*\textbf{TOÉÉÉÉT!}*

Tiếng còi giữa sân vừa điểm, hai khối đầu tiên đã bước vào khoảng cỏ như hai dải lụa đỏ-trắng được kéo căng đến mức rung mình. Không khí không ồn ào, chỉ lặng lẽ căng ra như mặt trống mới căng da, chờ một nhịp phách đầu tiên rơi xuống.

Các đội thi đấu nghiêm túc, dây thừng giữa sân trở thành ranh giới mong manh giữa hai thế giới: bên này là bước chân dồn xuống cỏ, bên kia là hơi thở gấp gáp đẩy ngược lên trời. Tiếng hò reo của khán giả không ào ạt, mà tuôn ra từng đợt, như thủy triều lên xuống theo từng xăng-ti-mét dây lùi về phía đối thủ.

Có trận kết thúc nhanh như cơn lốc: một lực kéo bất ngờ, vạch giữa sân bị nuốt chửng chỉ trong nháy mắt, để lại bên thua cuộc ngỡ ngàng như kẻ vừa bị giựt mất khăn tay. Có trận kéo dài đến hiệp ba: dây thừng rung lên từng hồi, mồ hôi rơi xuống cỏ thành những vệt sáng loang lổ, hai bên giằng co đến mức cả sân như nghiêng theo từng hơi thở dịch chuyển của vạch trắng giữa sân.

Nhưng dù nhanh hay chậm, dù áp đảo hay giằng co, tất cả vẫn nằm trong khuôn khổ của một cuộc thi thể thao bình thường: không có bất ngờ nào vượt ra ngoài dự đoán, không có khoảnh khắc nào khiến trái tim khán giả phải ồ lên. Chỉ là những bước chân dồn xuống cỏ, những bàn tay siết chặt dây, những tiếng reo vang lên rồi tắt lịm, như những nốt nhạc đã được viết sẵn trong một bản giao hưởng quen thuộc.

Đó là cho đến khi đến lượt khối của Kaigou.

Khi dây thừng còn nằm im trên cỏ, đội hình đã được xếp, nhưng khoảng trống giữa những bàn tay vẫn rung lên một nhịp lạ. Kaigou đứng ở cuối hàng, thân hình nhỏ bé như một cọng cỏ thép giữa bãi lau, mái tóc đen phủ kín gáy được buộc búi cao, đôi mắt xanh nước biển biểu hiện sự lạnh lùng. Cô là hạt cát duy nhất trong cốc thủy tinh đầy sỏi cuội, cùng với bảy nam sinh, bảy thân hình cao lớn, bảy bộ óc chưa kịp hiểu vì sao mình lại bị một cô gái dẫn dắt.

Một người trong số họ quay lại, nụ cười cố gắng như bật cửa sổ giữa đêm đông: ``Cố gắng nhé, đội mình--''

Tiếng ``Ừ'' của cô nàng rơi xuống đầy hời hợt như hòn đá lạnh, không vỡ, không vang, chỉ làm mặt hồ rung một vòng tròn rồi lặng. Cô chẳng thèm nghiêng đầu, mắt vẫn dán vào dây thừng như thể đang đọc một đoạn mật mã duy nhất trên thế gian. Nụ cười của chàng trai cứng lại, chưa kịp nở đã thành sương.

Không khí giữa họ như tấm lụa vừa bị cắt đứt chỉ, mềm mại nhưng không thể khâu lại.

Bên ngoài, Kuon dựa vào thành khán đài, hai tay khoanh trước ngực như ôm một bản giao hưởng dở dang. Cậu thở dài, ngán ngẩm với cô nàng: ``\dots Biết là ghét đàn ông rồi, nhưng ít ra cũng nên cho họ một chữ `cảm ơn' chứ.''

Suzuki đứng cạnh, môi cong như cánh cung nhẹ, mắt nhìn theo bóng Kaigou như người nhìn theo một ngôi sao xa, chỉ biết cười gượng: ``Onee-sama là thế đấy. Chị ấy cẩn thận tới mức đó mà\dots''

Nhưng điều khiến bãi cỏ như nghiêng mình không phải việc cô là cô gái duy nhất trong đội, mà là một bóng đen đứng cuối dãy dây, một khối cơ bắp được đắp bằng đất sét nung, cao đến mức ánh sáng phải vòng qua vai hắn mới tới được mặt đất. Dây thừng trong tay hắn trông nhợt nhạt như sợi chỉ, nhưng cũng chính nơi đó, đội bạn cảm thấy mình đã cột chắc cả một trời thắng lợi.

Danh xưng ``vô địch kéo co'' không còn là biệt danh, mà là một bản án treo lơ lửng trên đầu mọi đối thủ: mười người lớn gắn chặt vào đầu kia cũng chỉ bằng một mình hắn. Yae đứng ngoài vạch vôi, môi khép hờ như người vừa nếm phải trái bồ hòn: nó chát, nhưng không thể nhổ.

``Này, con nhóc lạnh lùng kia!'' tiếng hắn vang lên, đầy nhựa sống, đập vào tai Kaigou như búa tạ, ``Lần trước bà đấm thằng đại ca của tao một cú, giờ tao kéo cho bà biết tay. Cứ bám vào dây mà cầu nguyện đi!''

Không ai đáp. Gió chỉ kịp cuốn mùi cỏ non qua kẽ áo. Kaigou nhìn xuyên qua hắn như xuyên qua một tấm kính mờ; trong mắt cô, bóng hắn chỉ là một vết loang đậm trên bức tranh đã cũ. Cô khép miệng, thầm thì với chính mình: ``\hl{Thế quái nào mình lại phải đứng chung với đám cặn bã này chứ\dots}'' rồi quay đi, để lại sau lưng một khoảng trống lạnh lẽo, nơi cả tên khổng lồ cũng bất giác thu nhỏ lại thành hạt bụi.

*\textbf{TOÉÉÉÉT!}*

Tiếng còi vang lên, hiệp một bắt đầu.

Cả hai đội đồng loạt dồn chân xuống cỏ, dây thừng giật căng như dây đàn mới thay. Nhưng giữa nhịp giật ấy, Kaigou vẫn đứng thẳng, hai tay chỉ nắm lỏng dây, hầu như còn chẳng động đến. Ánh mắt cô trôi qua khoảng không phía trên đầu đối phương, hờ hững trước những gì xảy ra trước mắt.

Chẳng mấy chốc, lực kéo toàn đội lập tức mất cân bằng. Nam sinh phía trước giật ngược về sau, chân trượt trên cỏ, còn chưa kịp hét thì vạch trung tâm đã lùi qua khỏi mũi giày họ. Hiệp một kết thúc chỉ trong bảy giây, với chiến thắng thuộc về Đội Trắng.

``Cái gì vậy trời!?'' một nam sinh phía trước quay lại, mặt hiện rõ sự bất bình.

``Cậu đùa à!?''  một nam sinh khác lên tiếng.

``Tôi đổ mồ hôi ra đây nè!'' một cậu nam sinh cuối hàng quay lại, mặt đỏ bừng, dây vẫn còn rung rung trong tay.

Tên khổng lồ bên kia vạch cười hở cả hàm răng trắng nhởn, giọng vang như pháo: ``Hửm? Có phải con nhỏ kia sợ bị tao kéo bay không? Cứ việc bám chặt vào, đừng có mà khóc nhé!''

Đồng đội hắn huýt sáo:  ``Yên tâm, có `đại bàng' ở đây rồi, đội kia chỉ cần cầm dây cho đẹp mắt thôi!''

Kaigou chỉ liếc qua, khóe môi không nhúc nhích. Không khí như có ai siết chặt dây thừng giữa chừng, im lặng đến nghẹt thở. Cô dường như không để tâm tới những lời chế nhạo của đội bạn, ánh mắt vẫn chìm vào nơi xa xăm mà chính cô cũng không hề biết.

Ayame bước tới trước, míc-rô tắt, đưa cô về thực tại: ``Kaigou-san, cậu định để cả lớp mình thành trò cười cho khối bên kia à?''

``Gì?'' cô nàng mắt xanh căm căm, nhìn đồng đội cô với ánh mắt ghẻ lạnh. ``Cậu muốn tôi đồng hành với đám con trai dơ bẩn kia á?''

Suzuki đứng sát bên, nhỏ nhẹ hơn: ``Onee-sama, em biết chị ghét đàn ông, nhưng bảy người kia vẫn là đồng đội của chị đấy.''

``Ít nhất cũng phải có tinh thần đồng đội chứ,'' Kana nói.

Kaigou thở dài, mũi chân xoay vòng trên cỏ như đang tính góc gió. ``\dots Rồi rồi.'' Cô siết nhẹ dây, lần đầu tiên cánh tay lộ đường gân xanh.  Giọng cô vẫn lạnh, nhưng dây thừng đã nghe được tiếng lòng cô: nó sắp được kéo thật.

\ct

*\textbf{TOÉÉÉÉT!}*

Hiệp hai gọi tên bằng một hơi thở dài của còi, âm thanh như lưỡi cưa kéo qua mặt trống căng. Ở hai đầu sợi dây, mười bảy cái lưng cúi xuống cùng lúc, mười bảy đôi giày đóng cọc vào cỏ, mười bảy cái tim đập tràn vào nhau qua một sợi gai bện. Dây thừng giật căng, căng đến mức ánh nắng cũng phải trượt dài trên đó mà rơi xuống cỏ thành những hạt sáng lăn tăn. Không ai nhúc nhích, như thể cả hai bên đang giữ một con thú dữ trong lồng ngực, chờ nó gầm lên trước thì thôi.

Hai bên khán đài bỗng trở thành bờ biển. Tiếng reo không ào mà cuốn, không ồn mà sâu, cứ thế dâng lên từng đợt, vỗ vào lưng những người đang còng. Suzuki đứng trên bậc cao nhất, hai bàn tay nắm chặt như búa tạ không biết đánh vào đâu, chỉ biết gọi: ``\textbf{Onee-sama! Chị làm được mà!}'' Tiếng cô vừa dứt, Kaoru đã nâng lên thành điệp khúc, ``\textbf{Đội Đỏ cố lên! Đội Đỏ vô địch! 1-C bất diệt!}'' những âm tiết rơi xuống sân như hạt mưa đá, nảy lên vai người, rồi lăn vào khoảng trống giữa hai hàng, nơi dây thừng đang rung lên từng hồi như dây đàn chưa tìm được nốt thuộc về mình.

Touka đứng ngoài vòng xoáy, tay khoanh, mắt nheo lại thành khe sáng hẹp. Cô không hò reo như các bạn, chỉ quan sát, như người cầm thước đo giữa cơn bão. Và từ đây, cô nhận ra có gì đó không đúng. ``\dots Lạ thật.''

Shino nghiêng đầu, mái tóc nâu trượt xuống vai, hỏi khẽ: ``Sao vậy?''

Cô nàng Thần đồng không trả lời ngay, mà chỉ bằng cả khoảng lặng đang đặt giữa cô và Kaigou. ``Cậu ấy\dots\ gần như không cử động.''

Quả thật, Kaigou vẫn đứng đó, thân hình nhỏ bé như cọng cỏ thép giữa bãi lau. Hai bàn tay nàng chỉ buông lơi trên sợi dây, không siết, không kéo, không đẩy, như thể nàng đang cầm lấy một dòng sông bằng hai ngón tay. Thế mà đối phương - tám gã trai đang cúi rạp, tám cái lưng đẫm mồ hôi - không thể kéo nổi thêm một phân. Sợi dây căng ra giữa họ như một đường chân trời không thể dịch chuyển, và Kaigou chính là cái mỏ neo đã quên (hoặc không để tâm) mình đang cắm vào đá từ thuở nào.

Shino thở nhẹ, tiếng thở mỏng như lưỡi dao cạo qua mặt giấy: ``\dots Chỉ riêng việc đứng yên thôi mà cũng tạo ra lực cản như vậy sao?'' Câu hỏi không cần lời đáp, vì những người chứng kiến trên sân đã hiểu: có những người không cần kéo, họ chỉ cần hiện diện để khiến đất trời nghiêng theo.

Đội đối phương bắt đầu dồn toàn lực mạnh hơn. Gã to con cuối dãy gầm vang, giọng đắp đầy keo nhựa sống: ``Một, hai\dots\ kéo! Đứt dây cũng phải kéo!''  Tiếng hưởng theo của đồng đội vang lên như mõ trâm: ``Kéo nát họ ra!'' ``Để họ nếm mùi cỏ!'' ``Đại ca, xéo thêm một bước nữa là về đích!''

Nhưng chính lúc đó\dots

Mũi giày của thành viên đứng đầu trượt khỏi mảng cỏ ướt, kêu một tiếng ``chít'' nhỏ xíu như thắt lòng ai.
Cậu nam sinh ấy ngã về trước, hai tay vẫn siết dây, kéo theo bảy người sau như đồi tuyết đầu mùa bất ngờ lở. Dây thừng rung lên một nhịp *BỤP*, rồi cả hàng người đổ ập xuống cỏ, mồ hôi bắn tung tóe, giày dép văng vãi, như thể có ai giật tấm thảm dưới chân họ.

Kaigou vẫn đứng thẳng, chỉ nhường nửa bước, mắt xanh lạnh băng nhìn thảm cảnh trước mắt. Cô không cần kéo thêm, vì đối thủ đã tự ngã vào lưới của chính mình.

Tiếng còi kết hiệp vừa dứt, khán đài phía Đỏ nổ tung. Suzuki nhún hai lần trên bậc cao nhất, giọng vỡ òa: ``\textbf{Onee-sama, chị làm được rồi!}''

Kaoru vung tay vẽ nửa vòng cung đỏ: ``\textbf{Hay lắm, Kaigou-san! Một pha bẫy hoàn hảo!}'' Tiếng vỗ tay rơi xuống như mưa đá nung nóng, những tràng hò reo nối đuôi nhau, lấp đầy khoảng trống vừa bị dây thừng bỏ lại.

Không khí như có người thổi than hồng vào lò, bừng lên đỏ rực trong nháy mắt. Tinh thần của đội Đỏ, từng tưởng đã lụi tàn sau hiệp một, nay bật dậy như ngọn lửa được rót thêm dầu.

Nhưng gã khổng lồ kia không chịu chấp nhận. Vai hắn rung lên, bước chân nặng như chày giã cối, ngón tay đâm thẳng về phía Kaigou như mũi giáo. Hắn chĩa về phía Kaigou đang đứng chống hông: ``Đội này gian lận! Con nhỏ ấy đứng chơ vơ, chẳng kéo, chẳng đẩy\dots\ Chơi kiểu gì thế!''

Ayame chỉ nhướng mày, mắt lạnh như băng mùa đông. ``Không.'' Cô thả tiếng xuống thềm cỏ, nhẹ mà chắc.
``Luật chỉ cấm buông tay và phạm lệnh; còn việc đứng im, nếu đôi tay vẫn ôm dây, là quyền của cậu ấy.''

Gã nghiến răng, tiếng ken két như đá dần mòn. ``Có bản lĩnh thì kéo đi, đừng núp dưới cái bóng luật lệ nữa, con khốn!''

Kaigou quay sang, ánh nhìn xanh thẳm, phẳng lặng như mặt hồ chưa từng biết gió. Cô thở dài, một hơi thở dài đủ làm lá cỏ úa thêm một phần mùa. ``\dots Được thôi. Kéo thì kéo.''

Nàng khẽ nghiêng đầu, mái tóc búi cao lăn xuống một sợi, lẩn khuất bên má. ``Nhưng mà đừng hối hận đấy.'' Lời nói mỏng như lưỡi lam thái không khí, đủ để gã đối diện nhăn mặt, như thể vừa cảm thấy hơi rát trên da.

\ct

Hiệp ba quyết định mở ra bằng một hơi thở nén lại của cả sân. Không khí như dây đàn kéo căng đến mức chỉ cần thêm một nhịp thở cũng đủ đứt phăng. Ánh mắt người xem đổ dồn thành một dòng thác lặng, chảy từ khán đài xuống nơi hai hàng học sinh đang cắm chân vào cỏ như những cọc tiêu sống.

*\textbf{TOÉÉÉÉT!}*

Tiếng còi vỡ òa.

Chưa kịp thành âm, Kaigou đã kéo. Đúng hơn, là giật nhẹ. Một cú giật ngược về sau, bất thình lình như ai vừa rút lưỡi dao khỏi tim đêm. Dây thừng căng cứng, rung lên một nốt trầm chưa từng được viết trong bản nhạc trường học, rồi bật ngược.

Tám bảng thân người bên kia nhấc bổng, thậm chí cả tên to con kia cũng bị nhấc lên trời. Họ treo giữa không trung một tíc-tắc, những đôi giày còn vương cỏ non, những cái miệng há thành hình trăng khuyết. Không tiếng reo hò (mà đúng hơn là chưa kịp để ai reo một câu), không tiếng gió, chỉ có mồ hôi rơi từ mái tóc ai đó xuống cỏ, như thời gian vừa đứt gãy.

Rồi ``Ể!?'' vang lên, những nét mặt biến sắc hiện lên trước đội bạn và cả những khán giả xung quanh.

Touka há miệng, lời nghẹn lại thành hơi lạnh:  ``\dots Cậu ấy\dots\ khỏe đến đâu vậy?''

Aku nhìn cảnh tượng trước mắt, miệng mấp máy không nói nên lời: ``C-C-Chỉ một cái giật nhẹ thôi\dots\ Mà đã như thế này\dots''

Yae đứng bên cạnh, mồ hôi lăn dài trên thái dương, lòng bàn tay ướt đẫm như vừa bóp nát một viên đá. Cô trố mắt, lẩm bẩm: ``Không đùa được nữa rồi\dots''

Kuon, từ xa, đưa tay lên trán, che khuất nắng, hoặc có lẽ che khuất nỗi choáng: ``\dots Mình biết cô ấy mạnh, nhưng đến mức này\dots\ là quá đáng thật.''

Cả sân vẫn còn nghẹn ở cổ họng, như thể ai vừa nhìn thấy một vì sao rơi thẳng xuống giữa ban ngày, và chưa kịp ước. Khi đội đối phương rơi xuống mặt đất, trận đấu đã kết thúc. Chiến thắng thuộc về đội của Kaigou.

Tên to con lồm cồm đứng dậy, mặt đỏ bừng vì tức giận, rõ ràng không thể nào tin được cách đội thua cuộc một cách chóng vánh. Hắn sồn sống bước tới, rõ ràng là định ``nói chuyện'' với cô nàng duy nhất của Đội Đỏ.

Nhưng Ayame đã chặn trước. Cô giơ tay ra, cười tươi: ``Rồi rồi, không đánh nhau, không đánh nhau\~{} Nếu làm ai đó bị thương thì sẽ bị truất quyền thi đấu luôn đấy\~{}'' Giọng cô nhẹ nhàng, nhưng cách đôi mắt đỏ lóe lên cũng đủ để khiến tên kia khựng lại. Hắn tặc lưỡi, quay đi, hậm hực chửi bới trong miệng.

Ở phía đội Đỏ, bầu không khí hoàn toàn trái ngược. Tiếng cười, tiếng vỗ tay, tiếng gọi tên vang lên khắp nơi. Chiến thắng này không chỉ là một trận thắn, mà còn là một cú bùng nổ tinh thần.

Suzuki chạy tới, mắt sáng rực: ``\textbf{Onee-sama! Tuyệt lắm!}''

Kaigou chỉ khẽ gật đầu, tay xoa đầu cô em gái, như thể đó là điều hiển nhiên. Các thành viên bắt đầu giải tán, chuẩn bị cho môn tiếp theo. Một số người vẫn còn bàn tán về cú kéo ``không tưởng'' vừa rồi.

Kuon đứng ở rìa sân, thở dài nhẹ. Trong lòng cậu, chiến thắng của cô nàng không chỉ là một pha thể hiện sức mạnh - cái đó cậu thừa hiểu từ chính kinh nghiệm hút chết của bản thân - mà còn là một lời nhắc nhở rằng cô ấy luôn ở một vị trí xa cách, lạnh lùng, và không cần ai đồng hành.

``\textit{Cô ấy\dots\ vẫn như vậy,}'' cậu thầm nghĩ. ``\textit{Đứng giữa một đội toàn con trai, chẳng nói một lời, chẳng cần một ai. Nhưng khi ở gần mình\dots\ cô ấy lại khác.}''

Từ phía xa, cậu nhìn thấy cô nàng chạy tới, mặt hiện rõ sự tươi tỉnh, trái ngược với nét mặt như căm hờn ban nãy. Cô gọi lớn, tay vẫy lên cao: ``Kuon-kun! Cậu thấy tôi thế nào?''

``\textit{Thật kỳ lạ,}'' cậu mỉm cười mỏng, ``\textit{một người có thể mạnh mẽ đến vậy\dots\ lại lại cảm thấy dễ chịu khi ở gần một thằng con trai như mình.}''

Cậu xoay người, bước về phía khu vực nghỉ. Môn thi tiếp theo đang đến gần, và Kuon biết, dù có là Kaigou hay bất kỳ ai khác, cậu cũng phải chuẩn bị cho cuộc chơi của riêng mình.

Cậu siết chặt băng đô trên tay, ánh mắt lấy lại sự tỉnh táo. ``\textit{Đến lượt mình thể hiện rồi.}''

\begin{center}
  {\color{vol} Môn thi thứ ba: \uline{Bóng né.}}
\end{center}

Khác với hai môn thi trước, bóng né và bóng rổ được tổ chức trong nhà thi đấu của trường Aogami.

Đây là một công trình lớn nằm ở phía sau khu học chính, với mái vòm cao và kết cấu kim loại chắc chắn. Bên trong, sàn gỗ được đánh bóng sáng loáng, phản chiếu ánh đèn trần thành những vệt sáng dài. Các đường kẻ sân được phân chia rõ ràng, vừa cho bóng rổ, vừa cho các môn thể thao trong nhà khác. Hai bên là khán đài bậc thang, đủ sức chứa hàng trăm học sinh cùng lúc.

Trên cao, hệ thống loa và bảng điểm điện tử được lắp đặt gọn gàng. Những tấm lưới bảo vệ treo dọc theo tường để ngăn bóng bay ra ngoài. Không khí bên trong hơi ấm hơn bên ngoài, mang theo mùi đặc trưng của gỗ, mồ hôi và cao su. Tất cả tạo nên một không gian vừa khép kín, vừa sôi động, rất phù hợp cho những trận đấu tốc độ cao như bóng né.

Khi các lớp bắt đầu di chuyển vào trong, bầu không khí nhanh chóng nóng lên.

Tiếng bước chân, tiếng nói chuyện, tiếng cổ vũ vang vọng khắp không gian kín, tạo thành một âm thanh dội lại liên hồi. Các đội tập trung ở hai bên sân, vừa khởi động vừa quan sát đối thủ. Đội cổ vũ cũng nhanh chóng chiếm vị trí trên khán đài, chuẩn bị băng rôn và khẩu hiệu. Ánh mắt của nhiều người ánh lên sự háo hức, vì khác với tamaire hay kéo co, bóng né là môn thi đòi hỏi phản xạ, chiến thuật và cả tinh thần đối đầu trực diện.

Ayame bước ra, bóng đèn trần in xuống sàn gỗ một cái bóng thon dài, mảnh như lưỡi kiếm chưa rút. Cô dừng giữa vạch trung tâm, hai bàn tay chắp sau lưng, yên ắng đến mức tiếng hơi thở của khán đài cũng phải khép lại. hông khí trong nhà thi đấu như tấm lụa vừa được giật phắt, mọi nếp nhăn âm thanh được gấp gọn vào góc.

``Tiếp theo,'' cô nói, giọng vang lên khắp hội trường, ``là môn bóng né.''

Cô giơ tay phải lên, như thể đang mở một cánh cửa chỉ có thể nhìn thấy bằng ánh mắt. Ánh đèn trượt qua ngón tay cô, rơi xuống sân thành năm dải sáng nhỏ, đủ để hai đội nhận ra chỗ mình sắp đứng.

``Luật cơ bản vẫn là luật bóng né quen thuộc,'' cô tiếp, ``nhưng hôm nay chúng ta thêm vào vài điểm khác biệt để trận đấu trở nên hấp dẫn hơn.''

Ayame giơ ngón tay cái, móng sơn màu đỏ thẫm lóe lên như hạt pháo đầu mùa. ``Một là, chỉ một quả bóng. Hai đội cử đại diện lên giữa. Mình tung bóng, ai giành được và đưa về sân nhà, người ấy điều khiển trận đấu.''

Ngón tay thứ hai vươn ra, mảnh dẻ nhưng cứng cỏi. ``Hai là, thời gian ném không bị giới hạn. Các bạn có thể trao đổi, phối hợp, thậm chí lập chiến thuật ngay trên sân, nếu muốn.''

Ngón thứ ba. ``Ba là, bóng trúng chỉ hợp lệ khi chạm vào thân từ bả vai xuống. Chạm mặt là bất lịch sự, bị coi là phạm quy. Bóng vọng tường trước khi trúng cũng không được công nhận.''

Cô khép ba ngón lại, nắm thành một chấm nhỏ giữa lòng bàn tay, như thể vừa gom cả trận đấu vào trong quyền. ``Người bị trúng sẽ rời sân, không có phép màu cứu rỗi, không có bắt bóng để đánh chén tử thần. Khi một đội cạn người, trận đấu chấm dứt, khối lớp ấy được vinh danh.''

Nàng Hội trưởng mở lòng bàn tay, hạt chấn chỉ còn là đường chỉ đỏ mờ trên da. Cô mỉm cười, nhẹ đến mức nụ cười như chỉ dành cho người đối diện duy nhất. ``Rõ rồi chứ?''

\ct

Hai trận mở màn trôi qua như cơn gió lướt qua mặt hồ, chưa kịp gợn sóng đã tàn.

Khối năm hai bước vào sân trong tiếng reo nhẹ của khán đài, nhưng chỉ ít phút sau, không khí đã nặng như chì. Đội Trắng có những cú ném cực kỳ chính xác - những đường bóng thấp, xoáy, cắt ngang không trung như mũi dao lam mỏng. Mỗi lần bóng chạm áo đỏ, một tiếng thở dài lại rơi xuống sàn gỗ. Đội Đỏ cố gắng uốn mình, lách sườn, lùi vội, nhưng cuối cùng chỉ còn biết đứng nhìn hàng rào đồng đội mỏng dần, như thể ai đang xé từng mảnh lưới giữ họ lại với ánh sáng.

Khối năm hai thua. Không ai nói gì, chỉ có tiếng bóng lăn tròn trên sàn, như hòn đá lăn xuống giếng sâu.

Rồi khối năm ba bước ra, không hứa hẹn, không khẩu hiệu. Họ chỉ đứng thành hàng, những đôi mắt đã nhìn thấy quá nhiều trận đấu, những phản xạ xuất thần khiến khán đài phải ngạc nhiên. Khi bóng được tung lên, họ đón bằng cả lòng bàn tay rộng. Những cú ném của họ không nhanh như tên bắn, nhưng chính xác như lời hứa cũ: bóng tìm đúng chỗ áo trắng, chạm rồi rơi, im lặng và không thương tiếc. Khán đài bắt đầu thở lại, từng đợt, từng đợt, như thủy triều lên vào đêm trăng thu.

Khối năm ba thắng. Họ không reo hò, chỉ quay lại, nhướng mày với đám đàn em vẫn còn thất vọng. Một trận thua đã được gỡ, nhưng cả hai bên đều hiểu: bóng né không chỉ là trò đỡ bóng, mà đó là cách cả lớp học cách thở giữa những mũi tên bay tới tấp.

Và cuối cùng, tới lượt khối năm nhất, nơi Kuon, Shiina và Koishin theo học.

Mười lăm cái bóng xếp thành hai hàng dài, vẻ mặt ai nấy cũng thể hiện rõ sự quyết tâm, thậm chí có người còn đập hai tay lại với nhau để răn đe đối phương. Kuon đứng giữa dãy, mắt như mặt hồ đêm không gợn, phản chiếu cả khán đài lẫn ký ức. Không khí trong nhà thi đấu vừa nóng, vừa ẩm, mùi gỗ mới và cao su non quyện lại thành một tấm màn mỏng, đủ để ai cũng nghe thấy tiếng tim mình đập.

Ở đầu kia, kẻ từng đứng trước cổng trường chặn đứng cậu vì ga-tô, nay lại hiện ra giữa tầm mắt. Hắn xoay cổ tay trong lòng bàn tay, như kẻ thợ săn ngắm nghía viên đạn duy nhất còn lại. Nụ cười nham hiểm của hắn không dài, chỉ vừa đủ để lộ chiếc răng khểnh sắc, như mũi dao găm giữa miệng hổ. Mỗi lần cổ tay hắn rung lên, tiếng *RẮC* nhẹ lộm cộm kêu lên như đang chuẩn bị cho đại tiệc sắp tới.

Những cú ném của hắn đã được truyền thành giai thoại: nghe đâu có thể xuyên thủng cả bức tường lớp học, để lại trên vách một vết tròn đen như mắt quỷ. Người ta đồn rằng chỉ cần hắn nâng cao khuỷu tay, không khí sẽ tự khắc rách toạc thành hai nửa, và quả bóng lao tới mục tiêu như sao băng rơi vào đêm rằm.

Kuon nhớ rõ cái ngày hắn chặn cậu ngoài cổng: bóng chiều nhuốm đỏ tường gạch, tiếng còi xe máy rền vang, và câu nói ``Lần tới tao sẽ đập mày ra bã!'' được thả xuống như một tảng đá chèn cổ. Ký ức ấy vẫn còn nguyên mùi bụi đường, vẫn còn vương vết dầu nhớt trên mặt đường nhựa. Cậu đã quá quen với những lời đay nghiến như thế hồi còn học tiểu học, và cũng không phải là người dễ gì bị động với mấy lời như thế.

Bên cạnh Kuon, Shiina khoanh tay, thân hình mảnh như thanh kiếm còn trong vỏ. Ánh mắt nàng bờm Sư tử không nhìn vào bóng, cũng chẳng nhìn vào hắn; cô nhìn vào khoảng trống giữa hai đội, nơi quả bóng sắp vẽ nên quỹ đạo của riêng nó. Cô biết rõ đối thủ này: hắn từng ném bóng thẳng vào lưng một nữ sinh chỉ vì cô dám cười khẽ khi hắn trượt chân. Cú ném ấy để lại vết bầm tím hình trăng lưỡi liềm, như một dấu ấn của bóng tối lên làn da non.

Tiếng hắn cất lên, vang như thanh la gõ giữa đêm khuya: ``Cuối cùng cũng gặp lại rồi nhỉ, Hikawa.'' Âm cuối còn vương trong cổ họng, hắn đã ném thêm một câu như ném đá xuống giếng: ``Lần này\dots\ mày không có ai cứu đâu.''

Shiina không ngoái, không chớp, chỉ nhẹ nhàng mấp máy đôi môi như thể cắt dây một bản giao hưởng dở dang: ``Mày cũng mạnh mồm gớm nhỉ? Nói trước bước không qua đâu.'' Lời cô mỏng như lưỡi dao cạo, đủ để không khí bị rạch một đường chỉ đỏ, mùi sát khí thoang thoảng giữa sàn gỗ.

Koishin đứng lùi nửa bước, hai tay siết lại trong túc áo. Đôi mắt cô không dõi theo hắn, cũng chẳng nhìn Shiina; cô nàng tóc vàng nhìn vào khoảng trống giữa hai hàng người, nơi quả bóng còn chưa tung lên đã ngửi thấy mùi thuốc súng. Cơ hàm cô cứng lại, như thể đang cắn một hạt bí ngô chưa kịp nở, nuốt không trôi, nhả không xong. Một ngọn lửa nhỏ le lói trong đồng tử xanh lam, vừa sợ vừa muốn, vừa muốn vừa sợ.

Kuon nghiêng đầu, khẽ như gió thổi qua lá tre: ``Chúc may mắn.'' Ba tiếng ấy không nặng, không nhẹ, chỉ như một chiếc lá rơi xuống mặt hồ, nhưng đủ để mặt hồ kia dậy sóng. Chính sự điềm nhiên đó, không quá ngạc nhiên thay, lại khiến đối phương càng khó chịu hơn.

Hai hàng người như hai dải lụa đang được kéo căng giữa sân, mỗi bên chọn ra một mũi nhọn để giành lợi thế. Không khí trong nhà thi đấu như tấm da trống vừa được tháo mặt, mỗi bước chân của người được cử ra là một nhịp gõ khiến căng thẳng thêm dày.

Kuon không nhìn thẳng vào đối phương. Ánh mắt cậu trượt ngang, dừng lại ở bờ vai nhỏ bên cạnh. Akaba Koishin. Mái tóc vàng của cô bị gió quạt thổi tung một lọn, che đi nửa gò má đang nóng. Cô không chớp, mắt xanh lam như hai hòn bi thủy tinh ép chặt vào cổ họng đối thủ, chứa đựng một lời thách thức không cần thành tiếng.

Cậu cúi xuống, giọng nhỏ nhẹ như hơi thở vừa thoát khỏi kẽ lá: ``Cậu muốn giành chiến thắng đúng không?''

Câu trả lời đến ngay, không cần mượn thêm khoảnh khắc nào. ``Tất nhiên.''

Kuon gật nhẹ như xác nhận câu trả lời. Cậu không nói thêm, chỉ nhường một bước, để ánh đèn trần rơi xuống vai cô như mũi tên cuối cùng đã được lắp nỏ. ``Vậy thì lên đi.'' Dứt câu, Koishin chạy lên vị trí trung tâm sân, đối đầu với một cậu con trai ở đội đối phương.

Ayame khẽ nhún cổ tay, tung quả bóng lên trên cao, vẽ một vòng cung ngược sáng rực dưới đèn trần. Lập tức, Koishin bật về phía trước, mái tóc vàng tung bay sau gáy, hai mắt xanh lấp lánh ánh điện của kẻ săn đêm. Cô chụp lấy bóng giữa không trung, bàn tay nhỏ bé siết chặt vỏ cao su nóng hổi như bắt được trái tim đang đập loạn nhịp.

Chân cô chạm sàn gỗ, tiếng *CỘP* khô khốc vang lên như mộtt nhịp trống mở màn. Trong khoảnh khắc ấy, thời gian như co rút lại: cô quay nửa vòng, vai trái hạ thấp, cổ tay vặn nhẹ, cánh tay phải vung ra như lưỡi kiếm rút khỏi bao. Quả bóng lao đi, xoáy một đường thấp sát mặt gỗ, rít lên tiếng xé gió, vụt qua khoảng trống giữa hai hàng người.

*BỘP!* *BỘP! *BỘP!*

``Ái da!'' ``Ối!'' ``Đau!''

Ba bóng dáng của Đội Trắng giật nảy. Họ chưa kịp nhăn mặt, quả bóng cao su đã đập vào thân mình họ: một cú đánh vào vùng xương sườn, một cú va vào đầu gối, một cú gõ lên bắp chân, ba tiếng nổ khô khốc nối nhau như ba hồi trống báo tử.

Tiếng còi *TOÉÉÉÉT!* vang lên, xé toạc màng nhĩ khán đài. ``Trúng ba người!'' tiếng Ayame vang lên.

Khán đài đội Đỏ như bức tường đá bỗng biến thành biển lửa: tiếng reo, tiếng vỗ tay, tiếng gọi tên Koishin dội lên thành một bản giao hưởng nửa đêm giữa sàn gỗ, nơi quả bóng còn lăn tròn, lặng lẽ hút lấy ánh đèn như hòn pha lê khóa lấp lánh trong lòng chiến thắng.

``Chuyện nhỏ như con thỏ,'' cô nàng tóc vàng nói, chống hông hãnh diện.

Song bên kia, Đội Trắng chẳng chút lúng túng. Tên trùm trường - quân át chủ bài được cả khối trông chờ - bước ra giữa sàn. Chẳng cần khởi động, chẳng cần dặn dò, hắn chỉ rướn mắt nhìn đối thủ rồi phóng tay.

Một cú ném duy nhất.

*VỤT!*

Bóng lao đi như tia chớp cắt mặt gỗ, xoáy tròn, rít gió, đập thẳng vào hàng ngũ đội Đỏ. Không ai kịp phản ứng\dots

*BỘP!* *BỘP!* *BỘP!* *BỘP!* *BỘP!*

\dots năm người từ đội của Kuon loạng choạng, rồi lần lượt bước ra ngoài vạch hoặc ngã chỏng vó xuống mặt sàn. Khán đài nín thở, mấy giây sau mới dám thở ra tiếng xì khí như bể bơi vừa vỡ mặt.

*TOÉÉÉÉT!* ``Trúng năm người!'' tiếng Ayame reo lên, vẫn với giọng phóng khoáng đó.

Kuon đứng ngoài vòng chiến, hai tay khoanh, khẽ gật như thể vừa ghi nhớ một chiêu bài hay. ``Ném đẹp đấy,'' cậu thì thầm, giọng bình thản đến mức kẻ đối diện tưởng đó là lời tán dương chân thành.

Tên côn đồ xoay người, nở nụ cười vừa tự đắc vừa khinh miệt. ``Đẹp?'' Hắn khẽ phùng má, phát ra tiếng *hức* như con chó săn vừa xé xong miếng mồi. ``Mày chưa thấy gì cả, Hikawa. Đây mới chỉ là khởi động thôi.'' Hắn vung tay lên trời, bắt lấy bóng đèn trần trong lòng bàn tay ảo, rồi siết chặt như muốn nghiền nát ánh sáng. ``Lần tới, tao sẽ ném xuyên qua cả lớp học của mày.''

Koishin bước lên một bước, vai nhỏ rung lên. Đôi mắt xanh lam vốn dịu dàng nay loé lên tia lửa đỏ. ``Nghe nè,'' cô cất giọng, từng tiếng như viên đạn bạc được nạp vào nòng, ``cậu mạnh miệng quá rồi đấy.'' Cô nghiêng đầu, lọn tóc vàng trượt xuống che nửa gò má đang nóng. ``Cậu sẽ biết tay tôi.''

Không khí trong nhà thi đấu như bị xé rách một đường dài. Cậu Tổng nhếch mép, lùi nửa bước để nhường sân cho hai luồng khí đang va chạm, đáp gọn: ``Shiina-san, lên đi.''

``Để tôi.''

Nàng bờm Sư tử bước lên, mũi chân trượt trên gỗ như lưỡi kiếm rời vỏ. Ánh mắt cô không còn là ánh mắt, mà đó là lưỡi dao găm vào bóng tối, lạnh đến mức khiến đèn trần cũng phải khựng lại. Cô đứng thẳng, thân hình mảnh như cây tùng giữa bão tuyết, rồi vung tay.

Một đường bóng thẳng tắp, không xoáy, không vòng, chỉ đơn thuần là đường thẳng duy nhất giữa vô vàn đường cong. Quả bóng lao đi như mũi tên bạc bị bẻ gãy cánh, cắt ngang không trung, đâm thẳng vào hàng ngũ đối phương.

*BỘP!* *BỘP!* *BỘP!* *BỘP!* *BỘP!*

Năm tiếng nổ khô khốc nối nhau như hồi trống báo tử. Năm bóng dáng áo trắng loạng choạng, rồi lần lượt bước ra ngoài vạch, vai cúi, mắt nhìn xuống sàn gỗ như thể vừa bị tước mất danh dự.

*TOÉÉÉÉT!* ``Trúng năm người!''

Khán đài nổ tung. Tiếng reo không ào mà cuốn, không ồn mà sâu, cứ thế dâng lên từng đợt, vỗ vào lưng những người đang còng. Không khí nóng đến mức mùi cao su như bốc cháy, mồ hôi bắn tung tóe thành những hạt sáng lăn tăn trên sàn. Shiina lặng lẽ khẽ xoay cổ tay, như thể vừa gập lại một trang sách dày.

Khán đài phía Đỏ dậy sóng, tiếng reo như thác đổ. Aku đứng trên bậc cao nhất, hai đầu gối rung như lá tre trước gió, nhưng vẫn cố gắng gào thật to để lấp đầy khoảng trống trong cổ họng: ``Đ-ĐĐừng thua nhé! Không phải game đâu, nhưng\dots\ nhưng chúng ta sẽ chiến thắng!''

Bên cạnh, Shion chống hông, mắt sáng quắc như đèn pha, giọng vang lên như diễn viên kịch đang đọc thoại quyết định: ``Hỡi các chiến binh ánh sáng đỏ rực! Hãy nghiền nát lũ địch nhân đang lởn vởn trước mặt!''

Aku giật mình quay sang, mặt tỉnh bơ như vừa nghe thấy câu thần chú lạc loài: ``Cậu vừa nói cái gì thế, Shion-chan?''

Trận đấu tiếp tục.

Những cú ném được tung ra dồn dập. Quả bóng cao su bay đi rồi quay lại như con chim bị thuần bằng lửa, mỗi lần chạm tay lại tạo thành tiếng *bộp* khô rát. Koishin ném xong xoay người, tóc vàng vẫn còn vương góc sáng; cô không cần nhìn cũng biết bóng sẽ tìm đúng chỗ hiểm hóc trên vai đối phương.

Shiina đứng cách đó ba bước, mắt lạnh như thép vừa rót khuôn; cô chỉ cần khẽ vung tay, bóng đã cắt ngang hàng ngũ, để lại sau lưng năm bóng áo trắng lủi thủi bước ra ngoài vạch. Những người còn lại trong đội Đỏ thỉnh thoảng cũng ``góp vui'': có người ném trúng bắp chân, có người lại để bóng lông vũ bay lọt qua khe hở, khiến khán đài vừa reo vừa tiếc thay cho một cơ hội bị bỏ lỡ. Mỗi lần tiếng còi *toét* vang lên, Ayame liền giơ cao cánh tay, đọc tên như thể đang gọi linh hồn bước qua bên kia thế giới.

Ở giữa tâm bão, Kuon lại chọn cách trở thành bóng ma. Cậu không ném, không chụp, cũng chẳng la hét gì vội; cậu chỉ lắng nghe tiếng gió rít qua viên bóng rồi để cơ thể nhớ kỹ nhịp điệu ấy, quan sát luôn cả cách mà tên côn đồ ném và ghi nhớ từng động tác, từng dấu hiệu một. Có lần bóng lao thẳng về phía cậu, cậu chỉ khẽ nghiêng cổ, để mũi bóng vụt qua mái tóc, sát đến mức vài sợi tung bay lên ánh đèn. Có lần quả bóng lao tới mạnh, cậu lùi nửa bước, vặn hông, để nó tông vào tường. Đồng đội thấy thế, cứ ngỡ Kuon chỉ may mắn; chỉ có những kẻ đối diện mới hiểu, cậu đang đọc vị họ từng hơi thở. Trong gang tấc, khoảnh khắc ấy, Kuon không né bóng thông thường nữa, mà còn né cả ý đồ của người ném.

Đỉnh điểm, tên trùm lùi sâu vào cuối sân, khuỷu tay vung cao quá đầu, tung ra cú ném từ phía sau lưng đội mình. Quả bóng lao đi như tia chớp rạch không khí, xoáy một đường thẳng tắp về phía Kuon. Khán đài kịp thở gấp một nhịp; mọi con mắt dán chặt vào cậu. Nhưng Kuon chỉ khẽ cúi người, một động tác mềm đến mức dường như cậu đang cúi chào bóng tối. Viên cao su sượt qua mái tóc, va đập mạnh vào tường lưới phía sau, rung lên tiếng *VÙ* như búa bổ.

Aku ôm miệng, mắt tròn xoe: ``Trời ơi, Kuon-kun tránh được quả ấy kìa\dots!'' Cô nàng tóc hồng-xanh nhảy tung trên khán đài, hai má đỏ ửng. Suzuki thì hai tay nắm chặt khẩu hiệu, reo nhỏ như thầm: ``Onii-sama thật phi thường\dots''

Cùng khoảnh khắc ấy, Shiina đã xoay người, cánh tay vung tới tận cùng giới hạn khớp vai. Quả bóng phóng vút thẳng tắp như mũi tên bạc nhắm vào ngực tên trùm. Nhưng một thành viên đội Trắng bất ngờ bị kéo xê dịch, lấy thân mình làm lá chắn. Và, *BỘP!* viên cao su đập vào lưng cậu ta; cậu gục ra ngoài vạch, còn kẻ cầm đầu vẫn đứng nguyên, môi nhếch đắc thắng.

Kuon chau mày, thầm nghĩ: ``\textit{Chơi kiểu này hơi không fair-play nhỉ\dots}'' Koishin bước tới, giọng run nhẹ nhưng đủ lớn: ``Cậu dùng đồng đội làm lá chắn sao?'' Tên trùm nhún vai, cười khẩy: ``Luật không cấm đâu. Đúng không, Hội trưởng?''

Ayame gật nhẹ, đôi mắt đỏ vẫn hiện rõ sự tinh nghịch: ``Đúng luật mà. Người bị trúng vẫn phải rời sân.''

Không khí trong nhà thi đấu như tấm màn vừa bị xé ngang: bên này khán giả nín thở, bên kia tiếng reo nghẹn lại nửa chừng. Những cú ném cuối cùng vẫn vang lên đều đặn, mỗi lần bóng chạm áo lại có thêm một bóng hình lủi thủi bước ra khỏi vạch, để lại khoảng trống toát lên mùi gỗ nóng.

Chỉ trong vài nhịp thở, sàn đấu đã trống trải đến lạnh người. Cuối cùng, ánh đèn trần chỉ còn chiếu thẳng vào sáu cái bóng còn trụ lại.

Đội Đỏ: Hikawa Kuon - đôi mắt lim dim như đang đếm nhịp tim đối phương; Fukami Shiina - vai thả lỏng, cánh tay mảnh dẻ như mới rút kiếm khỏi bao; Akaba Koishin - lọn tóc vàng ướt mồ hôi, dính lên gò má đang đỏ ửng vì hưng phấn.

Đối diện, tên trùm khum khuỷu tay, ngón cái gãy gọn vào lòng bóng, nụ cười vẫn còn nguyên vị máu tanh từ những trận trước; hai tên đàn em kế bên lùi nửa bước, để lại khoảng trống đủ cho cơn bão cuối cùng.

Ba đỏ. Ba trắng. Không còn chỗ trốn, không còn ai để làm lá chắn. Chỉ còn những đường bóng sắp cắt ngang không gian, và tiếng còi báo hiệu hiệp đấu cuối chưa kịp vang, khán đài đã như chiếc phổi khổng lồ giữa chừng bị bóp nghẹt.

Tên trùm đứng thẳng trước mặt Kuon, cổ tay vặn nhẹ như đang siết cổ một con rắn vô hình; đôi mắt hắn đỏ ngầu, ánh lên cơn khát không che giấu: hăng máu hay thắng lợi, chẳng còn khác biệt.

``Chỉ còn mày thôi, Hikawa.'' Giọng hắn rít qua kẽ răng, mỗi chữ như mảnh thủy tinh vỡ cào vào không khí. ``Lần này, tao sẽ bóp nát cái vẻ thờ ơ của mày.''

Kuon không dịch chân, chỉ hơi buông vai, như thể cơn giận đối phương là làn gió thoảng qua. ``Ồ, nhiệt huyết thế mới có cốt cách chứ.'' Cậu cười nhẹ, điệu bình thản như đang bình luận thời tiết. ``Tôi bắt đầu thích cậu rồi đấy, loại người máu lửa luôn khiến trận đấu thêm hấp dẫn.''

Koishin giật mình quay sang, mắt xanh tròn xoe: ``Kuon-san, cậu đùa lúc này à!?''

Shiina vẫn lặng. Cô chỉ liếc cậu Tổng một cái, một ánh nhìn thoáng qua như lưỡi kiếm rút nửa chừng, rồi khép mí, dường như đã đọc được kế hoạch đang chạy trong đầu cậu. Không tán thành, không phản đối; cô chỉ đứng đó, kiếm chưa ra khỏi vỏ nhưng đã chọn điểm rút.

Im lặng bao trùm sân, chỉ còn tiếng bóng cao su nảy trên gỗ như nhịp tim đập loạn. Tên trùm nâng cánh tay, mắt đỏ lóe lên tia săn mồi; hắn không cần nói thêm câu nào, chỉ vung tay như thể đang quăng dao. Viên bóng lao tới, xoáy một đường thẳng tắp, gió rít lên như lưỡi lam mới mài.

Kuon khẽ hạ thấp trán, tóc đen xõa xuống che nửa khuôn mặt. Bóng sượt qua mái đầu, gọn đến mức chỉ có vài sợi tung bay lên ánh đèn trắng. Khán đài còn chưa kịp thở, hắn đã xoay cổ tay, phóng cú thứ hai từ góc chéo, nhắm thẳng vào vai trái. Cậu nghiêng mình, lướt qua đúng khoảnh khắc viên cao su vụt tới, để lại trên áo một vệt gió lạnh.

Cú thứ ba tới ngay lập tức, thấp hơn, rát hơn, nhắm thẳng vào đầu gối. Kuon nhún nhẹ, bóng lướt dưới đế giày, kêu *BỘP* xuống sàn như tiếng búa bổ. Ba lần ném, ba lần hụt, khoảng cách giữa hai bên giờ chỉ còn là đường chỉ đỏ căng cứng. Khán đài bắt đầu râm ran, tiếng xì xào lan như sóng trước những lần tránh đạn của cậu Tổng: ``Cậu ta đọc được đường bóng sao?'' ``Không đơn thuần là may mắn nữa rồi\dots'' ``Cậu ấy có phải là con người không vậy?''

Tên trùm nghiến răng, khuỷu tay vung cao quá đầu, tung cú thứ tư với góc xoáy cực hẹp. Bóng lao tới, đổi hướng đột ngột giữa chừng, như mũi tên bị gió thổi lệch. Kuon lùi một bước, nghiêng nửa thân, để viên cao su vụt qua sát mép áo, rồi rơi xuống sàn, lăn tròn trong im lặng. Bốn lần. Bốn lần đều sượt qua. Tiếng reo hò bên Đỏ bùng lên như thác đổ, còn bên Trắng thì nín thở, mắt tròn xoe như vừa chứng kiến điều không tưởng.

Rồi, trong gang tấc, bóng vạch một đường xoáy sắc lẹm, đập vào vai Kuon. Koishin giật mình, tiếng thở gọi vang: ``Trúng rồi--!'' Tên trùm nở nụ cười hở cả hàm răng, như kẻ săn đã ghì được con mồi.

Shino đứng phắt dậy, hai lòng bàn tay nắm chặt, thốt lên trong kinh kinh hoànghoàng: ``K-Kuon-san!''

Nhưng tiếng còi *TOÉÉÉÉT* cắt ngang mọi reo hò.

Ayame giơ cao bàn tay phải, giọng vẫn nhẹ nhưng đủ lan khắp nhà thi đấu: ``Không tính! Bóng chạm tường trước khi chạm người!'' Không khí như bị bóp nghẹt một giây. Kẻ thắng thành kẻ thua, người tưởng mất lại được cứu.

Tên trùm gầm lên, giận dữ đạp mạnh xuống sàn, gỗ kêu *RẮC* như xương gãy, rõ ràng cay cú vì bị tước mất bàn ghi điểm.

Cô nàng tóc vàng thở phào, hai vai buông xuống, lòng bàn tay ướt đẫm mồ hôi lạnh: ``Tưởng cậu đi rồi chứ\dots'' Cậu Tổng lồm cồm bò dậy, thở dài nói, giả vờ thất vọng: ``Chúc mừng vì đã trúng tôi, chỉ tiếc là bàn ghi không được công nhận. Tiếc ghê.''

Shino chậm rãi ngồi lại, tay vẫn đặt trước ngực, mắt không rời khỏi quả bóng đang lăn tròn, như thể nó có thể bật lên bất cứ lúc nào. Nhưng, cô đã có thể thở nhẹ vì sóng gió đã tạm trôi đi.

Chỉ trong nháy mắt, Koishin đã phản đòn. Cô xoay người, cánh tay vung ra như lưỡi kiếm bạc, quả bóng lao đi sát mặt sàn, xoáy một vòng cung nhỏ. Hai bóng Đội Trắng chưa kịp né đã giật nảy, ôm bụng, ôm đầu gối, lần lượt bước ra khỏi vạch. Tiếng còi vang lên, giọng nói lém lỉnh của Ayame nổi lên: ``Trúng hai người!'' khiến khán đài nổ tung, nhưng không ai kịp reo, vì chỉ còn một kẻ đứng lại phía bên kia.

Tên trùm.

Hắn đứng giữa sân, lồng ngực phập phồng, mắt đỏ ngầu dán chặt vào Kuon như thể muốn ăn tươi nuốt sống.   Không còn đồng đội, không còn lá chắn, chỉ còn hắn và ba người đối diện.

``Mày\dots\ Mày là thằng quái quỷ nào vậy?!'' tiếng hắn rít ra từ kẽ răng, vỡ vụn trong cổ họng. ``Tại sao mày lại\dots\ Sao có thể\dots''

Kuon nhún vai, mặt vẫn bình thản như đang trả lời thầy giáo bài tập về nhà. ``Quỷ gì đâu.'' Cậu cười nhẹ, đôi mắt lim dim như người vừa tỉnh giấc.  ``Tôi chỉ là người thường thôi. Chẳng qua có tí phản xạ, nên mới vậy chứ.''

Trên khán đài, Aku há hốc miệng, tay che trước ngực như vừa bị bóp nghẹt. ``\dots `Tí' cái kiểu gì vậy trời\dots?''

Tên trùm nghiến răng, tiếng ken két vang lên như thép rèn trong lò. Hắn không nói thêm lời nào, chỉ siết chặt quả bóng trong tay, mười đầu ngón tay như muốn nghiền nát cả viên cao su.

Trận đấu chưa kết thúc. Nhưng cả nhà thi đấu đã biết: đây không còn là trò chơi nữa. Đây là đấu trường giữa kẻ có những cú ném mạnh nhất, với kẻ có khả năng phản xạ né bóng tốt nhất.

Thế rồi, bất ngờ, Kuon lại lên tiếng, nhưng dường như không có ý mỉa mai châm chọc gì. ``Sao cậu không thử hạ hai người kia xem?'' Không chần chừ, cậu chỉ nhẹ về phía Shiina và Koishin, rồi tiếp tục nói nói với tông giọng nhẹ tênh: ``Biết đâu bóng bật ngược lại mà trúng tôi đấy.''

Khán đài nổ một tiếng ``ồ'' ngắn, như búng tay vào mặt trống.

Koishin quay lại, mắt trố ra đầy ngỡ ngàng với câu nói chẳng khác nào như đang chỉ bài cho đối thủ: ``Cậu muốn chết không vậy?!''

Suzuki trên khán đài níu chặt lan can, không tin nổi vào những gì cô nghe được: ``Onii-sama, anh đang nghĩ gì thế?''

Ngược lại, Shion lại có suy nghĩ khác. Cô gập cằm xuống tay, mắt híp. ``\dots Tâm lý chiến. Hay đấy.'' Aku gật như máy rung, đồng tình với cô bạn.

Shiina nhếch mép, tỏ vẻ hứng thú: ``Ý tưởng không tồi.''

Tên trùm cười khùng khục: ``Sao tao không nghĩ ra từ sớm nhỉ?''

Kuon thở dài, nhìn hắn như thầy giáo nhìn học trò lớp một: ``Suy nghĩ cậu đơn sơ quá.''

Hắn cười to, đập vai đàn em. ``Cảm ơn đã nhắc!''

Cả sân im một nhịp. Không một ai bảo hắn rằng: trong thổ ngữ Etsunan miền Bắc, ``suy nghĩ đơn sơ'' là cách nói tròn mép thay cho ``đầu đất''. Kuon đã chửi xong, mà kẻ bị chửi còn cười rúc rích, khiến cậu chỉ lắc đầu ngán ngẩm.

Cậu Tổng vừa dứt lời, tên trùm đã quay ngoắt sang Koishin. Cô nàng tóc vàng còn chưa hết cáu vì bị Kuon ``bán đứng'', nên chỉ kịp trợn mắt khi thấy bóng lao tới như viên đạn. Không gian như co rút lại trong một nhịp thở, tiếng gió rít qua viên cao su khiến khán đài nín lặng.

``Sang trái.''

Chỉ hai tiếng, gọn như lưỡi dao thái, Kuon ném thẳng vào tai cô nàng. Cô không kịp nghĩ, chỉ lệch vai theo phản xạ. Bóng vụt qua, sượt tà áo, đập mạnh xuống sàn, lăn tròn như con thú bị bỏ lỡ.

Koishin chậm rãi quay lại, mắt xanh loé lên tia sững sờ. ``Cậu\dots\ vừa cứu tôi ư?''

Kuon nhếch mép, điệu cười như vừa giật mất một lá bài trong tay đối phương. ``Chưa gì đã hoảng thế, cậu dễ sợ quá.''

Phía bên kia, tên trùm cũng kịp nhận ra mình bị đá xéo. Hắn nghiến răng, giọng như thanh sắt cà vào đá: ``Mày đùa tao à, Hikawa!? Mày lại giúp nó làm cái củ quặc gì!?''

Kuon xoay nhẹ cổ tay, dáng điệu thư thái như đang chỉ đường cho khách lạc. ``Tôi chỉ nói cậu có thể ném về phía họ.'' Cậu khẽ nghiêng đầu, đôi mắt lim dim. ``Chứ tôi đâu bảo họ đứng yên để làm bia đâu. Thông minh lên chứ, cái não đơn sơ này.''

Một làn sóng ``ồ'' lan nhanh khắp khán đài, như có ai vừa búng dây đàn giữa đêm. Trên khán đài, Kana bật cười khúc khích, tay che miệng: ``Kuon-kun vẫn thích chơi chữ nhỉ, đáng yêu quá đi mất!''

Shion đẩy ảo kính không tròng, giọng trầm như thầy pháp: ``Hỡi kẻ mang danh trùm, ngươi vừa bị phù thủy ngôn từ đánh lừa! Một chiêu tâm linh cao cấp đấy!''

Aku nghe xong chỉ biết gãi đầu, thì thào: ``Cậu ấy\dots\ vừa mắng người ta mà người ta còn cảm ơn, đúng là ma mãnh ghê.''

Gương mặt hắn đỏ bừng như vừa bị tát vào niềm kiêu hãnh, hai bàn tay siết chặt quả bóng đến mức các đốt ngón trắng bệch. ``''\textbf{Được thôi, con chó Hikawa! Ăn miếng này đi!}''

Một tiếng rít gió. Quả bóng lao về phía Shiina như đạn pháo.

Không thèm nhíu mày, nàng bờm Sư tử giơ tay trái, đón bóng bằng lòng bàn tai dẻo như da trống.

*BỐP*

Tiếng vang ngắn gọn, khô rát. Quả bóng nằm im trong lòng bàn tay, như con chim bị bóp cổ. ``Chỉ có thế thôi ư?'' Shiina nhếch môi. ``Mạnh lên đi chứ, tao còn chưa kịp ngáp này.''

Kuon cười khẽ bên cạnh: ``Cậu ấy nói đúng đấy, đừng có mà tiếc sức.''

Tên trùm gầm lên, nước bọt văng tung tóe: ``\textbf{Được, tao cho cả lũ bay biết tay!}''

Từ đó, sân bãi biến thành dây chuyền lửa. Hắn ném - Koishin khựng lại, mắt xanh trống rồng. ``\textbf{Sang phải!}'' cậu Tổng quát lớn. Cô nàng tóc vàng lăn theo tiếng hô, bóng vụt qua tai, gió lạnh táp vào gáy.

Hắn ném - Shiina đỡ, bóng nằm gọn, cô thở dài như người đang dạo chơi: ``Chậm quá.''

Hắn ném - Kuon lùi nửa bước, viên cao su cắt sát góc áo, chỉ có vài sợi vải tung bay. ``Gần đấy,'' cậu thì thầm, ``nhưng chưa đủ.''

Rồi một quả bóng xoáy thấp, lạnh lùng như dao rọc giấy. Koishin bước hụt, mắt cá vặn ra, cô ngã sấp, tay chống xuống gỗ. Nhân thời cơ đó, hắn ném quả bóng lao thẳng vào trán cô.

Khán đài nín thở, tiếng reo nghẹn lại nửa chừng.

Trong khoảnh khắc đó, cậu con trai bật người, bàn tay vung lên như cánh diều đứt dây.

*BỤP*

Bóng nằm gọn trong lòng bàn tay cậu, còn rung rinh như con cá tuyết vừa bị kéo lên khỏi lỗ. Koishin ngẩng lên, tóc vàng rối tung, mắt trợn ra kinh ngạc: ``Cậu\dots\ Sao lại?''

Cậu chỉ liếc qua, gắt giọng cắt ngang: ``Đứng dậy đi. Trận đấu vẫn chưa kết thúc đâu.''

\ct

Sau chuỗi phóng thích liên hoàn, vai tên trùm chùng xuống như cánh cửa bị bỏ quên bản lề. Hơi thở hắn rít qua kẽ răng, từng đợt gấp gáp đánh vào không khí nóng. Lực ném giờ chỉ còn là cơn gió thổi qua mùng, quỹ đạo thẳng tắp bỗng chịu nghiêng theo ý trời.

Một quả bóng vừa chạm tay Koishin đã bị cô bẻ ngược, nằm gọn trong lòng bàn tay như con chim ngủ quên.   Kuon lim mắt, khóe môi nhếch nhẹ: ``Có vẻ như hắn đã bắt đầu xuống sức rồi.''

Cậu khẽ nhặt quả bóng cao su, xoay nhẹ cổ tay, ném về phía mái tóc vàng. Cô nàng Tsundere bắt được, ngón tay siết lại, móng trắng bật cả vào vỏ bóng.

Kuon cúi xuống, giọng thấp như tiếng lá rơi: ``Cậu đã thấy điểm yếu chưa? Đường thẳng cuối cùng luôn nằm trong tay cậu.'' Rồi cậu đổi sang tiếng Anh, nhẹ tênh như kẻ mời bạn dạo chơi: ``\eng{Finish him.} (Hạ gục cậu ta đi.)''

Koishin nâng bóng lên trước trán, tóc vàng ướt mồ hôi dính vào thái dương. Cô nhìn quả bóng cao su, mắt liếc về phía đối thủ đang thở dốc vì mệt:  ``Đừng tưởng tôi chỉ để cho ngắm.''

Tên trùm gầm lên, giọng khàn đặc: ``\textbf{Còn tao đây! Đừng hòng--}'' Tiếng hét bị cắt ngang.

Koishin bước dài, vai xoay mạnh, cánh tay vung tới tận cùng giới hạn khớp. ``\textbf{Ăn đòn đi!!!}'' cô hét vang, giọng xé gió như lưỡi kiếm rút khỏi bao. Quả bóng lao đi như mũi tên bạc rời cung, không xoáy, không lạng, chỉ thẳng tới đúng vị trí mà cô neo tới.

*BỤP!*

Quả bóng đập trúg vào bụng hắn, lõm cả lớp áo, không khí trong phổi bật ra thành tiếng ``ặc'' ngắn ngủi. Tên trùm bật ngửa, hai chân khuỵu xuống, lưng đập xuống sàn gỗ kêu *RẦM* như trống trận dứt nhịp.

Shiina khép mi, giọng ngân nhẹ như ru: ``Chiếu tướng.''

Tiếng còi dài vút lên như mũi tên xé tan màn đêm, kéo theo cả dòng cảm xúc dồn nén từ đầu trận. Ayame nhếch môi, đôi mắt đỏ lấp lánh niềm hả hê khó che giấu; cô vung tay, giọng vẫn ngọt như mía lùi nhưng đủ sức gióng khắp nhà thi đấu: ``\textbf{Đội Đỏ giành chiến thắng!}''

Chỉ một khắc, khán đài như bình nước nóng được mở vung: tiếng reo, tiếng cười, tiếng vỗ tay hòa vào nhau thành bản nhạc dậy sóng. Các thành viên đội Đỏ đổ xuống sàn, áo đỏ chập chờn như ngọn lửa bốc cao. Koishin bị đồng đội nâng lên, tóc vàng rối tung dưới ánh đèn, mồ hôi lẫn nước mắt hạnh phúc lấp lánh trên gò má. Tiếng gọi tên cô vang lên từng hồi, như thể muốn in vào không gian lời khẳng định: ngày này thuộc về họ.

Koishin được đồng đội nâng cao, giọng nói rộn ràng vang lên từng hồi:

``Koishin-chan, đường bóng cuối cùng đó đỉnh quá! Cậu đã hạ được tên trùm khét tiếng nhất trường!''

``Cú ném quyết định của cậu thật tuyệt vời! Không hổ danh là tay bóng ném số một!''

``Chúng ta đã làm được rồi! Koishin, cậu là người hùng của ngày hôm nay!''

Tiếng reo hò không ngớt, những cái vỗ tay rộn ràng, và những nụ cười rạng rỡ bao quanh cô gái tóc vàng. Koishin mỉm cười, mắt xanh long lanh trong niềm hạnh phúc tột cùng.

Ở ngoài vòng xoáy ấy, Kuon và Shiina đứng lặng. Hai bóng hình tách biệt nhưng cùng chung nhịp thở; họ không la hét, chỉ mỉm cười, nụ cười của kẻ biết mình đã gieo hạt, rồi để cây đại thụ tự mình đâm chồi. Ánh mắt họ giao nhau một giây, đủ để hiểu: chiến thắng này không chỉ là điểm số, mà là lời khẳng định cho những ngày tập luyện đầy vất vả và gian truân.

Bên kia sân, tên trùm lồm cồm gượng dậy. Mồ hôi lạnh dính trên thái dương, hắn thở dốc, mắt vẫn đỏ ngầu nhưng đã mất điểm tựa. ``Trận này--'' giọng hắn khàn, ``không công bằng--'' lời nói nửa chừng như bị không khí cắn đứt.

Ayame bước tới, kiễng người xuống, đôi mắt long lanh tròn xoe: ``Luật rất rõ ràng mà\~{}'' Cô nghiêng đầu, ngón tay chạm cằm, giọng vẫn có gì đó cợt nhả nhưng đủ sắc lẹm: ``Nếu còn thắc mắc thì sao không than phiền lên hội học sinh xem? Biết đâu mọi người ở đó sẽ xem xét đấy\~{}''

Hắn cứng họng, lời phản kháng chết lặng trong cuống họng. Cuối cùng, chỉ còn tiếng gỗ kêu ọp ẹp dưới chân khi hắn quay đi, bóng lưng còng xuống, nhưng đôi mắt vẫn ghim chặt vào Kuon, như thể muốn in hình cậu vào đáy mắt, để nhớ lấy kẻ đã khiến hắn nếm mùi bại trận.

Suzuki, Kaoru và Shino lao xuống sàn như sóng vỡ, ba bóng hình nhỏ nhắn hòa vào thác khán giả đang tràn ra.

``Onii-sama! Tuyệt lắm!'' tiếng cô nàng mắt hai màu vang lên trước tiên, mắt cô ánh lên như hạt cườm thủy tinh dưới đèn.

Kaoru vỗ tay đều đặn, gật gù: ``Cậu né ghê thật! Mấy quả xoáy sát nách mà vẫn đứng yên như tượng. Tài thật.''

Shino khẽ nhún chân, hai má ửng hồng: ``Khả năng chỉ huy cũng vậy nữa! Ba người bọn cậu phối hợp cũng nhịp nhàng thật đấy!''

Kuon đứng thấp hơn họ nửa cái đầu, chỉ cười cười rồi gãi gãi mái tóc đen rối: ``Có gì đâu mà phải khen. Chuyện thường ở xã thôi mà.'' Sau khi nhận những lời chúc mừng, cậu liền xoay người, định bước về lối ra. Ở đó, ánh đèn huỳnh quang đã tắt một nửa, để lại khoảng hành lang xám xịt như màn đêm chờ sẵn.

``Ê.''

Một tiếng gọi nhỏ, đủ để khiến cậu dừng chân. Cậu ngoái cổ.

Koishin đứng giữa bóng râm và ánh sáng, tay khoanh trước ngực, vai hơi rung = có lẽ vì mồ hôi lạnh, có lẽ vì điều khác. Mặt cô đỏ không kịp che, lan từ gò má xuống tận mang tai.

``K-Không phải tôi tới để cảm ơn cậu đâu đấy!''

Kuon phì nhẹ, như nghe được câu thoại quen thuộc từ những nhân vật bằng mặt nhưng không bằng lòng điển hình trên màn ảnh. ``\dots Rồi rồi. Khỏi cần nói tôi cũng biết là cậu định làm vậy mà.''

Cô nàng đưa tay phải ra, lòng bàn tay hướng lên, ngón tay còn dính vệt cao su trắng. ``\dots Đập tay.''

Kuon nhìn một giây, đủ để nhận ra ký ức nào đó vừa lướt qua đáy mắt cô, rồi cũng giơ tay.

*BỐP*

Tiếng vỗ nhỏ, ngắn, nhưng đủ làm Koishin khép hờ mi, đủ để cô hiểu: Nếu không có Kuon, cô đã gục ngã từ sớm, chứ đừng nói đứng đến phút cuối để tung cú ném quyết định.

Cậu rút tay trước, xoay lưng. Bóng áo đen lẫn vào dòng người, biến mất sau khung cửa.

Ngoài hành lang, Suzuki và Kaoru đang chia băng đỏ, hô khẩu hiệu cho đội tuyển kế tiếp. Phần thi của Kuon đã khép lại. Giờ, cậu sẽ ngồi ở hàng ghế sau, nhường ánh đèn cho người khác, nhưng vẫn sẵn sàng nhảy xuống, nếu ai đó cần một lời chỉ điểm, hay một cái gãi đầu hiền lành.

\begin{center}
  {\color{vol} Giờ nghỉ trưa.}
\end{center}

Đồng hồ điểm 11 giờ 30 phút đúng.

Tiếng còi nghỉ trưa réo vang, lan tỏa khắp khuôn viên Aogami như tín hiệu tạm nghỉ. Chỉ trong nháy mắt, sân trường biến thành một bãi cờ đủ sắc: khăn hoa, khăn kẻ, khăn ca-rô được trải thẳng tắp, chồng lên nhau thành những tấm thảm ăn trưa rực rỡ. Đây không đơn thuần là giờ ăn, mà còn là ca nạp năng lượng trước khi bước vào bốn môn còn lại: đấu vật, bóng rổ, chạy vượt chướng ngại vật và chạy tiếp sức - những môn thi sẽ đẩy từng đứa trẻ tới giới hạn thể lực.

Kuon, Kaigou, Suzuki, Yae, Kanade, Uzuki và Koishin lẻn ra sân tập sau nhà thi đấu, nơi họ đã dọn dẵn đồ ăn từ sáng một góc cỏ yên tĩnh. Mỗi người một món, mỗi món một vẻ: Kuon gọn ghẽ với hộp cơm hâm nóng và chai nước lọc; Kaigou mở khăn gói cơm nắm ruốc trứng cuộn nhỏ xíu nhưng đầy ắp tình cảm.

Yae khoe ổ bánh mì kẹp thịt còn nóng hôi hổi, kèm theo hộp dâu tây đỏ au; Kanade ``sống ảo'' với hộp salad bảy sắc cầu vồng, trên rắc vài hạt óc chó lấp lánh; Uzuki nhẹ nhàng với bánh ngọt trà xanh matcha và chai nước ép cam ép tươi; còn Koishin thì ``hầm hố'' hẳn một khay cơm đùi gà chiên, rau củ xào và trứng onsen, đủ no đến tận tối.

Riêng Suzuki, cô nàng vừa mở hộp đã khiến cả hội tròn mắt. Kanade cầm đũa lên, chỉ chỉ: ``Khoan đã, phần cậu hơi khủng đấy! Cậu có ăn được hết không vậy?''

Trước mặt Suzuki là ba lớp cơm nắm, kèm theo đùi gà nướng, xúc xích, salad trứng, khoai tây chiên và cả miếng bánh flan nhỏ làm tráng miệng. Cô chỉ cười xoà: ``Hôm nay mình tăng cao năng lượng một chút thôi mà. Đội cổ vũ bọn mình cũng tiêu tốn nhiều năng lượng lắm chứ.''

Yae chớp chớp mắt, cắn nhẹ quả dâu: ``Cậu không sợ lên cân sao?''

Suzuki nhún vai, mỉm cười: ``Cơ địa tớ khó béo lắm, ăn thoải mái cũng chỉ đủ chất thôi.''

Một khoảng lặng nhỏ. Kanade thì thào, vừa gắp salad vừa lẩm bẩm, má hơi phồng: ``Ghen tị ghê\dots'' Koishin khoanh tay, quay mặt ra chỗ khác, nhỏ nhẹ nhưng đủ mọi người nghe: ``Hừ, đúng là người được trời cho\dots\''

Kaigou nhìn hai cô bạn đó, bật cười: ``Em gái tôi là như thế đấy. Ăn bao nhiêu cũng chẳng mập nổi.'' Còn Kuon thì thở dài, gắp miếng cơm lên: ``Đúng như Kaigou-chan nói đấy.'' Cậu vẫn chưa thể nào quê việc cô bé ăn nguyên cả một mâm cơm gia đình trong The Void, trước khi đánh chén cả hộp cơm mà cậu mang đi.

\ct

Không khí bãi cỏ sau nhà thi đấu như được tháo van khi mùi cơm nóng bốc lên. Bảy chiếc khăn picnic xếp thành vòng tròn khép kín, đồ ăn được mở ra từng lớp, tiếng giấy bạc xé rào rào thay cho tiếng trống trận. Họ không còn là những tay đấu vừa rượt nhau trên sân, chỉ còn là học sinh đói bụng, mắt sáng lên trước miếng đầu tiên. Nhìn chung, các cuộc hội thoại cơ bản chỉ xoay quanh việc họ thi đấu như thế nào, đối phương mang tới những món gì để có thể chia sẻ, và những chuyện phiếm xung quanh.

Chẳng hạn như\dots

\il

Yae ngồi gọn trong tấm khăn ca-rô, cằm tựa lên bàn tay mềm như đang tạm dừng một cảnh quay. Cô nhớ lại từng nhịp né của Kuon, giọng thốt ra như thở: ``Trận bóng né lúc nãy căng thật đấy.''

Koishin vừa nhấp một ngụm trà lạnh, cơm trong miệng còn nóng. Cô gật gù, đáp lại như người vừa chạy xong marathon tinh thần: ``Ừ, nhưng dù gì đội mình cũng thắng mà. Ổn hết rồi.''

Kanade khúc khích, đũa gõ nhẹ vào thành hộp salad. ``Ổn mà cậu lại là người kết liễu? Từ `ổn' của cậu bây giờ cao cấp thế nhỉ?''

Má cô nàng tóc vàng bỗng nhuốm màu cánh sen. Cô quay đi, giọng lúng túng như ai vừa bắt gặp bí mật: ``T-Tôi chỉ làm phần của mình thôi, có gì đâu mà\dots''

Kuon ngồi bên cạnh, lưng dựa vào thân cây sồi nhỏ. Cậu không nhìn ai, chỉ buông một câu như thả lưới xuống mặt hồ: ``Cú ném cuối của cậu chắc tay lắm.''

Koishin cứng người, đũa giữa không trung một giây. Cô khẽ ``hừ'', đáp như đã chuẩn bị sẵn trong đầu: ``\dots Biết rồi, khỏi cần nhắc.'' Nhưng đuôi mắt cô hé mở, lấp lánh một góc trời nhỏ vừa được ai đó gõ cửa.

\il

Uzuki cầm bình nước cam, ngoái sang Kaigou, khẽ hỏi: ``Kaigou-chan, cậu còn thi đấu trận nào nữa?''

``Đấu vật trước, rồi chạy tiếp sức cuối ngày.'' Kaigou đáp gọn, vừa nói vừa vặn nắp chai nước.

Kuon nghe xong gật nhẹ: ``Tôi thì lo quay phim và làm phó đội cổ vũ lớp mình.''

Kanade cau mày, chỉ hai người: ``Vậy hai cậu cùng sân à?''

``Kiểu vậy,'' cậu con trai nhún vai, ``một người ra tiền tuyến xung trận, một người đóng vai hậu phương hò reo sau hàng rào. Lý tưởng thì chỉ có một thôi: hướng tới chiến thắng.''

Kaigou phá lên cười: ``Nghe như phim hành động chiến tranh ấy! Kuon-kun, cậu và Suzu phải hét thật to nhé!''

Suzuki giơ cao hai tay, mắt sáng: ``Em biết rồi! Em và Onii-sama sẽ cổ vũ hết mình cho lớp!''

\il

Suzuki khẽ gỡ chiếc dây thun, mở hộp cơm thứ hai. Mùi gà nướng thơm lừng thoát ra, quyến rũ như lời mời gọi muộn màng. Yae ngước lên, đôi mắt to tròn như hai hòn bi thủy tinh lăn vào giữa đường viền: ``Cậu còn ăn nữa á?!''

Suzuki gật nhẹ, đôi mắt hai màu long lanh: ``Mình chưa đủ năng lượng mà. Ít nhất phải thêm một suất nữa mới đủ.'' Kanade buông đũa, thở dài não nề: ``Cơ địa thần thánh thật\dots\ người ta uống nước lọc cũng tăng cân.''

Uzuki mắt sáng lên, chen ngang: ``Vậy chia bọn mình một chút đi! Bọn mình cũng muốn ăn thử các món cậu làm!''

Suzuki cười khúc khích, gập nắp hộp lại một nửa như khúc cao bồi đóng khung cửa quán: ``Được thôi, cứ bình tĩnh, ai cũng có phần!'' Bốn đôi đũa đồng loạt vươn tới, như đàn chim sẻ đậu xuống bậc thềm.

Yae gắp miếng đùi gà vàng ươm, mắt híp lên sung sướng; Kanade nhận trọn vẹn khoai tây chiên giòn rộp, cắn ``rôm rốp'' rồi khúc khích khen ``ngon như mẹ làm''; Uzuki ôm lấy miếng bánh flan mềm mịn, khẽ rung lên một nhịp ``ưm'' ngắn gọn như bản giao hưởng thu nhỏ.

Cả nhóm xì xầm, tiếng cười lan ra thành vòng xoáy ấm, nắng trưa hắt qua tán lá, rọi sáng những miếng ăn được trao đi: những món quà bất ngờ giữa ngày thi căng thẳng.

\il

Yuki bước tới, tay cầm chiếc khăn lạnh lau mồ hôi. Cô khẽ vẫy chai nước như thể đó là tấm thẻ thông hành duy nhất cho phép cô xuất hiện giữa trưa nắng. ``Chào các cậu\~{} Mọi người sum vầy ghê!''

Kuon ngẩng lên, mắt dừng ở đôi găng dài qua khuỷu. ``Chuẩn bị cho hai môn chạy thế nào rồi?''

Yuki khẽ ngồi xuống, váy xếp thành nếp gấc trên thảm cỏ. ``Mình giữ sức từ sáng giờ. Chưa bung hết sức đâu, nhưng cũng cần nạp tí.''

Kaigou nghiêng đầu, tò mò hỏi: ``Chạy tiếp sức có áp lực không?''

``Có chứ\~{}'' cô gật nhẹ, giọng vẫn mềm như bông gòn. ``Nhưng mình không muốn làm cả đội thất vọng.''

Cậu Tổng lim mắt, chợt hỏi thêm: ``Cậu không thấy nóng sao? Găng tay dài đến cánh tay, tất cao quá đầu gối, trời hôm nay cũng chẳng phải dễ chịu gì cho cam đâu.''

Cô nàng tóc đen xoay cổ tay, sợi vải đen co giãn theo động tác như lớp vỏ bảo vệ thứ yếu, cười nhẹ: ``Mình đã quen rồi. Da mình nhạy lắm, thay đổi thời tiết một chút là nổi mẩn. Phải che kín thì mới chạy được.''

Kuon ``ừ'' một tiếng, nhưng trong đầu cậu chợt lóe ý nghĩ: ở thế giới gốc, Yuki là Roboco - tay chân hoàn toàn là cơ giới, chỉ chừa phần thân và phần đầu trông như con người. Có lẽ ở đây, cô cũng được tái hiện lại, dù thân thể đã là da thịt.

Suy nghĩ vừa kịp thành hình thì Uzuki gọi vang từ đầu dãy khăn: ``Kuon-kun, qua đây giúp mình xách nước một chút!''

Cậu đứng dậy, để mặc ý niệm vụt tắt. Yuki nhìn theo, lặng lẽ siết chặt thêm một đường chỉ may trên găng, như thể cả hai thế giới đều cần cùng một lớp áo giáp mỏng.

\il

Ba cô gái Ayame, Shino và Sakura cũng khệ nệ bưng khay cơm len lỏi tới nhóm, rồi khẽ khàng ngồi xuống khoảng cỏ trống bên cạnh, càng làm cho nhóm bạn sôi động hơn. Hội trưởng hội học sinh là người cất tiếng trước, nụ cười tươi như hoa chợt nở: ``Khí thế ăn uống thật vui nhỉ!''

Suzuki khép hờ nắp hộp cơm, ngước lên nhìn hội trưởng, giọng nhỏ nhưng đủ nghe: ``Ayame-chan, cậu làm MC cả buổi rồi, có thấy hơi mệt không?''

Ayame mím môi, đôi mắt đỏ lấp lánh như vừa được ai đó lau sạch bụi bạc. Cô lắc nhẹ đầu, lọn tóc đỏ rung theo nhịp: ``Mình quen rồi, còn lâu mới bằng lúc tập hát liên tục ba tiếng đâu.'' Rồi cô nghiêng đầu, khẽ cười như thể vừa nhớ ra điều gì đó thú vị: ``Chỉ hơi khát thôi, chứ mệt thì chưa đến lượt mình được nghỉ.''

Cô nàng gật gù, mắt hai màu lấp lánh đầy thông cảm: ``Vậy cậu uống nước mình nhé, mình mang cam ép đấy.'' Cô vừa nói vừa đưa chai nước cam tươi về phía Ayame, như thể đó là một món quà nhỏ cho người vừa đứng trên sân khấu cả buổi.

Nàng tóc trắng-đỏ nhận lấy, khẽ nghiêng chai, miệng tươi cười tỏa ánh nắng: ``Cảm ơn nhé Suzu, mình sẽ cần nó để hét to hơn chiều nay.'' Cô nói xong, nâng chai lên như một lời hứa ngầm: cô vẫn sẽ còn làm cả sân vỡ òa thêm nhiều lần nữa.

Shino không đáp lời bạn, chỉ chăm chú nhìn Kuon; đôi mắt cô ánh lên vẻ lo âu nhỏ nhặt. ``Cậu có cảm thấy mệt không?''

Kuon chỉ lắc đầu thật nhẹ, giọng điềm đạm: ``Tôi vẫn còn sức. Không lo.''

Sakura chờ hai người dứt lời mới cất tiếng, cô nghiêng đầu như chim sẻ: ``Vậy còn khâu cổ vũ thì sao?''

Suzuki ngồi kế bên gập nắp hộp cơm, đáp thay cho cậu conn trai, đôi mắt hai màu lấp lánh đầy tự tin: ``Phần đó bọn mình đã xếp lịch kĩ rồi, cậu yên tâm!''

\il

Tới lượt Wata cũng chung vui, hai tay nâng khay đựng những chiếc bánh xếp chồng như núi nhỏ, mỗi chiếc bọc giấy bạc lấp lánh dưới nắng tàn. Cô khẽ cúi đầu, mái tóc đen tết đuôi sam trượt xuống vai, giọng nhỏ như tiếng lá rơi: ``Mình mang tráng miệng đây\~{} Bánh táo nướng, có ai ăn không?''

Shino đang gấp khăn picnic, nghe vậy giật mình ngẩng lên. ``W-Wata-san\dots'' cô lắp bắp, đôi mắt xanh trố ra, ``cậu làm hết mấy cái này từ sáng giờ à?'' sự ngạc nhiên thấy rõ từ cô nàng. Dù biết cô ấy là người thích nướng bánh, nhưng mang hẳn một khay bánh to bằng một cái mâm cơm lại là điều không thể tin nổi.

Ayame từ phía sau chạy tới, vòng tay qua vai Wata, cười toe toét: ``Ôi bạn hiền, cậu đúng là thiên thần đường bánh!'' Nàng hội trưởng nháy mắt tinh nghịch, rồi quay sang Shino, giọng hạ thấp như chia bí mật: ``Nhớ lần trước cậu ấy làm bánh quy hình con cá không? Cậu ăn xong còn giấu vỏ hộp dưới gối tận ba ngày kìa.'' Mặt Shino bỗng đỏ ửng, cô vội quay đi, nhưng tay đã giật lấy một chiếc bánh táo, cắn nhỏ một miếng, mắt lim dim như mèo gặm kem.

``Wata-chan!'' Suzuki vỗ tay vào nhau, đũa vẫn còn kẹp miếng trứng cuộn, ``cậu đến đúng lúc lắm! Đội mình vừa thắng, bây giờ cần ngọt để nạp năng lượng!'' Cô nhón một chiếc bánh, xé giấy bạc cẩn thận rồi cắn ngập, má phúng phính như sóc tích trái cây. ``Ngon quá đi!'' Cô quay sang Shino, đưa thêm một chiếc, giọng réo rắt: ``Shino-chan, ăn đi, không lát nữa hết cả đấy!''

Cô nàng tóc vàng nhạt mỉm cười, khẽ gật. Nắng đầu giờ chiều rọi xiên qua tán sồi, phủ lên khay bánh một lớp vàng ươm. Ba cô gái ngồi quanh, tiếng giấy bạc xé rào rào như nhạc đệm nhẹ, gió thổi thoảng mùi táo quế, khóe mắt ai cũng long lanh niềm vui nhỏ, như vừa tìm được mảnh ngọt giữa ngày thi căng thẳng.

\il

Không khí bãi cỏ trưa cứ như vậy vẫn tràn ngập tiếng cười rôm rảrả. Nắng gắt phủ lên vai học sinh thành lớp áo lụa vàng, tiếng đũa gõ vào hộp nhịp nhàng như trống nhỏ. Mùi gà nướng, bánh mì thơm, cam ép mát vừa bốc hơi đã bị gió đưa đi khắp vòng, khiến bụng ai cũng réo thêm một nhịp. Trên tán sồi, chim sẻ nhảy từ cành này sang cành khác, như đang đếm những câu chuyện vừa kịp mở đầu.

``\dots'' Sakura ngồi yên đó, im lặng nhìn các bạn học.

Bất chợt, cô khép cánh tay, giọng cô nhỏ xuống đến mức chỉ còn là vệt sương mỏng: ``Kuon-san, Kaigou-san, Suzuki-san, và\dots\ Shino-san này.'' Cô khựng lại, mắt lướt qua bốn gương mặt vừa điểm tên, rồi thêm một tiếng gọi như mời ra khỏi vòng xoáy ồn ào. ``Tôi có chuyện này muốn hỏi riêng một chút.''

Họ đứng dậy, không cần trao đổi, bước ra phía sau gốc sồi già. Bóng lá rải xuống thành tấm màn xanh lờ mờ, đủ che những ánh nhìn tò mò. Sakura hạ giọng đến mức gió phải nín thở: ``\dots Các cậu có thấy `người đó' không?'' Cô nhấn mạnh tới hai chữ ``người đó'' ở cuối câu, như thể đang nhấn mạnh tới sự bất thường ở câu chuyện.

Không ai hỏi lại, bởi họ đều hiểu: cô bé mắt xanh, kẻ lạ mà quen, cũng là người tạo ấn tượng nhất. Kuon gãi cằm, trả lời trước: ``Có. Lúc sáng, khi đang chỉnh camera.'' Giọng cậu bình thản như kẻ kể chuyện cổ tích buổi trưa.

Suzuki siết nhẹ tay, móng trắng in lên vỏ chai nước: ``\dots Cô bé làm gì?''

``Giơ nắm đấm,'' cậu đáp, ``cổ vũ.'' Hai tiếng sau như hòn sỏi thả xuống giếng, nghe được cả tiếng vọng mơ hồ.

Im lặng dâng lên, Kaigou nhíu mày: ``\dots Chỉ vậy thôi à?'' Kuon gật, đôi mắt đen như khẳng định một bí mật không cần lời: không có dấu hiệu nào cho thấy mối nguy. Shino thở phào nhẹ nhõm, như thể vừa tháy được đáy giếng: ``\dots Vậy chắc là không vấn đề gì nhỉ?''

Suzuki vẫn không buông, lời nói mềm nhưng day: ``\dots Nhưng tại sao cô bé lại ở đây\dots?'' Câu hỏi treo lơ lửng, không ai đỡ lấy. Mọi ánh nhìn đổ dồn tới người con trai, hi vọng tìm ra được câu trả lời. Nhưng đáp lại họ chỉ là cái nhún vai, cười khẽ như kẻ vừa đọc xong một trang sách thiếu giấy: ``Biết chết liền à. Đến anh cũng không hiểu cô bé tới đây với mục đích gì nữa.''

Khoảng lặng ban đầu chỉ như hơi thở ngắn, nhưng rồi kéo dài ra mỏng tang, đủ để tiếng gió lọc qua tán sồi nghe như ai đang lật từng trang giấy mỏng. Bốn người đứng thành một vòng tròn khép, mỗi người ôm một mảnh nghi vấn, nhưng không ai tìm được đường ghép chúng lại thành lời giải thích trọn vẹn. Trên khuôn mặt họ, ánh nắng trưa in ra những vệt sáng loang lổ, như thể chính ánh sáng cũng đang do dự không biết nên dừng lại ở đâu.

``\textbf{Mọi người ơi\~{}! Ăn nhanh lên, sắp đến giờ thi đấu rồi đấy!}''

Tiếng Yae vang lên từ phía bãi cỏ, ngắt dứt mạch bí ẩn đang chực buông xuống. Câu gọi ấy như một cái vỗ nhẹ vào vai, đủ để họ giật mình rời khỏi mê cung suy tư, quay về với thế giới đơn giản của khăn picnic và hộp cơm. Họ đáp lại bằng một tiếng ``À, được rồi!'' rồi kéo nhau trở lại chỗ ngồi, để lại phía sau gốc sồi già và câu hỏi chưa kịp có lời đáp, lủng lẳng đung đưa theo từng cơn gió nóng.

Không khí bãi cỏ nhanh chóng trở lại nhịp cũ: tiếng cười rơi vào nhau, tiếng đũa gõ nhịp, những câu chuyện nửa chừng được nối tiếp. Cô bé mắt xanh lục như một hạt bụi nhỏ vừa lọt vào mắt, gây xốn xang rồi lặng lẽ trôi đi; mọi người tạm gác nghi ngờ, chuyển sang bàn chuyện đội hình, bàn chuyện thời tiết, bàn chuyện trời chiều sẽ nắng gắt hay mát trời. Bữa trưa chỉ còn là những miếng cơm cuối cùng, những ngụm trà lạnh, và cảm giác no ấm đang lan từ dạ dày lên đầu lưỡi.

Một lúc sau, họ bắt đầu gấp khăn, đóng hộp, vuốt phẳng nếp áo. Tiếng nhựa va nhau, tiếng vải xào xạc, tiếng nắp hộp khép lại như những dấu chấm nhỏ khép lại một trang sách. Mọi người đứng dậy, vươn vai, hít căng lồng ngực; trong làn khí nóng vẫn còn thoang thoảng mùi gà nướng và cam ép, họ nhìn về phía sân thi đấu, nơi bốn môn còn lại đang chờ, nơi những trang mới của ngày hội thể thao sắp được viết tiếp.

\il

Ở đằng xa, dưới bóng cây sồi già, cô bé mắt xanh lục đứng lặng như một bóng mờ. Nắng trưa xuyên qua tán lá, rọi thành những đốm vàng lất phất trên mái tóc đen tết đuôi sam. Cô không cử động, chỉ nhìn năm người vừa rời khỏi gốc cây, vừa kết thúc một cuộc trao đổi khẽ như gió.

Trong mắt cô, ba khoảnh khắc hiện về như thể làn khói mỏng. Trận tamaire, Kuon bị chèn ép, bởi có kẻ đẩy vai định chơi xấu; cô bé đứng sau khán đài cao nhất, bị cột cờ che nửa mặt, mắt không chớp.

Trận kéo co, Kaigou siết nhẹ, thị uy sức mạnh của mình trước đối thủ; cô bé lẩn trong bóng của cờ phướn, nhìn thấy đối thủ đứt đoàn hàng trước khi tiếng còn kịp vang.

Trận bóng ném, Kuon giơ tay chỉ đạo, Koishin xoay vai, Shiina lùi bước; cô bé núp dưới mái tôn hành lang, khuất sau bức pano cổ vũ, thấy từng mảnh ghép chiến thuật được đẩy đúng vị trí.

Và ở cả ba trận, Suzuki luôn ở mép sân, hai màu mắt lấp lánh, giọng hét vỡ òa trong gió. Cô bé nhớ rõ tiếng hô ``Onii-sama!'' hay ``Onee-sama!'' vang lên như một nốt cao xé tan mặt đất.

Cô gật nhẹ, cằm chạm vào vạt kimono đã sờn. Không một lời thoát khỏi khóe môi. Chỉ có hơi thở mỏng manh làm lá cây trước mặt rung động. Rồi bóng cô tan dần, biến mất trước khi nắng kịp đổ xuống gốc sồi.

\il

\begin{center}
  {\color{vol} Môn thi thứ tư: \uline{Đấu vật.}}
\end{center}

Sau khi tấm thảm gỗ được quét dọn vội vã, sân vận động trở thành võ đài ngoài trời. Khán giả vẫn còn vang vọng tiếng reo từ môn bóng né, giờ đã chuyển sang trò chơi mới: đấu vật - môn thi chỉ mới xuất hiện ở Aogami kể từ hai mùa trước.

Ayame bước ra, xoay chiếc míc-rô như đũa. Cô mở màn bằng nụ cười thừa nắng, giọng ngọt như kem tan: ``Nào, trước khi ai đó lầm tưởng, thìthì đây là đấu vật, không phải sumo đâu nhé!''

Hai cánh tay cô vẽ hình tròn trước ngực, minh họa sự khác biệt: ``Sumo thì đẩy đối thủ ra khỏi vòng tròn, còn đấu vật là quật, khóa, ghì để khiến đối phương không thể tiếp tục thi đấu!''

Cô nghiêng đầu, nháy mắt tinh nghịch: ``Và quan trọng nhất: không đấm, không đá, không boxing bất ngờ nha\~{}''

Từ bên ngoài, Shiina lườm nguýt, khẽ tặc lưỡi. Cấm nắm đấm với cô gái luôn coi bạo lực là ngôn ngữ mẹ đẻ quả là trói tay chân. Kaigou đứng bên, mắt nhìn xa xăm, không nói lấy một câu, nhưng cũng không giấu nổi sự ngán ngẩm.

Cô nàng tóc trắng vỗ nhịp ba ngón tay: ``Kỳ hội thao nay có ba bảng: nam đấu nam, nữ đấu nữ, và\dots\ nam đấu nữ!'' Ba giây im lặng đủ để nghe tiếng gió lùa qua khán đài.

Touka nhíu mày, đôi mắt tím như bị mây che: ``Nam đấu nữ? Trong võ đài ư?'' Cô khẽ lắc đầu, như thể muốn gạt đi một ý nghĩ vừa lọt vào tai. ``Bình thường làm gì có môn đối kháng nào cho phép hai giới đấu  nhau nhỉ?''

Kuon đứng bên, tay đút túi, mắt vẫn dán vào sân. Cậu nghiêng đầu một góc vừa đủ để mái tóc đen rối che khuất nửa mắt: ``Nghe như kiểu luật ai đó vừa nghĩ ra lúc nửa đêm, rồi quyết định luôn cho vui.''

Kana bật cười nhẹ, nhưng nụ cười chỉ kịp nở nửa chừng đã bị câu hỏi bóp nghẹt: ``Thế quái nào mà ban tổ chức dám ký duyệt thế nhỉ? Hay có ai đó đang thử xem học viện này chịu đựng được bao nhiêu điều kỳ quặc?''

Ayame hất tóc, đỏ mặt vì tự hào: ``Ý tưởng của tớ mà\~{} Muốn đổi gió một chút cho thêm không khí ấy mừ\~{} Tất nhiên là nếu ai tự tin với khả năng của mình thôi, còn cái này không bất buộc!''

Yae vỗ tay, phá vỡ băng: ``Thôi cũng hay, cùng lắm thì coi như\dots\ gia vị thêm!'' Cô quay sang đội nhà, nháy mắt: ``Shiina, Kaigou, ta cũng cùng nhau cố gắng hết mình nhé!''

Shiina xoay cổ, khớp răng rắc như lò xo: ``Chuyện nhỏ.''

Kaigou gật đầu, cái gật của kẻ đã thấy trước kết cục.

\ct

Tiếng còi vang, võ đài gỗ lập tức trở thành trung tâm của mọi ánh nhìn. Ở bảng nam, từng cặp đấu sĩ bước vào như những mũi tên đã lên cung: vai hạ thấp, mắt dán vào đối thủ, bàn chân trần siết chặt thảm tatami. Không khí nặng hơn, tiếng reo hò như bị bóp nghẹt trong cổ họng khán giả, chỉ còn vang lên những tiếng ``hự'' khẽ khi ai đó bất ngờ bị quật ngửa, lưng đập xuống sàn phát ra tiếng gỗ rúng động. Những tiếng vỗ tay mới nổ ra như thác lũ, xua tan bầu không khí căng cứng.

Chuyển sang bảng nữ, sân khấu như được thắp sáng bởi một thứ năng lượng khác: mềm mại nhưng đầy sức sống. Các cô gái tung người, tóc tết bay vút theo quán tính, ánh mắt sắc như dao cạo qua không khí. Một pha nghiêng người né cú vật tay, rồi bật ngược lại bằng đòn khóa khuỷu, khiến đối phương gục xuống như cánh hoa rụng. Khán đài nổ tung, tiếng huýt sáo và reo hò hòa quyện thành bản giao hưởng giữa trưa nắng. Không ít người phải thốt lên: ``Mình tưởng con gái chỉ dịu dàng, ai ngờ lại dẻo dai và dữ dội đến thế!''

Trang phục cũng góp phần làm nên màu sắc riêng. Nam sinh trần trùng trục, từng múi cơ cuồn cuộn lấp lánh mồ hôi dưới nắng, như những bức tượng sống động. Nữ sinh khoác lên mình bộ đồ thể thao bó sát, vừa khoe đường cong vừa tôn vinh sức mạnh ẩn giấu: bờ vai thon nhưng chắc nịch, cặp đùi săn gọn cuộn tròn trong từng bước dậm chân. Mỗi lần họ vật nhau, khán giả như được xem cả hai thái cực: cơ bắp cuồn cuộn đối đầu với sự uyển chuyển dẻo dai, tạo nên bức tranh đối lập đầy cuốn hút.

Khi trọng tài giơ cao tay, công bố người thắng cuộc, ánh mắt tất cả như được dệt lại thành một mạng lưới cảm xúc: ngưỡng mộ, tiếc nuối, phấn khích, và cả sự tự hào xen lẫn. Vài cô cậu học sinh còn giơ điện thoại, lia máy ảnh lia máy quay, cố gắng bắt trọn khoảnh khắc đối thủ bị ghì xuống, khoảnh khắc mà sau này, khi xem lại, họ vẫn thấy tim mình đập rộn ràng như đang chứng kiến trận đấu ấy lần nữa.

Khán đài phía nam bỗng nổi sóng. Suzuki giơ cao tay phải, Kaoru vung cờ đỏ, cả hai đồng thanh hô vang như hai thanh kiếm chắn trước mũi tên thời gian.

``\textbf{1-C! Cố lên!}'' tiếng vang vọng lên mái tôn, rơi xuống thành mưa nhiệt huyết.

Phía sau, Kuon không cần loa, chỉ cần lòng bàn tay.

Cậu khẽ đập từng nhịp ``ba-hai-một'', điều tiết hơi thở cả trăm con người thành một bản giao hưởng duy nhất. Mỗi câu hô được cắt đúng mạch, mỗi tiếng vỗ tay nổi đúng phách, như thể cả khán đài chỉ là một chiếc trống da khổng lồ và cậu con trai là người cầm dùi.

Tiếng reo không chỉ bay vào tai, nó thấm qua lớp áo thấm mồ hôi, len xuống bắp chân các tuyển thủ, biến thành lửa ở gan bàn chân họ. Trên sàn đấu, có đứa vừa bị quật ngã bật dậy ngay, nhưng không phải vì trọng tài, mà vì nghe thấy giọng Suzuki nhắc: ``Đội Đỏ chiến thắng!'' như một liều thuốc doping tinh thần hữu dụng.

Tiếng reo của 1-C như sóng lớn, lớp này chồng lớp kia, cuốn bay cả cơn nắng gắt đầu buổi chiều.   Kuon vỗ tay đều đặn, nhưng khi thấy vài đôi vai đội Đỏ đã hơi gục, cậu nghiêng đầu nói nhỏ: ``Chưa hết giờ đâu, còn hai hiệp nữa mà.''

Suzuki nghe thấy, liền giơ cao lá cờ đỏ, hét vang: ``Không sao đâu, ình còn bảng nữ đấy! Đội Đỏ, vực dậy tinh thần đi!!''

Tiếng cô như cây đũa thần, mấy cái đầu vừa rủ xuống bật lên ngay. ``Dạ!'' một dàn nữ sinh đáp lại, vỗ vào nhau như chuẩn bị ra trận mới.

Khi trọng tài thu dây vạch, Ayame bước giữa võ đài, phất tờ bảng: ``Tổng kết\dots\ Đội Trắng thắng nội dung nam!''

Khán đài bên kia nổ tung, cờ trắng phấp phới. Dù vậy, Đội Đỏ không ai bỏ cuộc; họ huýt sáo, vỗ tay thật to như muốn gửi luôn năng lượng sang cho bảng nữ đang chờ khởi tranh.

Suzuki nắm tay Kuon, giật nhẹ: ``Sang hiệp của chị em mình rồi, Onii-sama!''

Cậu Tổng đập tay cô: ``Ừ, cùng nhau hét thật to nhé!''

\ct

Sau giờ nghỉ hơi, khán đài lại nóng dần lên với bảng nữ đấu nữ.

Yae, Shiina và Kaigou lần lượt bước lên tấm thảm gỗ, ba bóng dáng quen thuộc khiến khán khán đài phía nam bỗng rộ lên tên họ như sóng vỗ. Khác với vẻ ``đấm đá bất chấp'' lúc tập, Kaigou hôm nay như chỉ dùng lực ``vừa đủ'': những pha quật vẫn dứt khoát, vẫn thu về điểm, nhưng không hề kéo dài hay khiến đối thủ nằm lâu.  Cô gái tóc đen rối khẽ cúi đầu xin lỗi sau mỗi lần ghì vai, rồi kéo bạn đứng dậy như thể nói: ``Chỉ là thi đấu thôi, đừng giận.''

Yae thì ngược lại, mềm mại như dải lụa nhưng ẩn trong đó là lò xo: cô nàng né khóa khuỷu bằng cách lùi nửa bước, xoay người như đang múa, rồi bất ngờ hạ gục đối phương bằng đòn quật hông cổ điển.  Tiếng ``Ồ'' từ khán giả vang lên, xen lẫn tiếng cười khúc khích của mấy bạn nam ngồi hàng đầu.

Shiina đơn giản hơn: đến, thấy, ghì, rồi về.  Mỗi lần cô bước xuống thảm là một lần đối thủ ngồi thở dốc, còn khán đài lại được dịp réo tên cô như bản trống trận. So với những trận đánh nhau như cơm bữa, cô vẫn còn khá nương tay, vì dù gì họ cũng chỉ là những tuyển thủ và không có ý thách thức cô.

Khi trọng tài đưa tay về phía đội Đỏ lần cuối, Ayame giật míc-rô, hét vang: ``\textbf{Đội Đỏ thắng bảng nữ!}''

Tiếng reo nổ tung, cờ đỏ phấp phới như sóng lửa.  Giữa trận bão âm thanh ấy, Suzuki nhảy xuống hai bậc khán đài, ôm lấy Yae vừa bước lên: ``Yae-chan, cú quật hông đó đẹp như phim!  Cơ bụng chắc lắm rồi đấy!''

Cô nàng tóc xám đỏ mặt, vừa thở vừa cười: ``Cậu mà cổ vũ to hơn nữa chắc tớ bay lên mất!''

Kuon đứng sau, đập vai Kaigou nhẹ: ``Cô chừng mực quá, đối thủ còn tưởng được chiêu đãi.''  Cậu nói nhỏ thêm, ``Nhưng đúng là đẳng cấp, không cần ra đòn nặng.''

Kaigou gãi má, có chút ngại: ``Con gái mà. Tôi cũng chẳng định làm họ bị thương đâu.'' Cô liếc về đám con trai ở đội đối diện, mặt tối lại như ăn tươi nuốt sống tới nơi: ``Nhưng mà đám đó\dots\ chắc tôi sẽ không nhẹ tay đâu.''

Suzuki liền xoay qua Shiina, mắt long lanh như đèn pin: ``Shiina-chan, ba lần đều ghì chuẩn luôn!  Chiều nay chạy tiếp sức, cậu giữ sức nhé, còn nước rút đấy!''

Shiina gật đầu, thoáng cười hiếm hoi: ``Ừ, cảm ơn cả hai cậu.  Đội mình còn một chặng nữa, đừng hét khản cổ sớm đấy.''

Kuon giơ cao chai nước, nhìn về phía sân điền kinh đang được kẻ vạch: ``Nam-đấu-nữ sắp tới rồi.  Ai thi thì cứ nhớ: thắng được thì thắng, thua cũng phải về còn nguyên xương, nghe chưa?''

Suzuki cười toe toét, vỗ tay vào nhau tạo thành âm thanh giòn: ``1-C, tiếp tục nào! Cho mọi người thấy sức mạnh của các cậu đi!''

Khán đài lại vang lên tên họ, như khúc ca chưa dứt, hướng tới trận cuối của ngày thi đấu.

\ct

Và rồi, môn thi đã đến phần được mong chờ nhất. Nam đấu nữ.

Yae bước vào võ đài như một dải lụa xám được thả từ trên xuống, mắt cô không nhìn khán giả mà nhìn thẳng vào đôi vai rộng của đối thủ - một nam sinh cao hơn cô cả cái đầu, cơ bắp cuồn cuộn dưới lớp áo thi đấu bóng mồ hôi. Cô mỉm cười, không phải để chào hỏi, mà như thể vừa đọc xong một câu đố và tìm ra đáp án.

Tiếng còi vừa cất, cả hai lao vào nhau không chút dò xét. Bàn tay của cậu bạn như càng cua bám vào cổ tay cô, nhưng cô đã lùi nửa bước, xoay hông, khiến lực bật ngược trở lại làm cậu loạng choạng. Họ quay vòng trên tấm thảm gỗ, những bước chân trần kêu ``sộp sộp'' như đang đếm nhịp cho một điệu tango chỉ có hai người biết. Khóa tay, giật khuỷu, đổi trọng tâm; mỗi động tác đều được cô ghi bằng cơ bắp chứ không bằng mắt; cô cảm lực của đối phương qua đầu ngón tay, rồi trả lại một góc nghiêng vừa đủ để bị mất thăng bằng mà vẫn tưởng mình đang kiểm soát.

Touka đứng ngoài, hai tay khoanh trước ngực, mắt vàng lia nhanh như máy đo quỹ đạo. ``Cân bằng lực tốt\dots\ chuyển trọng tâm linh hoạt,'' cô thì thầm, giọng vừa đủ cho Kaoru bên cạnh nghe thấy. ``Cô ấy đang kéo trận đấu về thế kiểm soát như kéo dây kéo, khóa chặt từng răng.''

Bất ngờ, Yae thả lỏng vai, như thể bỏ cuộc. Nam sinh giật mất đà, thân trên lao về trước; chỉ trong nháy mắt, cô vặn hông, đưa một chân vào giữa hai chân đối thủ, bật nắm áo, xoay tròn. Một ``đòn hông cổ điển'' nhưng được nêm thêm một chút mùi vị của nữ sinh 1-C. Cơ thể cậu bay lên, không cao lắm, chỉ đủ để không khí kịp lọt qua khe hở giữa lưng và tấm thảm. Tiếng ``đập'' khô khan vang lên, khán giả nổ tung, tiếng reo như bong bóng xà phòng vỡ giữa trưa nắng.

Nhưng trận đấu chưa dứt. Anh vận động viên nghiến răng, lăn người thoát khỏi đòn ghì, tay vẫn giữ cổ tay cô như thể đó là phao cứu sinh. Họ lại cuộn vào nhau, mồ hôi hòa thành một vệt bóng loáng trên thảm. Yae thở bằng miệng, nhưng mắt vẫn cười - cô thích cảm giác này: khi mọi mánh khóe đều được đối phương đọc vanh vách, và cô phải viết tiếp câu chuyện bằng một chữ ký mới.

Cuối cùng, cô tìm được khe hở: một khoảnh khắc khi cậu đưa vai về trước để tấn công, hông hở ra như cánh cửa bật. Cô luồn tay vào nách, gài khuỷu, đẩy trọng lượng cả cơ thể lên đốt sống anh: một đòn khóa vai biến tấu. Cậu nam sinh cố gồng, nhưng chỉ kịp nghe tiếng gỗ ``rộp'' dưới lưng và tiếng trọng tài đếm ``một, hai\dots''. Chưa đếm về ba, cậu đã vỗ thảm, đầu lắc nhẹ như thể vừa tỉnh giấc mộng.

Khán đài bùng nổ lần nữa, nhưng Yae không giơ tay ăn mừng. Cô đứng dậy, vuốt phẳng vạt áo, cúi đầu chào đối thủ: một cái cúi nhẹ, đủ để che nụ cười thỏa mãn. Phía ngoài, Touka thả một hơi dài, môi mấp máy: ``Hoàn hảo. Không dư một gam lực.''

\ct

Shiina bước vào giữa thảm, dáng đứng thong thả như kẻ đi dạo, không hề tốn sức làm nóng. Đối thủ vừa nhích chân lên, cô đã lướt tới sát mặt, khoảng cách co rút thành một hơi thở. Không có đòn dọa, không có giằng co; chỉ một cái chớp mắt, cổ tay đối phương đã nằm gọn trong lòng bàn tay cô, xương khuỷu bị gập về phía trước như khóa cửa tự động đóng sập.

Cô xoay hông một vòng cung nhỏ, trọng lượng cơ thể dồn xuống gót chân, đối thủ bị cuốn theo quán tính, đầu gối khuỵu xuống trước khi kịp nhận ra mình đang bay. Một tiếng gỗ khô rục vang lên, vai áo dính sát thảm; toàn bộ màn đấu kết thúc chỉ trong ba nhịp đếm ngắn ngủi.

Aku trố mắt, hàm rơi xuống như vừa chứng kiến ảo thuật. ``C-Cái gì vậy trời\dots\ nhanh quá!'' Tiếng thét thốt ra giữa khán đài, nhưng bị bóng dáng cô gái phủ lên. Shiina đứng dậy, lòng bàn tay phủi nhẹ cánh tay, dáng vẻ như vừa gỡ xong một chiếc lá dính trên vai.

Cô chỉ vọn vẻn nói ``\dots Xong rồi.''

Không ai bật reo, không ai trầm trồ; chỉ có làn sóng gật đầu khẽ của khán giả, như thể họ vừa xem lại một cảnh quen thuộc đã chiếu đi chiếu lại nhiều lần trong đầu; chẳng ai lạ lẫm gì với cô nàng bờm Sư tử - chị đại của ngôi trường Aogami này cả. Trên khán đài, Kuon nhếch mép, Suzuki khẽ huýt sáo, còn Kaigou khoanh tay, mắt lim dim: ``Vẫn vậy, chưa từng thay đổi.'' Trận đấu với Shiina không cần đến hiệp hai; chỉ cần cô muốn, đối thủ đã nằm xuống trước khi trọng tài kịp thở.

Cuối cùng, Kaigou - quân bài bí ẩn của giải - bước vào sân. Dáng cô thon thả, mảnh dẻ, trông chẳng khác nào cành liễu giữa đám sồi gậy.

Đối thủ là một nam sinh to con, vai nở, bắp tay cuồn cuộn. Hắn nhếch mép, cười khẩy: ``Nhóc này ổn không đấy? Trông nó không khác gì cái que củi khô ấy!''

Tiếng cười khúc khích vang lên sau lưng hắn. Cô nàng không đáp, chỉ đứng im, khẽ rướn vai xoay cổ để khởi động các cơ. Một làn gió thổi qua, tóc cô khẽ bay, phơi ra chiếc cổ trắng nhỏ nhưng cứng cỏi.

Tiếng còi vang lên. Nam sinh lao tới, bước chân nặng nề đập xuống thảm. Và\dots\ Một giây, hai giây\dots

*RẦM*

Cơ thể to lớn bất ngờ quay cuồng, đập xuống mặt thảm. Trận đấu đã kết thúc trước sự ngỡ ngàng và ngơ ngác của khán giả. Không ai thấy rõ cô đã làm gì, chỉ biết đối thủ nằm sấp, mồ hôi túa ra như vừa chạy marathon.

Touka lấy tay quệt trán, giọng run run: ``Lực siết tay vượt xa ngưỡng thường. Cậu ấy\dots\ không phải con người.''

Shion tròn mắt, hai tay chụm trước ngực như đang đọc thần chú: ``Phà! Đòn chí mạng của Ma Nữ Bóng Tối\dots\ kẻ thù ngã gục chỉ trong nháy mắt! Làm tốt lắm, thuộc hạ!''

Sakura úp miệng, thì thầm như sợ ai nghe thấy: ``\dots Không thể tin nổi. Cậu ấy có phải là con gái không vậy\dots?''

Kaigou buông tay, quay lưng. Bước chân cô nhẹ như không hề vừa quật một gã to gấp đôi mình. Không khoe khoang, không liếc nhìn, không một lời, chỉ có bóng dáng thon dài khuất dần sau hàng ghế, để lại hơi lạnh lẽo trên tấm thảm gỗ.

\ct

Tiếng còi dứt, khán đài còn chưa kịp ngừng vang, Ayame đã nhảy lên bậc cao nhất, giơ tấm bảng vẽ huy hiệu đỏ chói.  ``Hạng mục nam đấu nữ\dots\ Đội Đỏ thắng!'' Cô cười toe toét, hai tay giơ lên cao, hét lớn vào míc-rô:  ``Tổng kết lại: \textbf{Đội Đỏ giành chiến thắng chung cuộc!}''

Tiếng reo không thành lời, chỉ còn là sóng âm. Yae, Shiina, Kaigou - ba cái tên được khán giả nhắc đi nhắc lại như điệp khúc.

Kuon đứng ngoài rìa, khoanh tay, thấy gió nóng thổi qua cổ. ``Không tệ. Không tệ chút nào.''

Suzuki vỗ nhẹ vai cậu, mắt híp lại: ``Anh nên cảm thấy hãnh diện di, Onii-sama! Lớp 1-C chúng ta mà lị!''

Khán đài bắt đầu dọn ghế, tiếng bước chân lộp cộp. Người ta hối hả chạy sang sân bóng rổ, nội dung thi đấu tiếp theo đang chờ, nhưng trong lòng họ vẫn còn vang vọng tiếng thảm tatami vừa rúng động.

\begin{center}
  {\color{vol} Môn thi thứ năm: \uline{Bóng rổ.}}
\end{center}

``Môn thi tiếp theo của hội thao năm nay sẽ là\dots\ bóng rổ!''

Cả trăm cái đầu quay về phía sân bóng rổ, nơi vừa được kéo rèm, để lộ ba vạch sơn trắng còn thơm mùi nước bảng mới. Không cần giới thiệu dài dòng, Ayame đã cười toe, ném quả bóng lên cao rồi đón lại bằng hai ngón tay như diễn viên xiếc; luật thi được cô kể lại bằng những câu cắt nghĩa ngắn, lại xen thêm câu đùa khiến khán giả vừa ghi nhớ vừa bật cười, căng thẳng tan đi như bong bóng xà phòng.

Các khối lớp lần lượt bước ra, áo đỏ trắng xen kẽ, nhưng ánh mắt mọi người đều tìm về cô gái ``vận động viên'' của khối năm nhất - cũng là người duy nhất đại diện cho lớp 1-C tham gia: Yae.  Cô gái tóc xám nhạt không cao nhất, không nổi bật bằng cơ bắp, lại khiến người ta nhìn bằng cảm giác ``chỗ này sắp có chuyện hay''.  Và đúng như vậy, ngay lần chạm bóng đầu tiên, cô đã chứng minh mình không chỉ đến cho vui.

Yae không chạy nhanh, nhưng bước chân cô luôn đến trước quả bóng nửa nhịp.  Cô đưa tay ra, như thể mời gọi trái bóng vào đúng khoảng trống mình đã đặt sẵn; rồi khom người, xoay hông, nhảy lên; tất cả chỉ trong một dòng chảy duy nhất, không góc cạnh, không dư thừa.  Khán giả bên dưới vừa kịp ``ồ'' lên thì bóng đã rơi xuyên qua rổ, lưới rung nhẹ như vừa được gảy một nốt thăng.

Touka đứng khán đài trên, hai tay khoanh trước ngực, mắt dán theo từng đường cong cơ thể Yae. Cô thì thầm như đang đọc một bản nhạc được viết hoàn hảo. ``Cách di chuyển của cậu ấy tối ưu đến mức gần như không có động tác thừ.''

``Cậu làm tôi thấy sợ vì mấy câu nói máy móc đấy, Touka,'' Shiina đứng bên cạnh càu nhàu, khẽ húc cùi chỏ vào hông đối phương. ``Cứ tận hưởng đi.''

Bên cạnh, Shion khẽ nhếch mép, giọng đầy hàm ý: ``Quái vật nhỏ bé ấy đang viết bản hùng ca riêng trên sàn đấu.''  Kanade thì không nói, chỉ dí hai tay vào nhau mỗi khi bóng lọt rổ, đôi mắt cô sáng lên như đứa trẻ được thêm kẹo.

Hiệp đầu tiên diễn ra với tốc độ khá cân bằng. Đội đối phương cố gắng kèm chặt Yae, thậm chí cử hai người theo sát mỗi khi cô có bóng. Nhưng chính điều đó lại mở ra khoảng trống cho đồng đội, và Yae tận dụng điều đó một cách hoàn hảo. Cô chuyền bóng đúng lúc, đúng chỗ, tạo ra những pha ghi điểm không ngờ tới. Không phải lúc nào cô cũng là người ném rổ, nhưng gần như mọi đường bóng đều đi qua cô.

Khi tiếng còn báo hiệu hiệp một vang lên, Yae đã ghi năm điểm.  Cô như người lái đò, khéo léo đưa dòng nước đi đúng hướng mình muốn; phần còn lại của đội chỉ việc đợi ở cuối con nước để đón bóng và đưa vào rổ.  Khán đài phía Nam vỗ tay theo vòng tròn, tiếng reo không lớn nhưng đều đặn, nhịp nhàng như đang vỗ nhịp cho một bản giao hưởng mà Yae là người cầm đũa duy nhất.

Hiệp hai vừa bắt đầu, đối phương liền chuyển sang phòng ngự khu vực, nhằm ngăn cách triệt để lợi thế của đội bạn. Yae không chậm trễ: cô bất thình lình bứt tốc, luồn qua khe hở giữa hai cái bóng cao lớn, để lại phía sau tiếng giày ken két và cái giật mình của khán giả.

Ở tận cùng đường biên, cô nhảy vội, thân mình nghiêng đến mức tưởng chừng sắp đổ, nhưng cổ tay vẫn mềm mại tung bóng. Quả bóng xoáy nhẹ, vẽ một vòng cung thấp, rồi *ROẸT*, vào lưới. Tiếng reo vỡ òa, đặc biệt từ khối khán đài đỏ rực như được đổ thêm dầu.

Touka gật gù, đôi mắt tím khẽ long lanh: ``Khả năng điều chỉnh trọng tâm đến từng mi-li-mét. Không thể tin nổi.'' Kanade siết nhẹ tay vô thành lan can, thì thầm như tự nói với mình:  ``Không còn là năng khiếu nữa rồi. Đằng sau mỗi cái xoay người ấy là cả trăm buổi tập đêm, cả nghìn lần ngã xuống rồi đứng dậy.''

Trận đấu trôi về hiệp cuối, sân bóng như dây đàn căng cứng. Điểm số díu nhau từng con số, mỗi lần bảng tỉ số nhấp nháy, khán giả lại nín thở. Đối phương bỗng biến thành bầy sói: hai người bứt ra, lao thẳng vào giữa sân, tay quơ ngang địa đoạt bóng.

Yae lúc ấy đang ôm trái bóng bằng cả trái tim nóng. Cô khựng nửa bước, mũi giày ken két trên sơn, thân hình nghiêng đến mức tưởng chừng sắp đổ. Trong khoảnh khắc ấy, cô xoay cổ tay, bóng trượt khỏi tầm tay đối phương, lăn về phía sau lưng, một đường chuyền ngược mờ ảo như ảo thuật. Khán đài ``Ồ'' lên, tiếng reo chưa kịp vỡ thì bóng đã về tay đồng đội, rồi lập tức trở lại Yae.

Thời gian dần chạy về những giây cuối cùng. Không ai kịp đếm nhịp. Cô nhảy luôn, hai bàn tay nâng bóng lên trước mặt, mắt nhắm hờ như đang ngắm sao. Quả bóng rời tay, xoáy nhẹ, vẽ một vòng cung thấp, mềm mại như lông chim.
Không khí tưởng như dày lên, mọi tiếng ồn bị bóp nghẹt.

*BỤP* Lưới rung nhẹ, tán ra thành một chấm trắng xóa.

Shion bật cười, giọng vỡ ra giữa tĩnh lặng: ``Vào rồi!''

Kanade đập tay xuống lan can, gần như hét: ``Ba điểm!!''

Tiếng còi kết thúc nổi lên ngay sau đó, như mở nút chai xì-gon. Người ta ôm nhau, nhảy cẫng, còn Yae chỉ đứng giữa sân, tay vẫn giơ cao, ngón cái khẽ chạm vào trán, một cái cúi đầu nhỏ trước trời chiều đang bắt đầu tắt nắng.

Nàng Hội trưởng hội học sinh vung tay lên như thể vừa kéo rèm khép lại một vở kịch, giọng cô vang lên đầy khoái trá: ``Và đội chiến thắng\dots\ \textbf{là Đội Đỏ!}''

Tiếng reo vỡ òa, lan tỏa từ hàng ghế đầu lên tới mái tôn. Khán đài đại diện cho Taira như bừng cháy, những cánh tay giơ cao, những bờ vai va vào nhau, tiếng vỗ tay dồn thành một trận mưa nhiệt huyết. Giữa cơn bão ấy, Yae chỉ khẽ thở dài, lòng bàn tay vuốt qua trán, lau đi vài giọt mồ hôi còn vương. Cô không cười, không la hét, chỉ đứng im, như thể mọi ánh mắt đang hướng về mình chỉ là gió thoảng qua.

``Quả không hổ danh là `vận động viên' có khác\dots\ Đúng là quái vật\dots'' Touka lẩm nhẩm, giọng khàn vì đã thốt ra quá nhiều câu hỏi không lời đáp. Lần này, cô không nghi ngờ nữa, chỉ còn sự thừa nhận trọn vẹn.

``Nhưng là quái vật đội ta được thuần hóa!'' Shion đáp lại, vai rung nhẹ, nụ cười nhếch lên như thể vừa được chia kẹo sau giờ học.

Ayame không để cảm xúc kéo dài thêm giây nào. Cô vỗ tay hai cái, âm thanh giòn tan như tiếng gõ mõ, ra lệnh cho dàn diễn viên rời sân khấu. Các đội lần lượt bước xuống cầu thang, giày kêu lộp cộp, tiếng nói cười râm ran như sóng vỗ vào bờ. Họ mang theo hơi nóng của trận bóng, mang theo mùi gỗ thấm mồ hôi, mang theo cả những cái nhìn chưa kịp nguôi ngoai.

Một màn diễn khác đang chờ, và khán giả đã bắt đầu dời bước, như thể chỉ cần quay đầu là có thể quên đi vừa rồi, nhưng không ai thực sự quên được cái cách quả bóng rơi xuống lưới, cái cách Yae đứng giữa sân, và cái cách Đội Đỏ đã thắng không chỉ trận bóng mà cả trái tim khán giả.

\begin{center}
  {\color{vol} Môn thi thứ sáu: \uline{Thi cổ vũ.}}
\end{center}

Khán giả còn đang nghẹn lời vì cú ném ba điểm cuối cùng của Yae thì tiếng loa đã réo lên, kéo mọi người khỏi dòng hồi tưởng. Dòng người chưa kịp giãn, vẫn chen nhau trên khán đài, đã thấy Hội trưởng Hội học sinh đạp hai bậc thang một lúc, míc-rô vừa bật, vừa cười như đã giấu kẹo suốt buổi.

``Chưa xong đâu nha\~{}''

Tiếng cô vừa dứt, khán đài như bị kéo căng lại. Mọi ánh mắt đổ dồn về phía sân trống vừa được trải nhựa mới. Cờ đỏ, cờ trắng rung rinh trên hàng rào, gió thổi qua mang theo mùi bụi và mùi mồ hôi còn vương lại từ trận đấu vật.

``Tiếp theo sẽ là nội dung trận chiến cổ vũ giữa hai đội!''

Ayame không dài dòng. Cô phổ biến luật bằng những câu ngắn gọn, xen vào đó vài cú huýt sáo khiến khán giả vừa cười vừa ghi nhớ. ì hội thao của Aogami có nhiều nội dung, phần thi này được rút gọn thành hai đội Đỏ và Trắng thi đối kháng trực tiếp, thay vì từng lớp riêng lẻ. Thời gian biểu diễn có hạn: hô khẩu hiệu, hát cổ vũ, vũ đạo đồng bộ, và một màn bất ngờ do chính đội tự chế. Điểm số không dành cho kẻ la to, mà cho kẻ kéo được cả khán đài về phía mình bằng cả trái tim.

``Nhớ nha,'' cô nàng nháy mắt, ``khán giả cũng đóng góp quan trọng vào điểm thi đấy!''

Cô tung đồng xu. Trên không trung, đồng năm mươi yên xoay tít một vòng, rơi xuống lòng bàn tay rồi vỗ chặt.

``Ngửa!''

Đội Trắng reo lên trước. Cô nàng tóc trắng mở bàn tay ra. ``Đội Trắng sẽ biểu diễn trước, yo\~{}''

Họ bước lên khoảng đất trống như kẻ mở màn một vở kịch. Phía Đỏ, Suzuki nắm chặt cờ, khẽ nghiêng đầu về phía đồng đội, ánh mắt nói thay lời: ``Nghe thấy tiếng reo của họ chưa? Ta sẽ làm lớn hơn thế.''

Khán đài dần im phăng phắc, chờ xem bên nào sẽ biến tiếng vỗ tay thành sóng thần trước.

\ct

Tiếng trống vang lên, không dồn dập như mưa rào mà đều đặn như nhịp tim khỏe mạnh với ba nhịp mở đầu, hai nhịp giữa, rồi dứt khoát một cái *BỤP* cuối cùng khiến cả khán đài như được đánh thức cùng lúc. Đội Trắng bước ra, bước chân trùng nhau đến mức chỉ nghe được những tiếng gót giày khêu gỗ sàn như một thanh kiếm rút khỏi vỏ. Hàng trước giậm gót, tạo nền đất rắn; hàng sau hô vang, giọng xếp tầng như sóng biển. Không ai vượt ai nửa nhịp, không ai chậm một khắc. Họ như những chiếc đồng hồ được vặn chung một núm.

``\eng{\textbf{White! White!}}'' tiếng hô không lớn nhưng lan tỏa, len lỏi qua từng hàng ghế. Giữa đội hình, bốn người bất ngờ xoay tròn đồng trục, tay giơ cao như bốn cánh quạt mở tung, chân dậm xuống tạo thành một vòng xoáy trắng khiến khán giả như bị hút vào tâm bão. Không chỉ là xoay, họ xoay đúng 270 độ, dừng sát mép vòng tròn, gót chạm gót, vai song song, tạo thành một đường thẳng tắp không thể soi bằng thước kẻ. Khán giả bị cuốn theo, vỗ tay theo nhịp, không cần ai bảo ai.

Bài hát cổ vũ vang lên, giọng hát của họ dính vào nhau như keo, không thể tách riêng. Đoạn đầu là đồng thanh mi-nôr, đoạn sau tách thành ba bè: cao, trung, trầm, rồi lại gộp vào một hợp âm hoàn hảo. Không ai bị lạc giữa dàn âm, không ai cố tình nổi bật, sự đồng đều ấy lại trở thành thứ vũ khí đáng sợ. Đến đoạn cao trào, cả đội xoay người, dựng thành chữ V thuần khiết: hai cột chân là hai thanh niên cao nhất, ba người ở giữa nhún đầu gối, tạo mũi tên, còn lại giơ tay tạo thành hai cánh bay. Người cầm cờ trắng xoay một vòng, mũi giày khẽ chạm đất đúng một nhịp duy nhất 0 không thừa, không thiếu. Rồi tất cả dừng lại, khẩu hiệu cuối gọn như dao cạo: ``Trắng chiến thắng! Trắng vô địch''

Khán đài nổ tung, nhưng không phải vì ồn ào mà vì cái cách họ im bặt một giây rồi bùng nổ ngay sau đó: một khoảnh khắc chân không âm thanh khiến tiếng reo sau đó như vỡ tung trần nhà.

Kanade chớp mắt, hai tay siết chặt thành lan can đến mức gân trắng bật: ``Không tì vết\dots\ họ không chỉ đẹp, họ còn đúng đến từng milimet.'' Đôi mắt cô long lanh như vừa thấy một bản giao hưởng không cần dàn nhạc. Kana ngồi bên cạnh, gật nhẹ, giọng trầm xuống: ``Từng chi tiết của đội họ ăn khớp, không dư một góc cạnh.'' Cô lẩm nhẩm đếm nhịp, ngón tay gõ trên đầu gối như đang ghi chép bí mật.

Kuon thở dài, không phải thất vọng mà như kẻ vừa phát hiện ra bức tường không có chỗ bám: ``Không lộ ra điểm yếu nào luôn. Cả đội như một khối bê tông đúc sẵn, không khe hở để luồn lách.'' Cậu nghiêng đầu, mắt vẫn dán theo đội Trắng, giọng khẽ: ``Muốn thắng họ, ta phải khiến bầu không khí phải náo nhiệt và sung sức hơn nữa.''

Kaoru cắn môi đến mức máu nhỏ đầu lưỡi: ``Xem chừng có vẻ khó rồi đây\dots\ khó không chỉ là điểm số, mà là cách họ khiến khán giả quên mất mình đang thi đấu.''

Ayame vỗ tay, mắt híp lại thành hai khe xếch: ``Hoàn hảo ghê luôn đó\~{} Đúng kiểu `chuẩn bài sách giáo khoa' ấy! Nhưng sách giáo khoa này lại có phần bài tập không lời giải, mà tự chính họ đã tìm ra lời giải!''

Cô ngoái về phía Kuon và Suzuki. ``Nghe này, Đội Đỏ,'' giọng cô thấp, nhưng míc-rô kéo âm thanh chạy dọc sân, mỗi chữ rơi vào lòng người ta như hạt chì nóng. Mình muốn cả khán đài phải quên mất đội Trắng từng tồn tại. Được không?''

``Được!'' tiếng đáp vang lên rời rạc, chưa đủ.

Ayame gật nhẹ, quay lưng về phía khán đài, míc-rô hạ thấp. Cô đếm thầm: ``Một\dots''

Tiếng trống lệch nhịp nổi lên, không từ đâu, như ai đó đập vào mặt trống bằng lòng bàn tay trần.

``Hai\dots''

Cánh tay cô vung lên, vạt áo vẽ một đường cung đỏ. Toàn đội Đỏ bật về phía trước một bước, gót giày đập xuống đúng phách ``ba,'' một tiếng *CHÁT* duy nhất, gọn đến mức khán giả giật mình như vừa bị tát.

``Bây giờ,'' Ayame quay lại, giọng vừa đủ cho míc-rô hút, ``\textbf{hãy khiến họ lạc nhịp và theo nhịp của chúng ta!}''

Cô buông míc-rô, bước lui về mép sân, để lại khoảng trống trung tâm như mồi lửa. Đội Đỏ không cần lệnh nữa; họ đã thành dòng dung nham đỏ chạy theo nhịp trống lệch; chân tự động nhún, tay tự động vung, và khán đài bắt đầu gõ chân theo lúc nhanh, lúc chậm, lúc ngừng hẳn rồi bùng vỡ.

Ayame đứng ngoài, hai tay khoanh, mắt lim dim: ``Đẹp lắm. Giờ thì\dots\ đốt sân lên đi.''

``\textbf{Đội Đỏ, chuẩn bị!}''

Tiếng Kuon hô từ chiếc loa tay như mồi lửa thả vào đám rơm khô, khiến cả khán đài bỗng nóng ran. Không khí vừa còn lơ đãng vì màn Trắng quá chỉn chu, giờ bị xé toạc bởi một nhịp trống lệch đúng nửa phách, không sai, mà cố tình đánh lạc hướng mọi chiếc đồng hồ đang chạy trong lòng người xem. Một nhịp ngắn, gãy, như ai đó vừa giật phăng khóa cửa căn phòng máy lạnh, nhường chỗ cho luồng gió mùa hè thổi vào bất ngờ.

``\textbf{ĐỎ!}''

Tiếng đáp trả vang lên, không chỉ từ một mà từ cả trăm cái cổ họng. Kuon đứng trên bậc cao nhất, loa cầm ngược, đầu nghiêng, như thể cậu đang không hô mà đang thả từng con chữ vào không trung rồi bắt chúng phải nổ tung. Cánh tay cậu vung lên không theo khuôn mẫu nào cả - lên cao, chéo xuống, vẽ một đường xiên mơ hồ - nhưng mỗi lần cậu dừng lại, cả khối Đỏ lập tức đóng băng đúng tư thế ấy, rồi lại bật ra như lò xo. Không có kẻ chỉ huy, chỉ có dòng điện âm thanh chạy từ đôi mắt cậu sang bả vai người bên cạnh, rồi lan ra toàn đội như sóng thần.

``\textbf{ĐỎ! TIẾN LÊN!}''

Akane giật chiếc loa tay, giọng vỡ òa, nhưng không hề chói tai, mà như chiếc loa công suất lớn vừa được đặt đúng chỗ, đủ khiến tim người ta rung theo. Bên cạnh, Suzuki bất ngờ bật ra một tiếng hô trầm ấm, vang dày, làm đất sân dường như rung thêm một nhịp. Hai giọng cao-thấp đan vào nhau tạo thành hợp âm không cần chỉnh, kéo theo cả dãy khán giả vốn đang thờ ơ nay bắt đầu gõ chân xuống nền gỗ.

Động tác bắt đầu như sóng biển đêm: không đồng đều, nhưng có nhịp thở riêng. Người thì xoay hông trễ nửa nhịp, người thì giậm chân mạnh hơn đúng một phách, nhưng tất cả đều hướng về trung tâm, một nơi mà Kuon tạo ra bằng cách\dots\ nhún vai. Đúng, chỉ một cái nhún vai, cả đội như được báo hiệu: ``Đừng quá khư khư khuôn khổ, giữ lại tinh thần.'' Những bước nhảy không còn là động tác nữa, mà thành dấu chấm than di động trên sân.

Đoạn chuyển đội hình đến, ai cũng chờ một màn lệch nhịp, và nó đến thật: hàng sau bước tới trước nửa bước, hàng trước lại lùi về sau, khiến không gian như bị xóc nảy. Nhưng Kuon không cứu vãn bằng lời kêu gọi thông thường, mà đơn giản đưa ngón tay trỏ lên trời, vẽ một vòng tròn nhỏ. Ngay lập tức, cả khối quay đồng trục 180 độ, vai chạm vai, gót chạm gót, tạo thành khối xoáy đỏ rực giữa sân. Từng động tác tưởng chừng không liên quan và có gì đó lạc lõng, nhưng chúng lại ăn khớp tới nhau đến lạ. Khán giả bên dưới không kịp hiểu chuyện gì, chỉ biết tay mình đã tự động vỗ theo nhịp trống vang lên từ đâu đó.

``\textbf{Nhịp! Giữ nhịp nào! Gâu gâu!}''

Ngay sau tiếng hô của nàng bờm Chó, cả nhóm bật ra một tiếng vỗ tay đúng chuẩn muộn một giây, tạo cảm giác như cả khán đài vừa được lôi kéo vào một vũ hội chung. Không còn khoảng cách giữa sân và khán đài nữa, tất cả đang nhún nhảy trên cùng một dòng điện màu đỏ, đỏ rực, đỏ tươi, đỏ đến nỗi cả màn đêm cũng phải chớp mắt.

Kuon giơ cao cánh tay trái, bàn tay nắm chặt như cầm một thanh gươm vô hình; đồng thời, cả khối Đỏ bật về phía trước một bước dài, gót chấm đất đúng nhịp ``đập'' cuối của trống. Khán giả bên dưới giật mình, rồi reo lên thích thú. Một bạn nam khối năm hai bật cười: ``Nghe như cả khối mình đang đáp máy bay vậy!''

Akane xoay người 270 độ, tay phải giơ cao míc-rô, tay trái kéo dải khăn đỏ phất lên tạo thành một đường cong lửa. Cô hét tiếp: ``\textbf{Đỏ, quyết tâm!}'' Toàn đội đáp lại bằng một bước nhảy đan chéo: người cao nhảy qua người thấp, người thấp lùi về sau, tạo thành hình chữ X nổi trên sân. Mấy cô bạn khối một ú ớ: ``Kìa, họ nhảy cao quá!''

Suzuki, vốn được xếp đứng ngoài cùng, bất ngờ lao vào giữa, lưng cong như cung, hai tay đẩy về phía trước đúng lúc cả đội gập người xuống. Động tác tạo thành một ``làn sóng'' đỏ dập xuống rồi bật ngược lên, mượt như lụa. Kuon nở nụ cười hiếm hoi, quay sang khán đài phía Nam, tiếp tục hét: ``\textbf{Còn bạn nữa, cùng hát lên nào!}''

Khán đài bên dưới bắt nhịp. Một nhóm học sinh từ các lớp của khối năm hai và năm ba bắt đầu vung tay, hò reo: ``\textbf{Đỏ! Đỏ! Đỏ!}'' Tiếng vỗ tay không đều nhưng rộn rã, xen lẫn tiếng cười khúc khích khi thấy một bạn nam đứng hàng hai bị lệch mũ, vội vàng chỉnh lại trong khi vẫn cố giữ nhịp.

Đỉnh điểm là phần dựng tháp người. Ở động tác này, cá thành viên của đội sẽ trèo lên ba tầng. Đó cũng là lúc này bài bất ngờ của đội mới được tung ra, không chỉ cả những khán giả đang ngồi xem, mà có khi còn cả chính các thành viên của đội đối thủ.

Tầng dưới cùng gồm bốn đàn anh khối ba, chân dang rộng bằng vai, cẳng tay nổi gân xanh đỡ lấy đôi giày của tầng giữa. Ở tầng giữa, ba đành chị khối hai, đầu gối khụy xuống bám cổ tay tầng dưới, thân hình nghiêng 15 độ để tạo độ dốc an toàn.  Tầng trên cùng là hai đàn em khối một, tay bám vai nhau, chân trụ trước rồi kéo thân lên cao như chim vỗ cánh, trong đó có Suzuki.

Khán đài nín thở. Kaigou đứng ngoài, mắt mở to trước cô em gái đang ở trên cao, đổ mồ hôi lạnh: ``Này, cái đó có trong kịch bản không vậy!?''

Kuon cũng khựng lại một nhịp, sững người khi thấy người con gái hai màu mắt đang tạo dáng trên cao.``\dots Suzuki, em đang làm cái quái gì--''

Suzuki giơ cao lá cờ đỏ, đứng thẳng trên đỉnh tháp. Cô đếm thầm ``một, hai'' rồi nhún gối, lao về phía trước. Một cú rơi tự do chỉ hơn hai mét nhưng kéo dài như cả thế kỷ. Giữa không trung, cô xoay người 360 độ, tay giữ chặt cán cờ để giữ thăng bằng; đôi giày trắng vẽ một vòng cung đỏ rực. Hai đàn anh ở dưới cùng đồng loạt quỳ gối, tay đan vào nhau tạo thành lưới đỡ. Cô nàng mắt hai màu đáp thẳng giữa lòng bàn tay họ, lực hút gân cốt khiến cả tháp rung nhẹ nhưng không đổ. Cô lập tức bật ngược lên, lộn thêm một vòng gọn gàng, tiếp đất bằng hai chân trong tiếng ồ vang dội của khán đài.

Kaigou chết lặng. Kuon cũng không nói được gì. Ở dưới khán đài, khán giả như nổ tung, tiếng reo vang trời, nhưng hai anh em họ chỉ đứng im, mắt dán vào cô em gái vừa hạ cánh nhẹ nhàng giữa biển lửa đỏ. Cô gái vốn trầm lặng đó, trước nay chỉ lầm lũi núp bóng sau anh chị, giờ lại tự tin đứng giữa trung tâm, thở gấp nhưng mỉm cười, như thể hoàn toàn khác hẳn với con người thường ngày.

Ngay khi phần biểu diễn tưởng chừng kết thúc, Kaoru giật lấy chiếc loa từ tay Kuon, giọng vỡ òa:

``\textbf{Chưa xong đâu! Mọi người cùng hát lên nào!}''

Cô hô to, mỗi câu đều kèm theo một nhịp vỗ tay, kéo cả khán đài vào điệp khúc:

\tabs[6cm] ``Đỏ! Đàn anh dựng nền!
\tabs[6cm] Đỏ! Đành chị nâng bước!
\tabs[6cm] Đỏ! Đàn em tung bay!
\tabs[6cm] \textbf{Đỏ! Ta cùng hướng tới chiến thắng!}''

Tiếng reo không chỉ từ khối 1-C hay khối năm nhất đơn thuần nữa. Các khối hai, ba, thậm chí là các giáo viên cũng bắt nhịp. Cờ đỏ rực giơ cao như rừng lửa, nhịp vỗ tay đập vào nhau đều đặn như trống trận.

Suzuki quay về phía khán giả, mắt sáng quắc, reo lên trong sự phấn khích cao độ: ``Chúng ta không cần hoàn hảo, chúng ta cần họ in mãi trong tâm thức, rằng chúng ta là ai!''

Nhịp trống dồn dập về những nhịp cuối, khán đài như bị níu lại một giây rồi vỡ òa. Tiếng vỗ tay không thành nhịp, cứ liên tiếp như mưa đầu mùa, kéo dài mãi không dứt. Cờ đỏ tung bay, ai cũng đứng, ai cũng hét, như thể cả trường vừa được cắm thêm một trái tim mới.

Ayame cười toe toét, hai bàn tay đập vào nhau thật to: ``Oa\~{} Đội Đỏ đúng là hỗn loạn nhưng cuốn thật đấy!''

Cô đưa míc-rô lên môi, mắt nửa khép, giọng trầm bổng như diễn viên trên sân khấu: ``Và bây giờ, điều mà cả hai đội đều mong \dots'' Khán đài nín thở.

Cô nàng bỗng khựng lại, nháy mắt một cái thật duyên: ``\dots sẽ được công bố sau nhé! Giữ lại chút hồi hộp cho buổi chiều cuối buổi mới vui!''

Dưới sân, vài thành viên Đội Đỏ thở dốc, vai buông xuống vì những bước nhảy chưa trọn vẹn. Kuon gỡ chiếc khăn đỏ lau mồ hôi, nhún vai cười: ``Ổn rồi. Chúng ta kéo được cả khán đài rồi còn gì.''

Suzuki gật đầu, ngực còn phập phồng: ``Đúng là vậy ạ\dots\ nếu có thua\dots\ thì cũng không hối hận. Mọi người cũng đã cố hết sức rồi mà.''

Tiếng reo vẫn vang vọng trên khán đài, lẫn trong gió chiều. Ở phía xa, sân chạy đã được kẻ vạch trắng tinh, những đường thẳng như mời gọi. Kuon quay lại, mắt tìm chiếc máy quay quen thuộc - công việc mà cậu làm giỏi nhất.

Bên ngoài, sân chạy đã được kẻ vạch trắng, những đường thẳng thẳng tắp như những câu hỏi đang chờ lời đáp. Gió thổi qua, phơi phới cờ đỏ, phơi phới cả mùi đất mới xới.

Và tới lúc này, món chính của sự kiện mới được mang ra.

\begin{center}
  {\color{vol} Môn thi thứ bảy: \uline{Chạy vượt chướng ngại vật.}}
\end{center}

Sân trường Aogami được trải đất màu đỏ nâu được nung cứng dưới nắng giao mùa, phẳng lì như tấm da trống khổng lồ vừa được căng ra. Nắng đầu giờ chiều nghiêng nghiêng, phủ lên bề mặt một lớp bụi vàng li ti, khiến mỗi vạch sơn trắng năm phân đều lấp lánh như chỉ bạc. Có tám làn, cách nhau đúng một khoảng bằng chiều rộng của hai bàn tay giơ ngang; vạch đầu tiên được kẻ lại bằng bút dạ quang để không bị nhòe, còn vạch cuối cùng, cũng là nơi cán đích được đính thêm dải ruy-băng đỏ chỉ chực vỡ ra khi ai đó chạm xuyên qua. Tổng độ dài của sân chạy rơi vào khoảng hai nghìn mét.

Chướng ngại vật bắt đầu ngay sau mũi vạch xuất phát: bốn rào cao tới hông người, thân gỗ sồi được bọc cao-su mềm, sơn trắng, cứ cách mét rưỡi lại nhô lên một cái như những khúc xương sườn trắng hếu của con cá voi khổng lồ. Tiếp đến là tấm lưới bò dài bảy mét: dây ni-lông đan mắt, căng đúng độ cao một gang tay so mặt đất; dưới lưới, cỏ được cắt sát gốc để không vướng đầu gối, nhưng vẫn chừa lại một lớp đất mềm mới xới, hơi ẩm, sẽ dính vào áo như hình con cua nếu lăn qua. Đoạn zig-zag gồm mười cọc cao 1,2 mét, đặt thành hai hàng so le; mỗi cọc quấn băng phản quang màu vàng chanh, để khi nắng xiên vào sẽ lóe lên như mắt mèo, bắt mắt người chạy vào buổi chiều.

Ở giữa sân, bể nước nông rộng $3 \times 5$ m, sâu vỏn vẹn phần tư mét, nhưng lòng được lót bạt xanh dương, nước được bơm đầy đến mép; mặt nước rung lên từng gợn nhỏ mỗi khi gió qua, phản lại những đám mây bàng bạc như màn hình LCD khổng lồ. Cuối cùng là đường nước rút sáu mươi mét, vạch trắng nối tiếp như những nốt nhạc trầm dồn dập, dẫn thẳng đến điểm dừng được đánh dấu bằng lá cờ đỏ tung bay, phía sau là tấm bạt in hình ngọn lửa, để khi ai đó về đích sẽ có cảm giác đang lao vào một đám cháy đang chờ.

Hai bên khán đài được dựng bằng giàn khung sắt sơn xám, bậc thang gỗ thông, chỗ ngồi rộng vừa đủ cho ba người bắt đầu rung lên từng hồi nhịp nhàng khi khán giả nhún chân. Các học sinh Đỏ - Trắng chen nhau, cờ tay phất phơ, tiếng xì xào như sợi chỉ đứt: ``Liệu có ngã không?'', ``Nước trong bể có bẩn không?'', ``Ai sẽ chạy đầu tiên?''. Mùi bụi khô, mùi sơn mới, mùi mồ hôi hòa vào nhau, khiến không khí nóng ran như vừa được bật lò vi sóng.

Giữa vòng tròn trung tâm, không phải Ayame, mà là phó hội trưởng Wata bước ra, míc-rô được gài lên ve áo, giọng cô vang rõ từng tiếng như hạt thủy tinh rơi xuống đá: ``Được rồi, mọi người chú ý. Môn kế tiếp: chạy vượt chướng ngại vật cá nhân. Mỗi người sẽ xuất phát khi có hiệu lệnh và vượt qua các chướng ngại vật như vượt rào, lướt lưới, zig-zag, lội qua bể, nước rút về đích. Ai về nhanh nhất sẽ ghi điểm cho đội. Tuy nhiên\dots''

Cô dừng lại, để khoảng lặng kéo dài đúng một nhịp tim, rồi nói tiếp, nhỏ hơn nhưng đủ để cả sân phải nín thở: ``\dots trong quá trình thi sẽ có những 'yếu tố bất ngờ'. Ban tổ chức không tiết lộ trước. Các bạn tự cảm, tự ứng biến, để từ đó có những phản xạ tốt nhất.''

Cô gập míc-rô, gật đầu nhẹ, rồi bước lui một bước, để lại phía sau là khoảng sân im phăng phắc đang chờ tiếng còi xuất phát đầu tiên.

Tiếng xì xào lan nhanh như sóng nhỏ.

``Ayame-san đâu mất rồi?''

``Sao phó hội trưởng lại đọc luật?''

``Chẳng phải hội trưởng phải là người phổ biến chúng sao?''

Những tiếng bàn luận đó chưa kịp nguôi, một giọng cất lên từ đường xuất phát, vừa quen vừa tinh quái. ``\textbf{Yo, da yo!}''

Mọi đầu quay lại. Ayame đã đứng ở đó, tay chống hông, nụ cười như vừa giấu kẹo dưới lưỡi. ``Lần này mình cũng nhảy vào cuộc chơi nha\~{}''

Khán đài như bị ném lửa. Tiếng kinh ngạc, tiếng la, tiếng cười đan nhau thành một mớ dây điện rối.

``Hội trưởng cũng thi luôn á!?''

``Đội Đỏ được hỗ trợ hết cỡ rồi!''

``Không công bằng!!!''

Cô chỉ phất tay, mắt híp lên, như thể cả thế giới đang làm quá lên chuyện cỏn con.
``Ồ, mấy chiêu `bất ngờ' kia\dots\ mình tự tay bày ra đó, coi như bài kiểm tra đột xuất nha\~{}''

Một cơn sóng nghi vấn dâng lên.

``Thế chẳng phải thiên vị sao?''

``Người viết đề lại vào phòng thi luôn?''

Ayame nghiêng đầu, nheo mắt, ngón tay dựng lên như cây thước kẻ lý lẽ. ``Thiên vị á? Đâu có\~{}'' Cô khẽ lắc ngón tay, điệu bộ như bày mẹo cho trò chơi board-game. ``Mình cũng chẳng biết bất ngờ sẽ nhảy ra lúc nào mà\~{}''

Khán đài im bặt một nhịp. ``\dots Hả?''

``Chính vì không ai biết trước, nên mới công bằng chứ,'' ccô cười tươi như vừa thả một quân bài phăng-teo xuống bàn.

Ở ngoài sân, Kuon đặt máy quay lên vai, ống kính khóa chặt vào Ayame. Cậu hơi chau mày, như vừa thấy một lá bài lật ngửa đúng lúc cuối ván.

``\textit{Quả nhiên\dots}''

Hình ảnh cô nàng Oni ở thế giới cũ hiện lên: nụ cười xoắn ốc, cái nháy mắt đầy mồi nhử, hay những lần vung kiếm katana đầy ngẫu hứng, và cả pha\dots\ làm nổ văn phòng đầy dầu, những trò mà chỉ nhớ lại thôi cũng đủ khiến cậu rùng mình.

``\textit{Không khác mấy, nhưng\dots}''

Một luồng gió nóng cuốn qua, Kuon thở ra, tiếng thở như thả một quân cờ xuống bàn.

``\textit{\dots phiên bản này xem chừng còn liều hơn cả bản gốc.}''

Vạch xuất phát đã lên màu trắng loáng. Các vận động viên dàn hàng, giày siết xuống đất như những mũi tên vừa được lắp vào cung. Yuki vươn vai, đôi mắt lấp lánh ánh đèn khán đài. Ayame đứng cạnh, chỉ xoay nhẹ cổ tay, dáng điệu như thể trận đấu chỉ là phần mở đầu cho một trò chơi lớn hơn nhiều.

\ct

Wata bước lên vạch xuất phát, míc-rô gài trên ve áo rung nhẹ theo nhịp thở. Cô không cần gõ loa, chỉ đưa tay phải lên cao, bàn tay mở ra như bật công tắc vô hình. Khán đài chợt im, tiếng gió chen vào khe nón tai nghe rõ mồn một.

``Sẵn sàng\dots''

Cánh tay cô vung ngang, kẻ một đường thẳng đỏ trên không. Tiếng còi xe đạp vang lên, chát chúa, cắt đứt mọi lùi bước. Người đầu tiên bật ra khỏi vạch vụt thành mũi tên - Shiromi Yae.

Ba bước đầu, cô đã nửa thân trước đối thủ. Không ai kịp reo, chỉ nghe tiếng gót giày đóng đinh vào đất, đều như máy đếm tiền.

Rồi, chỗ nửa đường, *XOẸT!* vòi nước bật ngang, phun mù trắng xóa. Yae không chậm lại mà thấp người, bàn chân trượt một gang, cơ thể nghiêng như con sóng lách qua khe đá. Khán đài bật ồ, tiếng kinh ngạc còn chưa thành hình đã bị cô bỏ lại phía sau.

Phía trước, hàng rào in bóng, nhưng bóng lệch, không đổ. Yae nhếch môi, nhận ra hình chiếu. Cô không rẽ, xuyên thẳng. Làn gió vẫn vỗ vào lưng áo, rào ảo tan như khói.

Một nhánh đường hiện ra: ngắn đầy bẫy, dài trống trải. Cô chọn ngắn, mắt không chớp. Rào giả đổ rầm sau lưng, cô đã về đích, gót in vạch trắng như dấu chấm than cuối câu.

``\textbf{Đội Đỏ, dẫn trước!}'' tiếng hô vỡ òa, lan tỏa từ khối một lên khối ba, như sóng leo thang.

Yuki bước vào đường chạy như một chiếc đồng hồ vừa lên dây, từng bước chân dài đều đặn, không vội vã cũng không chậm chạp. Không có pha bứt tốc gây choáng như Yae, cô chỉ lặng lẽ duy trì tốc độ ổn định, để khoảng cách giữa mình và đối thủ luôn như một khoảng lặng an toàn. Đằng xa, tiếng reo của khán đài chỉ như gió thoảng qua tai, cô mải cảm nhận nhịp tim mình đập đều như máy đếm thời gian.

Đến khúc zig-zag, bỗng một tiếng *CẠCH* khô khốc vang lên, và cả hàng cọc lập tức xoay ngang, biến thành mê cung động. Yuki khẽ giật mình, gót chân chậm lại đúng nửa nhịp. Cô có thể chọn dừng hẳn để quan sát, nhưng thay vào đó lại mỉm cười như người vừa gặp trò đùa của bạn cũ: ``Hay đấy!'' Cô hạ trọng tâm, bước ngoặt chậm rãi, để đôi chân tìm lại điểm cân bằng mới, rồi lại trở về nhịp cũ, như thể chưa từng có cú xoay bất ngờ kia.

Khúc khích.

Tiếng cười trẻ con bật ra giữa sân, trong trẻo nhưng lạnh ngắt, như hơi nước bốc lên từ mặt hồ đêm đông. Cô nàng cảm thấy nổi da gà dọc cánh tay, cô dừng lại một giây, mắt tìm kiếm nguồn âm thanh.

``\dots Cái gì vậy?''

Không thấy bóng dáng ai, chỉ có cái bóng của mình in trên đất, méo mó bởi ánh nắng xiên. Cô lắc đầu, ép bản thân quay lại cuộc đua, nhưng tiếng cười vẫn lẩn khuất trong gió, như lời thì thầm của kẻ đi trước mở đường.

Ở ngoài đường chạy, Kuon đang chỉnh khung hình, bỗng cảm thấy có gì đó lạnh chạm vào tai nghe. Cậu khựng lại, nhíu mày: ``\vie{\dots Tiếng gì vậy?}''

Chiếc đồng hồ thông minh rung nhẹ, giọng Game vang lên bình thản: ``\eng{\dots Âm thanh đó có vẻ như không nằm trong danh sách hiệu ứng của ban tổ chức.}'' Kuon siết chặt máy quay, mắt không rời khỏi Yuki, như thể chỉ cần chớp mắt cô sẽ biến mất cùng tiếng cười kia.

``\vie{\dots Vậy là chính \emph{chúng} rồi.}''

Khán đài bỗng rộ lên những tiếng xì xào như lá khô xào trong gió lửa. Tiếng cười trẻ con vừa vang lên quá rõ, quá gần, không còn là âm thanh loa phóng thích nữa mà như ai đó đang đứng ngay sau lưng từng khán giả, thổi vành tai họ một hơi lạnh. Một bạn nữ khối một ôm chặt cờ tay, môi run: `` Mình\dots\ mình vừa nghe thấy bé gái gọi bằng chị\dots''

Bên cạnh, thầy giáo thể dục giật mình đánh rơi còi, mắt trợn tròn: ``Cái quái\dots\ rõ ràng chúng ta đã kiểm tra kỹ mà\dots'' Tiếng xì xào lan thành sóng, từng hàng ghế rung nhẹ như có ai gõ gót lên gỗ, nhịp đều đặn - bốn, năm, sáu - đúng nhịp bước chân trẻ nhỏ chơi nhảy lò cò.

Kaigou đang đứng ngoài hành lang tầng hai, tay siết lan can đến trắng xương. Cô không nghe bằng tai nữa, mà bằng cả làn da: tiếng cười ấy thấm qua kẽ tóc, bám vào gáy, len xuống xương sống. Nó không phải của một đứa trẻ nào cả, mà của tất cả những linh hồn nhỏ bé cô từng gặp trong thị trấn này kể từ ngày nó được tái sinh. Những đứa trẻ mắt tròn xoe len lỏi giữa dòng người tấp nập, những bóng hình xuất hiện dưới sân trường, hay việc chúng lù lù xuất hiện ở trong rạp chiếu\dots\ Tất cả hòa vào một âm duy nhất, reo lên như thể vừa tìm được lối về.

Kaigou cắn chặt răng, mồ hôi lạnh túa ra: ``Các em muốn nói gì vậy\dots?'' Gió thổi qua hành lang, mang theo mùi hoa sữa chưa kịp nở, và trong mùi hoa ấy, cô nghe rõ một lời thì thầm không thành tiếng: ``Mau tới đây chơi vói chúng em đi\dots\ Hãy giúp chúng em được giải thoát\dots''

Quay lại đường đua Yuki hít một hơi thật sâu, cảm nhận mùi cỏ non và mồ hôi lẫn vào nhau. Cô chọn đường dài, nơi không có bóng cọc xoay, không có tiếng cười lạnh. Từng bước chân lại trở về đều đặn, nhịp nhàng, như chiếc máy chạy bộ vừa được reset. Khoảng cách giữa cô và đối thủ không thay đổi, nhưng cô biết mình vừa vượt qua không chỉ chướng ngại vật bằng gỗ, mà cả thứ gì đó vô hình đang lảng vảng trên sân. Khi cô chạm vạch giao tiếp, đôi mắt vẫn còn đọng lại chút nghi hoặc, nhưng nụ cười đã trở lại: cô đã giữ lửa cho Đội Đỏ, không để ai vượt qua.

Tiếp theo, đến lượt Hội trưởng hội học sinh bước vào đường chạy. Vừa nghe tiếng còi, ngay từ bước đầu tiên, cô đã\dots\ đánh một bước chậm nửa nhịp, khiến cả khán đài ``dậy sóng'' bằng tiếng ngạc nhiên thốt lên.

Nhưng chỉ trong chớp mắt, cô nghiêng người, đổi hướng, lách qua tấm rào giả đang đổ ập xuống đúng chỗ cô vừa đứng. Khán giả còn chưa kịp hết ``ồ'', Ayame đã trở lại nhịp chạy, như thể cô có bản đồ bẫy in sẵn trong đầu.

Kuon nhíu mày sau ống kính: ``\dots Vậy là không phải đoán trước dữ chưa?''

Đoạn zig-zag, các cọc xoay ngang dọc bất thình lình. Không thắng gấp, cô bật tung người, nhảy vọt qua một trụ, lướt ván qua trụ khác, mũi giày chỉ chạm đất đúng chỗ trống duy nhất. Cả khúc cua gỗ trở thành sân khấu parkour thu nhỏ, mà nhân vật chính vẫn thở đều, tóc bay lất phất.

Kanade chép miệng: ``Thật\dots\ không tưởng.'' Koishin thì nhướng mày nghi ngờ: ``Người tự đưa ra luật chơi thì biết cũng là bình thường thôi, đừng khen quá làm gì.''

Tới ngã ba đường, Ayame bất ngờ phanh lại. Cô quay sang khán đài, nháy mắt như là người chơi chính trong trò chơi: ``Đường nào bây giờ ta? Ngắn hay dài đây mọi người?''

Tiếng hô tán loạn. Người la ``Đi đường ngắn đi!'', kẻ gào ``Không, đường dài an toàn hơn!''. Cô nhún vai, chọn đường ngắn, lao vút. Vừa chạy được hai bước, vòi nước phụt thẳng vào mặt. Cô trượt chân, suýt ngã, nhưng kịp xoay hông, giữ thăng bằng bằng một cười khúc khích: ``Ui cha, trơn quá\~{}''

Rồi cô phi tiếp, vượt qua vạch đích trong tiếng reo hò vỡ òa xen lẫn tiếng thở dài bất lực của những người dựng bẫy. Không thể không hoài nghi sao được khi chính cô là người đã đề xuất ra ý tưởng quái gì, và chính cô cũng đã hoàn thành phần thi của mình mà không gặp phải bất trắc gì quá nghiêm trọng.

Lượt cuối của Đội Đỏ gọi tên Hikawa Kaigou. Đối thủ kề vai cô là gã khổng lồ từng bị cô hất lên không trung ở môn kéo co. Nhìn thấy người đã đả bại mình, hắn nheo mắt, khẽ nhếch môi:  ``Sức khỏe của mày có vẻ rất tốt đấy. Nhưng tốc độ và phản xạ thì chắc mày chưa chắc nuốt được tao đâu.''

Kaigou chỉ thở dài, không đáp. Cô lầm bầm trong miệng, chẳng thèm nhìn vào hắn ta: ``\textit{\vie{Tại sao mình lại phải gặp cái tên này nữa chứ\dots\ Xui đến thế là cùng\dots}}''

Tiếng còi *TOÉT* xé không khí.

Hai bóng đồng thời bật về phía trước. Gã kia như xe tăng gầm rú, mỗi bước dài bằng hai bước người thường, còn cô nàng ắmt xanh thì lặng lẽ ép mình xuống thấp, gót chân chỉ vừa chạm cỏ đã bật lên, gần như trượt sát mặt đất.

Đoạn bể nước cạn chỉ đủ ngập mắt cá chân nhưng phủ một lớp bạt trơn bóng. Gã to con đạp phải, đôi giày thể thao ma sát kém, *xoẹt* một cái, thân hình loạng choạng. Kaigou đã nghiêng người từ trước, trọng tâm thấp, hai cánh tay dang ra như lá chắn; cô lướt qua khỏi mặt nước chỉ bằng ba nhịp chạy không dính một giọt.

Tiếp theo là dải hình chiếu 3D biến ảo: rào gỗ ảo, bóng ma cọc tiêu, và những tường lửa giả. Gã chệnh choạng dừng lại, mắt tìm điểm mù, quýnh quáng nhì những chướng ngại vật lỗn lộn thực ảo.

Kaigou không chậm một giây. Cơ thể cô nhớ lại mê cung tối tăm từng chạy suốt đêm để thoát Shino (hồi cô chưa công nhận) - những hố ``tử thần'' ló từ mặt đất, những bức tường mê cung xen lẫn gạch và ảo ảnh, những màn bóng đen truy đuổi. Ký ức đó giờ trở thành bản năng: cô lách, xoay, trượt, bật, mỗi động tác vừa vặn như đã được lập trình.

Ống kính Kuon rung nhẹ. Trong tai nghe, Game thì thầm: ``\eng{Chuỗi chuyển động này\dots\ trùng khớp với `vòng lặp truy đuổi'.} của các cậu.''

Khán đài bắt đầu ồn ào. Người ta không hiểu cô gái thanh mảnh kia đang làm phép gì, chỉ thấy bóng đỏ loang loáng vượt qua từng lớp bẫy như thể đường chạy được vẽ riêng cho cô.

Đoạn zig-zag, mười cọc gỗ xoay ngang dọc bất thình lình. Gã to con vấp phải một cái, đầu gối chạm đất, mất nhịp. Kaigou không thèm nhìu lại; cô nhún một cái, đầu ngón chân chạm đỉnh cọc, cất mình bay qua, mềm như múa.

Rào ảo cuối cùng đổ ập xuống trước mặt. Cô không rẽ, không dừng, chỉ đơn giản đạp thẳng vào trung tâm tấm bạt giả. *RỘP!* Tấm rào gục ngã, khói bụi bốc lên. Cô lao ra khỏi đám mây đất cát, vạch đích đã ở ngay trước mắt.

Phía sau, gã khổng lồ còn đang vật lộn với lưới bò, hai tay quờ quạng như gấu mắc bẫy. Kaigou gài bàn tay lên đầu gối, nhẹ nhàng bước qua dải ruy-băng đỏ. Còi hết giờ réo vang.

Khán đài nín thở một nhịp, rồn vỡ òa với cô gái nói ít làm nhiều kia, ngỡ ngàng với màn trình diễn vừa rồi: ``Cho\dots\ cho hít khói luôn kìa\dots'' một học sinh khối năm nhất lắp bắp, mắt tròn xoe.

Cô nàng không nghe thấy gì nữa. Tim cô đập chậm lại, mồ hôi mặn chảy vào khóe môi. Cô biết mình vừa chạy không chỉ vì một trận thi đấu; cô vừa chạy để xua đuổi cái bóng đã đuổi theo mình suốt bao tháng ngày. Và lần này, cái bóng ấy đã bị bỏ lại phía sau, cùng với gã khổng lồ còn đang ngồi thở dốc giữa đám lưới nhàu nát.

\ct

Yae bứt tốc như tên lửa, Yuki giữ nhịp đều như máy đồng hồ, Ayame bất ngờ như lá bài tẩy, còn Kaigou dập đối thủ bằng cú knock-out cuối cùng đầy lạnh lùng. Bốn mảnh ghép ấy lắp vào nhau, tạo thành chiếc thắng tất yếu cho Đội Đỏ. Khán đài vỡ òa, tiếng vỗ tay dập dồn như mưa đầu mùa.

Ở mép sân, gã khổng lồ vừa thua vật lộn gỡ lưới, mặt đỏ bừng, mắt tìm kiếm kẻ đáng trách. Ánh nhìn hắn lướt qua Kaigou, rồi dừng lại ở Kuon lâu hơn một chút, người đang cầm máy quay. Im lặng bốc lên giữa hai bên, nặng hơn cả khói bụi vừa bị giày đạp bay.

Cậu Tổng không đáp trả bằng lời, chỉ khẽ nghiêng ống kính, nhưng trong đầu cậu đã gióng lên hồi chuông cảnh báo. Game gằn giọng, lên tiếng: ``\eng{Cậu nên đề phòng bọn chúng.}'' Cậu gật, hầu như không tiếng động: ``\eng{Đúng với suy nghĩ của tôi. Chúng sẽ trả đũa, không hôm nay thì ngày mai, không ở sân thì ở nơi khác.}''

Gió thổi qua, cuốn theo mùi bụi và mồ hôi. Trong khoảnh khắc trước khi loa phát thanh cất tiếng, một tiếng cười trẻ con nhỏ xíu, như sợi chỉ thừng, vương lại giữa không gian, nghe rõ đến mức khiến vai Kuon giật nhẹ.

``\vie{Giờ thì\dots\ chuẩn bị cho môn thi cuối cùng nào.}''

\begin{center}
  {\color{vol} Môn thi thứ tám: \uline{Chạy tiếp sức.}}
\end{center}

Khán đài tại sân chạy của trường Aogami như chảo lửa thổi bùng lên lần cuối. Mười tiếng reo không kịp nghỉ, vang dội thành một dòng chảy âm thanh không ngớt, cuốn theo bụi đất, mùi sơn mới và cả mùi mồ hôi đã ngấm vào gió chiều. Không ai còn ngồi yên; họ đứng, họ nhảy, họ đập hai hàng ghế gỗ rền vang như trống trận. Tất cả đều hiểu: đây là hiệp cuối, nơi chỉ còn một cái tên được khắc vào cúp.

Ngay chấm phát súng, Ayame và Wata xuất hiện. Hai bóng áo đỏ-trắng đứng sát nhau, một nóng một lạnh, như đầu máy và toa phanh của cùng một đoàn tàu.

Ayame phóng giọng qua míc-rô, lời còn chưa tới tai khán giả đã bị tiếng vỗ tay cắt ngang: ``Đây là chạy tiếp sức! Mỗi đội sẽ có năm lượt thi đấu, mỗi lượt gồm bốn người. Người chạy sẽ chuyền gậy cho người tiếp theo trong khu vực quy định. Nếu làm rơi là mất thời gian đấy nhé\~{}!''

Wata nghiêng đầu, bổ sung không chút vội: ``Đội nào hoàn thành đủ năm lượt với tổng thời gian nhanh hơn sẽ giành chiến thắng.''

Dòng người xì xào như sóng vỗ vào bờ, ôn lại bốn lượt vừa qua.

Lượt một: Đội Đỏ vượt lên nhờ cú bứt tốc của cô nữ sinh tóc tém.

lượt hai: Đội Trắng san bằng khi chàng ``xe tăng'' nhảy rào bằng cả hai chân.

Lượt ba: một cú trượt chân giữa bể nước khiến đội Trắng nắm bắt thời cơ và giành chiến thắng.

Lượt bốn: bản lề cuối cùng nghiêng về phía áo đỏ chỉ bằng nửa bước chân.

Cả bốn đang xếp hạng ngang bằng: hai lần nở nụ cười, hai lần nghiến răng, cân bằng như đồng hồ treo tường chờ nhịp cuối cùng. Hai mươi người kia đã rút lui khỏi đường chạy, đem lại cho khán giả những cuộc rượt đuổi nghẹt thở và hấp dẫn.

Suzuki khoanh tay, mắt vẫn dán vào đường chạy. Cô thì thầm với Kana đứng cạnh: ``Mình hy vọng là Onee-sama và các bạn sẽ giành chiến thắng trong vòng thi này.''

Kana cười hì hì, vỗ vai cô bạn: ``Yên tâm, lần này còn có `bộ tứ thần ' của lớp 1-C mà. Gậy có muốn rơi cũng phải xin phép họ trước đã.''

Phía sau hai người, Koishin gõ gõ giày xuống nền gỗ, mặt đỏ bừng lên, má phụng phịu: ``Hừ, tôi chỉ đến xem chứ có cổ vũ gì đâu. Dù gì họ cũng thắng lớn rồi mà.''

``Không cần phải giấu đâu! Tớ biết là cậu cũng muốn đội nhà giành chiến thắng mà,'' cô nàng tóc tím nói, thóc nhẹ vào lườn của cô nàng tóc vàng, càng làm cho Koishin đỏ mặt hơn.

Sân khấu của môn thi cuối cùng giờ đây được nhường cho bộ tứ cuối được tinh chọn (dù khá là nhanh gọn) đại diện cho lớp đứng bởi Tsukishita Yuko: Yae, Yuki, Shiina và Kaigou. Tiếng loa rít lên, MC gọi tên bốn gương mặt cuối cùng.  Họ xếp hàng, gậy tiếp sức đỏ tươi nằm gọn trong tay Yae, như thanh kiếm chờ vị tướng ra trận.

Bốn cái tên đó, chỉ cần nghe qua thôi cũng đủ khiến người khác phải dè chừng.

Trong khi Đội Đỏ đang nóng lòng chờ đón bốn cái tên quen thuộc, bên kia sân, bầu không khí của Đội Trắng lại đặc quánh như mùi thuốc súng chưa kịp tan. Ba kẻ từng khiến khán giả phải nhăn mặt suốt buổi sáng nay lại xuất hiện, đứng sát nhau thành một hàng dọc đầy ám ảnh: gã bị cho hít khói lúc vượt rào, tên ``đội trưởng bóng né'' từng chặn đầu Kuon, và kẻ đẩy sau rồi giả vở không biết gì trong trận tamaire. Họ không cần nói gì, nhưng sát khí từ chúng cũng đủ khiến người ta nhớ lại cảm giác bị bóp nghẹt cổ họng. Đứng cạnh ba bóng ma ấy, cậu học sinh thứ tư trông như một con cừu lạc đàn, mắt lia xuống đất, chân khẽ lùi về sau mỗi khi tiếng cười khẩy của đồng đội vang lên.

Không khí bỗng nặng hơn. Mây chiều như đè thấp xuống sân, tiếng reo hò tắt ngóm giữa cổ họng. Ở vạch xuất phát cuối, Kaigou đứng thẳng băng, hai bàn tay nắm chặt đến mức gân xanh nổi chằng chịt dọc mu bàn tay. Không phải nhịp tim, cũng chẳng phải khát khao chiến thắng, đôi mắt cô chỉ hướng về phía ba kẻ kia, như thể cả thế giới chỉ còn lại bốn bức tường vô hình giữa họ và cô. Cơn giận không chỉ là làn khói bốc lên rồi tan, nó là than hồng cháy âm ỉ từ sáng, giờ đã đỏ lòm, dí sát vào lồng ngực.

Tên đeo băng đội trưởng của đối thủ nhếch môi, giọng vừa đủ để cô nghe: ``Nhìn gì mà như muốn ăn tươi nuốt sống vậy? Đừng quên luật đấy, cô bé. Chỉ cần một cú đấm cũng đủ bị ban tổ chức lôi cổ ra ngoài rồi.'' Hắn cố bắt chước lại giọng hồn nhiên của Ayame, pha thêm vài động tác tay vung vẩy, khiến lời cảnh cáo nghe như lời chế nhạo: `` `Làm bị thương người khác là out ngay đó\~{}' nghe quen tai chưa?''

Im lặng. Khán đài ồn ào nhưng với Kaigou, mọi thứ như bị bóp nghẹt. Một bước thôi, chỉ cần cô nhích thêm một bước, cái cằm của hắn sẽ gặp gót giày cô. Nhưng cô không nhúc nhích. Cơn giận bị bứt ra khỏi xương, nhai nát trong miệng, rồi nuốt xuống cổ họng đắng chát. Cô quay đi, hàm răng vẫn kẹp chặt, như thể vừa khóa cửa một con thú dữ vào trong lồng ngực.

Kuon đứng ở bên ngoài đường pitch, tay cầm chiếc máy quay, lặng lẽ chứng kiến mọi thứ. Không một tiếng động, chỉ có hơi thở nhẹ thoát ra vì tính khí nóng nảy của cô nàng.

``Kaigou-chan.''

Tiếng gọi của cậu khiến cô nàng như bị ai giật dây, vai cứng đờ.

``\vie{Lại đây chút.}''

Kaigou quay lại, bước tới, khoanh tay, giọng vẫn không giấu được vẻ khó chịu: ``\vie{Có gì nói nhanh.}''

Cậu nhìn thẳng, giọng gọn, không màu mè: ``\vie{Tôi vừa hỏi lại barem điểm từ ban tổ chức rồi.}''

Cậu giơ tay, đếm như đang đọc bảng cửu chương: ``\vie{Tamaire và kéo co chiếm 5\% tổng số điểm mỗi môn. Bóng né, đấu vật, bóng rổ là 10\%. Thi cổ vũ 15\%, chạy vượt chướng ngại: 20\%, còn chạy tiếp sức\dots\ 25\%.}''

Kaigou hơi nhíu mày. Cô không phải kiểu người thích tính toán, nhưng chỉ cần nghe qua cũng đủ hiểu môn này là then chốt. ``\vie{Vậy tổng lại thì sao?}''

Kuon không trả lời ngay. Cậu nhìn cô, chờ cô tự nói.

Kaigou thở nhẹ, rồi tự tính nhẩm.

``\vie{Chúng ta thắng tamaire\dots\ nhưng thua một vài lượt nhỏ, tổng thể vẫn tính là thắng.}''

``\vie{Chúng ta thắng kéo co,}'' cậu Tổng tiếp lời.

``\vie{Bóng né, ta thắng.}''

``\vie{Đấu vật, ta thắng phần nữ và hỗn hợp.}''

``\vie{Bóng rổ\dots\ cũng nghiêng về phía chúng ta.}''

``\vie{Chạy vượt chướng ngại vật, thắng.}''

Cô dừng lại một chút, rồi kết luận:

``\vie{\dots Chưa tính thi cổ vũ và chạy tiếp sức\dots\ thì Đội Đỏ đã có khoảng 60\% tổng điểm rồi.}''

Kuon khẽ gật đầu. Chính xác.

Cậu nói tiếp, giọng trầm ổn như máy đếm kim: ``\vie{Thử xem kịch bản tệ nhất nhé. Đội Trắng bứt phá ở chạy tiếp sức, tức là giành được hai mươi lăm phần trăm. Rồi họ còn thắng luôn phần cổ vũ, làmười lăm phần trăm nữa.}''

Mỗi con số được cậu nhả ra chậm rãi, như thả từng viên bi thủy tinh xuống đáy ly. ``\vie{Cộng lại, bốn mươi.}''

Kaigou khựng nửa nhịp, lông mày nhướn cao. Kuon khép lại bài toán:  ``\vie{Sáu mươi so bốn mươi. Chúng ta vẫn nhiều hơn.}''

Khán đài rền vang tiếng reo, nhưng quanh hai người như có tấm màn chắn vô hình, tách họ ra khỏi mọi ồn ào.

Cô nàng mắt xanh thở ra, vai buông lỏng hơn chút. ``\vie{Vậy là\dots\ dù có thua hiệp cuối này\dots}''

``\vie{Thì cúp vẫn về tay Đội Đỏ,}'' Kuon kết thúc, nhún vai như vừa gạt bỏ một hạt bụi trên ve áo. Cậu dừng một giây, ánh nhìn sắc lại. ``\vie{Nhưng đừng vì thế mà chủ quan.}''

Kaigou cười khẽ, khàn đặc: ``\vie{Cậu tưởng tôi định buông xuôi à?}''

Kuon không đáp trực tiếp, chỉ để lại một câu nhỏ như đóng dấu: ``\vie{Cô thuộc dạng xả vai trước, nghĩ sau. Còn lạ gì.}''

``\vie{\dots Tôi không phủ nhận,}'' cô nàng gầm gừ, mắt đảo ra chỗ khác.

Một hồi im lặng kéo dài. Cô nàng mắt xanh quay mắt về phía đội Trắng. Ba bóng đen kia vẫn dán chặt vào họ, ánh mắt như mũi dao cùn, không đâm nhưng cứa dai dẳng.

Cô khẽ nói, giọng chỉ đủ hai người nghe: ``\vie{Chúng sẽ chơi bẩn.}'' Kuon gật, không chút ngập ngừng: ``\vie{Tôi cũng nghĩ vậy.}''

Cậu nghiêng đầu, khóe môi lạnh tanh: ``\vie{Nhưng vải thưa không che được mắt Thánh. Dù có chơi trò tinh vi đến đâu, sớm muộn gì cũng bị phát giác thôi.}''

Kaigou liếc xéo, cười nhẹ:  ``\vie{Câu đó nghe quen quá nhể?}''

``\vie{Ừ,}'' Kuon đáp, mắt vẫn dõi xa, ``\vie{nhưng vẫn đúng.}''

Kuon quay lại nhìn đường chạy, rồi nói tiếp:

``\vie{Về thứ tự, nghe đây. Theo tôi, cứ để Yuki-san chạy đầu. Cô ấy ổn định, không bị cuốn theo tâm lý. Mở màn an toàn. Yae-san chạy thứ hai. Tốc độ và phản xạ tốt, xử lý tình huống bất ngờ ổn.}''

Cô nàng ``song sinh'' gật nhẹ. Không có gì để phản đối.

``\vie{Shiina-san thứ ba. Cô ấy đủ mạnh để kéo lại nếu bị tụt, và đủ lì để giữ nhịp nếu đang dẫn.}''

Cậu Tổng dừng lại một nhịp, rồi nhìn thẳng vào Kaigou: ``\vie{Cô sẽ là người chốt hạ.}''

Kaigou không hỏi ``tại sao''. Cô hiểu.

``\vie{Trong ba người chạy vượt chướng ngại vật vừa rồi, cô nhanh nhất. Và cũng là người khó bị lung lay nhất.}''

Cậu nói thêm, giọng thấp xuống một chút: ``\vie{Nhưng nhớ một điều. Đừng để mấy lời khiêu khích kia ảnh hưởng. Coi chúng như những lời gió bay thôi.}''

Kaigou nhếch môi: ``\vie{Tôi tưởng cậu phải hiểu rõ tôi rồi chứ.}'' Kuon đáp ngay: ``\vie{Chính vì hiểu nên tôi mới nhắc.}''

Một giây. Hai giây. Rồi Kaigou thở ra, xoa nhẹ gáy: ``\vie{\dots Được rồi.}''

Kuon gật đầu: ``\vie{Thế là đủ.}'' Cậu lùi lại nửa bước, giọng trở về bình thường: ``\vie{Dù sao thì,}'' một nụ cười rất nhẹ thoáng qua khóe môi cậu, ``\vie{tôi không coi trọng thắng thua đâu. Thứ quan trọng hơn là tinh thần đồng đội của mọi người.}''

Cậu nhìn về phía nhóm 1-C, những người đang khởi động chuẩn bị cho môn thi cuối cùng này. ``Cố gắng lên nhé.''

Kaigou không nói thêm lời nào. Cô xoay người, bước về phía nhóm, vai thẳng, gáy cứng. Nhưng chỉ được nửa đường, ba tên đội Trắng - kẻ đeo băng đội trưởng cùng hai gã luôn miệng cười khẩy - bất ngờ lướt qua, cố ý đụng vai. Chúng không đủ gần để nghe rõ cuộc trao đổi, chỉ kịp bắt được vài mảnh câu ``\dots sáu mươi\dots\ bốn mươi\dots''. Tên đeo băng nheo mắt, quay sang bạn, giọng vừa đủ để Kaigou nghe thấy: ``Nghe như chúng nó đang tính toán xem bao nhiêu phần trăm thua trận đấy nhỉ!''

Cả bọn phá lên cười, tiếng cười vỗ vào không gian như mảnh thủy tinh vỡ. Cô không dừng bước, chỉ hơi nghiêng đầu, để ánh mắt xanh lướt qua chúng như lưỡi dao cùn. Không có lời đáp, không có cả ánh nhìn dài, chỉ một khoảnh khắc lạnh toát, đủ khiến tiếng cười tắt ngóm trong cổ họng kẻ đứng giữa. Cô bước tiếp, để lại sau lưng ba cái bóng đang cố gắng giữ vẻ ngạo mạn nhưng đã lỡ nhịp.

Yuki, Yae và Shiina đồng loạt ngoái lại. Không ai hỏi ``sao lâu thế'', chỉ thấy Kaigou dừng ngay trước mặt, khoanh tay, giọng đã trở về độ lạnh thường trực: ``Thứ tự giữ nguyên. Yuki-san, Yae-san, Shiina-san, rồi đến tôi.''

Cô nàng ``Xã hội đen'' khẽ gật, mắt vẫn dán vào bảng tên đối thủ, nói vỏn vẹn đúng một từ: ``Ổn.''

Cô nàng mắt vàng chỉnh lại tất, siết nhẹ bàn tay trái như khóa chặt nhịp tim: ``\dots Tớ sẽ mở màn thật tốt.''

Cô nàng tóc xám cười nhẹ, khoe cương đều trắng bóng: ``Để lại phần khó cho bọn này.''

Kaigou liếc sang đội Trắng lần cuối - ba kẻ vừa bị cô bỏ lại đang cố nói gì đó để lấy lại thể diện - rồi quay đi, giọng thấp nhưng rõ: ``Cẩn thận bọn chúng.''

Shiina nhếch môi, tay đập nhẹ gậy tiếp sức vào lòng bàn tay: ``Thách chúng nó làm vậy đấy.''

Trên vạch xuất phát, bầu không khí đã nén lại như lò xo. Tám vận động viên bước vào ô riêng, gót chân chạm vạch trắng, đầu gối hơi khuỵu, cánh tay đong đưa tìm nhịp. Khán đài bỗng im phăng phắc, chỉ còn nghe tiếng gió kéo qua cờ phường, lách tách như giật thắt.

Một giây\dots\ hai giây\dots\ và\dots

Tiếng còi *\textbf{TOÉT}* vang lên, sắc như dao cạo, xé rách màng căng thẳng. Yuki của Đội Đỏ và gã trùm của Đội Trắng bật ra như hai mũi tên rời cung, bước đầu tiên nện cỏ cùng lúc, gót bùn bắn lên tung tóe. Cả sân như nổ tung trong tiếng reo, nhưng hai bóng dẫn đầu đã lao xa, chỉ để lại trên đất dấu giày in nghiêng nghiêng, hơi nóng bốc lên từng nhịp.

Yuki không phải người nhanh nhất lớp, nhưng cô là người ổn định nhất. Nhịp chạy của cô như guốc đồng hồ, đều đặn, không vội, không chậm,chắc chắn như một nhịp máy được căn chỉnh hoàn hảo.

Bên cạnh, tên trùm côn đồ lao đi như búa tạ bật lò xo ngay từ những bước đầu tiên. Sải chân dài, dồn dập, cơ bắp gồng lên từng đường gân,– hắn muốn dập tắt cuộc đua ngay khi vạch còn chưa kịp nguội.

``Cố kết thúc trong hiệp một đây mà,'' Touka nheo mắt.

``Nhưng đường dài còn nguyên, hụt hơi là khó đấy,'' Shino đáp lại, giọng phẳng như tờ.

Ban đầu, hai bóng song song. Chỉ vài giây sau, khoảnh khắc đầu tiên lộ rõ: tên trùm lấn sang làn Yuki, không đủ để bị phạt rõ ràng, nhưng đủ để ép cô lệch đường, như cái gờ nhỏ trên mặt bàn khiến thước kẻ vỡ nhịp. Một trò tiểu xảo từ đối phương để trấn áp.

Cô nàng mang găng và tất chợt khựng nửa nhịp. Rất nhanh sau đó, cô khép vai, đổi góc chạy, điều chỉnh lại nhịp thở. Mọi thứ trong tích tắc, êm như chuyển số trên xe tự động.

Khán đài thở phào khi cô đã vựot qua trò bẩn mà đối phương tung ra.
``Bị chèn mà vẫn giữ nhịp được kìa\dots'' Sakura thì thầo.

``Không tệ.'' Koishin gật, tay khoanh trước ngực, dấu hiệu cô đã chịu ấn tượng.

Tên trùm ngoái đầu, khóe môi hắn kéo lệch thành nụ cười nhợt nhạt rồi bất ngờ gia tốc lần thứ hai, đạp mạnh xuống đất như muốn nghiền nát đường chạy.

Yuki không bị cuốn theo cú vọt của đối phương; cô chỉ khẽ nghiêng thân, mở thêm một bậc nhịp đều đặn như đồng hồ vừa được lên dây cót. Tốc độ tăng thêm không ào ạt như đối phương, chỉ là một đường thẳng được kéo dài thêm vài centimet mỗi bước, đủ để hai bóng lại song song, đủ để khán đài vừa kịp thở ra.

Đến đoạn cuối của lượt chạy, khi tiếng reo còn chưa tan, trò hề thứ hai ập tới.   Trong khoảnh khắc gậy tiếp sức đổi chủ, hắn giả vờ mất thăng bằng, vung cả cánh tay quá đà, khiến đầu gậy vẽ một vòng cung ngắn hướng thẳng vào sườn Yuki. Một cú ``lỡ tay'' có thể biến thành cú hất ngã nếu cô phản ứng chậm.

Yuki thấy rõ quỹ đạo ấy.

Cô không nhảy sang, cũng không phanh lại. Những gì cô làm đơn giản là hạ hông, trượt thấp, để thanh gỗ vụt qua trên tóc. Cô lao về trước như chưa từng có va chạm, nhịp thở vẫn đều, tốc độ không đổi.

Khán đài chỉ kịp thấy một vệt bóng mờ cắt ngang đường đua, còn Kanade thì tròn mắt:  ``Cậu ấy\dots\ đọc được hết rồi.''

Shino lắc nhẹ đầu, điềm đạm sửa lại: ``Không phải đâu\dots\ là cậu ấy chẳng thèm để tâm.''

Khu vực chuyền gậy trở thành điểm nóng nhất sân. Yae đứng thẳng, hai đầu gối hơi khuỵu, bàn tay đỏ rực giơ về phía trước như tháp pháo chờ lửa. Tiếng gọi vang lên:

``\textbf{Yuki-san!}''

``Đây!''

Chỉ một nhịp thở, cô gái tóc đen-đỏ lao vào vùng chuyền. Gót chân cô chạm vạch trắng, cánh tay vung về sau rồi thả lỏng đúng lúc; chiếc gậy đỏ tươi trượt trong lòng bàn tay Yae, khớp như bánh răng cuối cùng được lắp vào đồng hồ.

Bên kia, tên trùm giao gậy cho gã lực sĩ từng bị Kaigou hất lên không trung môn kéo co. Cả hai đội hoàn tất gần như cùng giây, nhưng tiếng gậy của đối phương va vào kẽ tay, chậm nửa nhịp, tiếng *CỘP* khô khốc vang lên sau tiếng *TÁCH* nhẹ nhàng của Đội Đỏ.

Yuki thu nhỏ bước, tránh vạch, để cơ thể trôi dọc mép ngoài đường chạy. Tên trùm cũng rời làn, cố tình áp sát. Hai bóng song hành trong gang tấc, gió tạt vào tai như lưỡi dao.

Hắn quay đầu, khóe môi kéo lệch: ``Chạy cũng khá đấy.''

Cô nàng không liếc lại, chỉ để tiếng nói mỏng manh bay ngược gió: ``Cậu cũng vậy.'' Câu trả lời mềm như lụa, nhưng đủ lạnh để hơi nóng trên đường chạy bỗng nguội bớt.

Yae bật khỏi vạch như viên đạn, bỏ qua mọi thử nghiệm làm quen với nhịp độ; cô lao thẳng vào đường chạy bằng tốc độ tối đa, mái tóc đỏ tung bay sau lưng thành một dải lửa nhỏ. Đối thủ bên kia - gã vạm vỡ từng bị Kaigou tống lên không trung ở môn kéo co - cũng không chịu thua, mỗi bước chân hắn đạp xuống đều gợn sóng trên mặt sân, tiếng bịch bịch nện đây dồn dập.

Hai bóng song hành, khoảng cách chỉ bằng một gang tay. Shion bật cười, mắt sáng lên như thấy món đồ chơi yêu thích: ``Giờ mới là đường đua thực thụ này!'' Kana nắm chặt tay, thầm lặng bình luận: ``Yae-san không để ai bắt kịp đâu.''

Chưa đầy mười giây, họ bắt đầu vượt nhau từng chút. Gã lực sĩ dựa vào sải chân dài, còn Yae dùng tần suất bước cao. Mỗi lần cô vượt lên, khán giả lại nổ một tràng reo, rồi lập tức nín lặng khi bóng đỏ và bóng trắng trở lại ngang nhau, một vòng lặp khiến ngực người xem nhói theo.

Đúng lúc ấy, đối thủ nghiến răng, giậm thình thịch xuống vũng nước còn sót lại giữa đường. Tia nước bắn ngược, vẽ thành lưới mỏng bay thẳng vào làn của Yae. Không đủ để trọng tài phạt, nhưng đủ khiến mặt đất trơn như phủ lớp mỡ.

Bàn chân Yae chạm phải mặt ướt, khiến cô trượt về phía trước nửa bước. Cô loạng choạng, thân trên nghiêng sang trái, cánh tay giật ra để bắt nhịp. Tiếng ``Á!'' của Aku vừa cất lên thì cô đã xoay hông, biến đà trượt thành cú trượt dài sát mặt đất, y như vận động viên trượt băng. Trong chớp mắt, cô lấy lại trục thẳng, tiếp tục tăng tốc mà không cần chỉnh lại nhịp.

Khán đài vỡ òa với màn xử lý khéo léo của cô. Yae ngoái đầu cười khúc khích, giọng nghe như chọc tức: ``Chơi bẩn ghê nhỉ!'' Gã lực sĩ chỉ hừ mũi, máu dồn lên thái dương, nhưng hắn vẫn phải bật theo nhịp chạy, vì cô gái lực điền đã bứt lên trước mặt hắn nửa thân người, và khoảng cách đang dần nới rộng.

Yae bỗng tháo chốt, bật như cò bỏ đạn. Cô không còn giữ lại chút xăng nào trong bình, mỗi bước đạp xuống như muốn in họa tiết đế giày lên mặt đất. Khoảng cách giữa cô và gã vạm vỡ bắt đầu kéo giãn, từng khoảnh khắc như được kéo căng bằng lực kéo vô hình.

Nhưng ngay trước khi vạch trắng chạm gót, gã lực sĩ bất thình lình lấn sang, nửa bước đủ để cản đường chuyền gậy. Đó là một cái bẫy tinh vi, vừa đủ để qua mắt trọng tài, vừa đủ để Yae mất nhịp nếu cô lóng ngóng.

Cô đã đọc được ý đồ ấy từ ba nhịp trước. Không chậm rãi cũng không vội vàng, cô chỉ nghiêng thân mình một góc vừa đủ như con cò vượt dòng, rồi lao thẳng vào khu vực trao gậy. Giọng cô vang lên trong gió:

``\textbf{Shiina-san!}''

``Đưa đây.''

Cú bắt gậy diễn ra gọn ghẽ mà không gặp phải trục trặc gì. Ở bên kia, gã lực sĩ cũng vừa ép gậy vào tay đồng đội - kẻ đã chơi xấu Kuon nhưng tiếng *CỘP* khô khốc của chúng vang lên chậm hơn nửa nhịp. Hai đường chạy lại trở về vạch xuất phát ngang bằng, như hai dây đàn vừa được chỉnh cùng một nốt.

Yae thả chân, để gió lùa qua tóc, rời đường đua như mũi tên vừa cắm đích. Gã lực sĩ lướt qua, vai nặng trĩu mồ hôi, khịt mũi khoái chí: ``May cho mày là đã tránh được đòn của tao đấy.'' Cô không quay lại, chỉ nhún vai, đáp như thả một chiếc lông xuống nước: ``Không. Là kỹ năng của tớ thôi.'' Dứt câu, cô bước thẳng, để lại sau lưng hắn một khoảng trống mỏng manh nhưng đủ lấp đầy bằng im lặng.

Ở vạch trắng tiếp theo, lượt ba chưa kịp cất tiếng, không khí đã nặng như chì. Chiếc gậy vừa trao tay Shiina, mặt sân như thụt xuống một tấc, tiếng reo bị bóp nghẹt trong cổ họng khán giả. Nếu hai hiệp trước là cuộc rượt đuổi của nhịp tim và kỹ xảo, thì giờ đây, tốc độ nhường chỗ cho đối đầu trực diệncủa hai con người sắt đá nhất. Áp lực không còn nằm trong đồng hồ bấm giây; nó đang đứng thẳng, chắn giữa hai làn, chờ một cái gật đầu để lao vào nhau.

Shiina không vội bứt tốc. Cô bước vào đường chạy như người đi xem mặt đất, từng bước đặt xuống chắc nịch, nặng trĩu, đến mức người ta có cảm giác dù trời có sập cũng không lay nổi thân hình mảnh khẳnh mà chắc nịch ấy. Không rên, không gấp, cô chỉ lẳng lặng bám lấy nhịp chạy của riêng mình, một nhịp chậm hơn đối phương đúng một khoảng lặng của tim.

Bên kia, tên từng chơi xấu Kuon vừa nhận gậy đã quay đầu liếc sang. Đôi mắt hắn gặp đôi mắt cô, và trong khoảnh khắc ấy, bước chân hắn loạng choạng. Chỉ một tích tắc, nhưng đủ để cả khán đài nhận ra thứ cảm xúc hắn cố giấu sau lớp mặt nạ ngạo mạn: e dè, như thể con mồi vừa phát hiện ra thợ săn đang rình trong bóng tối.

``\dots Tch.'' Hắn nghiến răng, quay mặt đi, nhanh tới vội vàng như muốn chạy trốn chính khoảnh khắc yếu lòng ấy.

Shino thở phào, lau mồ hôi trán: ``Nhìn cậu ta kìa, như vừa gặp ma.''

Koishin hừ nhẹ, nhưng khóe môi khẽ nhếch: ``Ai bảo đầu năm cứ làm phiền Kuon-san cơ, rồi bị Shiina-san đã cho hắn một cú đá đít nhớ đời.''

Cô nàng bờm Sư tử vẫn không nói gì, bắt đầu tăng tốc đều đặn, như con sóng thủy triều lặng lẽ dâng lên, mỗi lúc một nuốt gần bờ. Khoảng cách giữa hai người thu hẹp lại từng khoảng, từng hơi thở, từng nhịp tim. Không ai reo hò, không ai dám thở mạnh - cả sân như đang chứng kiến một con thú săn lùi dần con mồi vào chân tường, bằng ánh mắt xanh lạnh tanh và những bước chân không thể cản.

Hắn nghiến chặt hàm, hai má căng cứng như thể chỉ cần buông lỏng một chút, cơn tức sẽ xông thẳng lên não.  Hắn không được phép cho bản thân thua, và nhất là không được phép gục ngã trước mặt Shiina, kẻ từng khiến hắn ôm hận từ đầu năm.

Khi bóng cô gái tóc xám vừa lóe lên bên trái, hắn chủ động đánh lái. Cánh tay phải giật mạnh, hông xoay ngang, toàn bộ trọng lượng dồn sang phần vốn thuộc về cô nàng. Một pha lấn tuyến tinh vi, giống hệt mánh khóe đồng đội hắn từng dùng với Yuki, chỉ khác lần này hắn bỏ thêm lực, thêm hi vọng.

Nhưng Shiina không chút chần chừ.

Cô không lùi, không né, cũng chẳng hạ tốc độ, mà thay vào đó, cô gia tăng sải bước, thân hình mảnh dẻ lao thẳng như mũi tên rời cung. Vai họ chạm nhau, khẽ thốt một tiếng *BỤP* khô khốc giữa tiếng gió.

Trong phần nghìn giây ấy, trọng tài không thấy gì sai, khán đài chỉ kịp thấy hai bóng áo trắng đỏ chồng lên rồi tách ra. Nhưng kẻ lảo đảo lùi về sau lại chính là hắn, gót chân trái trượt khỏi vạch, trọng tâm nghiêng ngả, hai tay giật vội để bắt lấy không khí.

``Gì\dots!?'' tiếng thốt lên của hắn bị gió nuốt chửng.

Shiina đã vượt qua, mái tóc xám phấp phới như cờ chiến thắng. Cô không quay đầu, chỉ để lại một câu thì thầm đủ lạnh đóng băng đường chạy: ``Trò cũ rích cho trẻ con.''

Hắn nghiến răng, gân má nổi lên như sợi dây thừng căng cứng. Gã vẫn chưa thể bỏ cuộc, lại tiếp tục lao tới, mắt lóe lên tia quyết tử.

Giữa đường chạy, hắn bất ngờ giả vờ trẹo cổ chân, nghiêng người sang phải, tạo ra khoảng trống chật hẹp như cái bẫy. Một bước sai lầm của Shiina thôi cũng đủ khiến nhịp cô loạn, đà bứt phá tan tành.

Nhưng cô gái tóc xám chỉ khẽ cúi đầu, mũi giày đẩy mạnh xuống đất.  Cơ thể cô bậc lên, vút qua đầu hắn như làn khói trắng lướt qua song sắt, như thể cô đã đọc kịch bản từ trước khi hắn nghĩ ra nó. Cứ thế, khoảng cách giữa hai người được nới rộng ra bằng những bước chạy đầy dứt khoát, không hề đả động gì tới đối thủ.

Cô chỉ tập trung vào vạch chuyền gậy phía trước, nơi người cuối cùng của nhóm đang đứng chờ để hoàn thành cuộc đua.

Touka trố mắt, giọng chỉ còn là hơi thở: ``Phản xạ ấy\dots\ quả thật là quá phi thường.''

Shion mỉm cười, khẽ đáp như ru: ``Sư tử có sống như thế giới trần gian chúng ta đâu mà\~{}''

Đối thủ cũng lao tới vạch, gương mặt căng cứng. Hắn nghiến răng, quay sang cậu bé cuối cùng - người vốn lọt vào đội hình chỉ vì thiếu người - gằn giọng: ``Chạy cho ra hồn! Dẫm lên đường mà ngã thì đừng trách tao!''

Cậu học sinh kia giật nảy, môi run: ``V-Vâng\dots! Mình sẽ cố hết sức!'' Nói đoạn, cậu cầm lấy chiếc gậy rồi vút đi trong tầm mắt.

Ở phía bên kia, Shiina không nói thêm, chỉ gật đầu nhẹ rồi trao gậy. ``Kaigou.''

``Ừ.''

Cả hai hoàn tất động tác trong tích tắc, gọn đến mức không kịp nghe tiếng gió.

Kaigou chỉ mới nhấc gót, khán giả còn chưa kịp thở thì cô đã dứt điểm ba bước, sau đó, như có lò xo nổ dưới lòng bàn chân, cô đẩy toàn thân lao vút thành một vệt mờ đỏ. Không phải ảo ảnh, chỉ là tốc độ vừa vượt ngưỡng có thể nhìn bằng mắt thường. Cậu học sinh bên Đội Trắng cũng bật ra, nhưng đất dưới chân hắn như lún xuống, mỗi nhịp đều chậm nửa khung hình.

Kanade ú ớ: ``Khoảng cách giữa hai người họ\dots\ lớn quá.''

Sakura lắc nhẹ đầu: ``Không chỉ nhanh, cô ấy mang theo cả tính khí chèn ép, đối thủ bị đè ngay từ bước đầu.''

Kaigou lao tới như mũi tên rời cung, chẳng mấy chốc mà đã bắt kịp đối thủ. Cô không hề để ý đến sự tồn tại của cậu học sinh bên Đội Trắng, hoàn toàn chỉ để ý tới việc băng băng về đích. Nhưng chính sự lạnh lùng ấy lại khiến cậu ta hoảng sợ hơn bất kỳ lời đe dọa nào.

Trong khoảnh khắc hai người song hành, cậu học sinh kia liếc sang. Ánh mắt cậu ta chạm phải đôi mắt xanh của Kaigou. Nó là ánh mắt không giận dữ, không căng thẳng, chỉ là một vực thẳm băng giá vô cảm. Cậu ta nuốt khan, nhịp chân loạn xạ, như thể có bàn tay vô hình đang siết chặt cổ họng.

``Đừng\dots\ đừng lại gần quá chứ\dots!'' tiếng lẩm bẩm yếu ớt thoát ra từ miệng cậu, nhưng Kaigou đã vượt qua. Không một lời nói, không một cái liếc nhìn cuối cùng. Cô chỉ đơn giản bỏ lại cậu phía sau, một khoảng trống lạnh lẽo và im lặng.

Shino nhíu mày, lòng thấy thương hại: ``Tội nghiệp\dots\ cậu ta như bị bóng ma vượt qua vậy.''

Aku rùng mình, mắt tròn xoe: ``Không phải bóng ma\dots\ là băng tuyết đấy. Cô ấy\dots\ lạnh đến đáng sợ.''

Suzuki nhìn người chị của mình, không biết nên cảm thấy ngưỡng mộ hay nên thương cảm nữa, dù rằng cô là người hiểu rõ tính cách của Kaigou hơn bất cứ ai.

Khoảng cách giữa hai người cứ thế bị kéo giãn, từng mét băng cỏ như được gấp lại dưới chân Kaigou. Cô đang bứt phá bằng cả ý chí, mỗi bước dài thêm, nhanh thêm, biến kẻ đuổi theo thành một chấm trắng mờ dần trong tầm mắt. Không còn khả năng rút ngắn, không còn đường lùi, chỉ còn cách để lại hậu phương cho đối thủ nuốt trọn bụi đất.

Vạch đích hiện ra như một vết cắt trắng giữa biển cỏ xanh. Tiếng reo của khán đài dồn lên thành một làn sóng âm, đập vào màng nhĩ, rung cả ngực. ``Kaigou!'' tiếng gọi vang lên không chỉ từ một chỗ, mà từ khắp các khán đài, hòa thành tên riêng duy nhất của chiến thắng.

Cô lao qua vạch như mũi tên đã cắt dây cung, không kịp nghĩ, không kịp ngoái. Không có màn ăn mừng vung tay, không có cú giậm chân khoe dằn, chỉ một luồng gió đỏ xẹt qua, để lại sau lưng tất cả những gì từng nghi ngờ, từng dè chừng. Đó là cách một kết thúc nên im lặng: gọn gàng, dứt khoát, không dư thừa.

Chừng năm, bảy giây sau, bóng trắng mới chật vật hiện ra. Cậu học sinh Đội Trắng gần như gục ngay trên vạch, hai tay chống đầu gối, mồ hôi rơi thành dòng xuống cỏ. Không còn sức để nuốt trọn nỗi thua, chỉ còn tiếng thở dốc như xin được phép gục ngã, và thế giới bên kia vạch đã thuộc về người khác từ lâu.

*\textbf{TOÉÉÉÉÉT!!}*

Tiếng còi dài vút qua sân cỏ, như mũi tên xuyên qua tấm màn căng thẳng cuối cùng.  Ayame và Wata chỉ cần trao nhau ánh mắt đã hiểu: kết quả phân thắng bại đã rõ mười mươi. Họ cùng gật đầu nhẹ, rồi họ cất giọng song thanh, vang như pháo hoa cuối mùa: ``\textbf{Đội Đỏ chiến thắng!}''

Khán đài nổ tung. Tiếng reo không chỉ là âm thanh, mà là một thứ chất lỏng nóng hổi đổ xuống từng bậc xi-măng, cuốn theo bụi cỏ, mồ hôi, cả những giọt nước mắt không kịp rơi.

Các thành viên Đội Đỏ như sóng thủy triều đỏ, ào ạt vỡ vào bốn bóng hình nằm giữa trung tâm đường chạy.

Kaigou bị nhấc bổng lên, không phải bởi đôi tay, mà bởi chính tiếng hò reo, cô bất giác cười, nụ cười hiếm hoi như bông tuyết giữa hè khi đã một lần nữa đánh bại đám cặn bã của trường.

Yuki bị kéo vào một cái ôm siết chặt đến mức cô phải vỗ vai mọi người trấn an: ``Tớ vẫn cần xương sườn để thở.''

Yae thì bật khóc, nước mắt hòa mồ hôi thành những vệt đỏ loang trên má, còn Shiina đứng ngoài vòng, tay khoanh, nhưng khóe môi khẽ nhếch, đủ để ai tinh mắt thấy cô đang rung lên theo nhịp trống ngực.

Phía xa, ba tên côn đồ như ba cái bóng bị quên bật đèn. Họ không dọn đi, cũng chẳng dám lại gần, chỉ đứng đó, mặt tối sầm như bức tường vừa bị khắc một dòng chữ xấu hổ bằng mực đỏ tươi.

Ayame vỗ tay ra hiệu, giọng vẫn còn vang vọng trong hầm trời khán đài: ``Được rồi~! Tập trung lại nào! Giải lao năm phút rồi tổng kết hội thao!''

Wata giơ cao chiếc loa cầm tay, như thể muốn gom hết cả không khí vào trong lời mình: ``Tất cả các lớp di chuyển về khu tập kết! Đường thẳng một hàng, không chen lấn!''

Dòng người bắt đầu chảy về trung tâm sân, tiếng bước chân rầm rập như mưa đầu mùa. Cờ đội phấp phới, loa phát thanh rè rè, mùi cỏ nóng hôi hổi bốc lên từng đợt. Âm thanh chưa hạ nhiệt, mà nó chỉ đổi chiều, từ reo hò sang xì xào, từ xì xào sang tiếng cười khúc khích, rồi lại reo, như con sóng không bao giờ muốn vỗ vào bờ.

\il

Nhưng ở nơi không ai để ý,
trên bậc thang khán đài cao nhất,
một hồn ma bé gái đang ngồi.

Cô bé mắt xanh như viên bi thủy tinh nằm trong ly nước đá, tóc tết bím dài chấm gót chân, lặng lẽ buông hai chân đung đưa trong khoảng không.
Không ai thấy cô, hoặc có lẽ ai cũng thấy, nhưng lại tưởng là ánh sáng lạ.

Cô nhìn Kaigou.
Kaigou đang thoát khỏi vòng vây, lững thững đi ven đường chạy, tay xoa nhẹ chỗ cổ áo bị kéo giãn.
Hồn ma bé gái nghiêng đầu, mái tóc bím lật sang một bên như dải ruy-băng bị gió thổi, và mỉm cười, nụ cười không dành cho người sống.

Cô nhìn Kuon.
Kuon đứng dưới gốc cây bàng, máy quay cầm tay vẫn hướng về phía đám đông, nhưng ống kính lại lệch một chút, để lọt vào khung hình bóng hồn ma bé gái.
Kuon không biết (hoặc giả vờ không biết) chỉ chau mày, chỉnh lại tiêu cự, như thể vừa bắt được một bóng sáng lạ thường.

Cô nhìn Suzuki.
Suzuki ngồi trên khán đài, hai tay nắm chặt lan can, mắt dán vào chị mình.
Hồn ma bé gái đưa ngón tay trỏ lên môi, ra hiệu ``suỵt'', rồi chỉ vào ngực Suzuki, nơi trái tim đang đập loạn nhịp vì niềm tự hào lẫn nỗi lo âu không tên.

Không ai nghe cô thì thầm.
Chỉ có gió, vừa khẽ đưa ba lọn tóc bím bay lên, như muốn viết một câu chào không thành lời:
``Em vẫn ở đây.
Và em thích cách các anh chị chiến thắng lắm.''

\il

\ct

Nắng xế tắt dần, phủ lên sân Aogami lớp vàng úa. Bóng học sinh kéo dài, chồng chéo như mạng nhện mỏng. Cả ngàu thi đấu đã qua, nhưng dòng người vẫn rạo rực, chỉ khác là từ reo hò chuyển thành nín thở chờ kết quả.

Hai khối Đỏ và Trắng đứng xếp hàng chỉnh tề, như hai dải ruy băng khổng lồ căng giữa sân. Trên bục gỗ cao, Ayame cầm mic, vai nhún nhảy, còn Wata đứng cạnh, ôm tấm bảng thiệt lớn như tấm khiên.

Nàng Hội trưởng khẽ xoay mic, cất giọng ngọt vang: ``Cảm ơn cả nhà vì một ngày cháy hết mình nha!''

Tiếng vỗ tay nổ ra như mưa rào, dội xuống từng bậc khán đài, vỗ vào mái tóc, vai áo, vỗ cả vào khoảnh khắc mà mọi người vẫn còn nghẹn lại trong cổ họng, như thể vừa được phép thở sau cả ngày thi đấu căng như dây đàn.

Ayame xoay nhẹ míc-rô, mắt liếc qua tấm bảng tổng sắp phía sau, giọng vẫn ngọt như kẹo bông nhưng đã thêm một nét sắc như lưỡi dao giấy: ``Nào, trước khi chúng ta cùng hò reo lần cuối, chúng ta sẽ cùng nhau tổng kết lại kết quả của các môn thi nha\~{}!''

Tiếng vỗ tay lại nổ ra, lần này không như mưa rào, mà như mưa đầu mùa, rơi lâu hơn, ấm hơn, và không ai muốn che ô.

Nàng Phó gật đầu, trầm ấm: ``Trên tay mình đây là kết quả chung cuộc của các môn thi. Bây giờ, mời mọi người cùng lắng nghe.''

Wata trải bảng kết quả, giọng trầm ấm đọc từng hạng mục như gõ vào trống lớp.
``Đầu tiên, tamaire, Đội Đỏ chiến thắng.'' Ayame nghiêng mic, mắt cười: ``Khởi đầu nóng, bể trời đỏ luôn!''

``Tiếp theo là môn kéo co, Đội Đỏ chiến thắng.'' Ayame vỗ tay, giọng như thả kẹo: ``Dây thừng về phe nữ sinh rồi, các anh cứ nằm dưới đất mà ngắm mây!''

``Tiếp tới là bóng né, Đội Đỏ chiến thắng.'' Ayame nháy mắt: ``Trái bóng mạnh thế còn né được, lại còn hạ gục được tên trùm!''

``Trong môn đấu vật, Đội Đỏ thắng nội dung nữ và hỗn hợp. Đội Trắng thắng nội dung nam. Đội Đỏ giành chiến thắng chung cuộc.'' Ayame thổi phù một lọn tóc: ``Đai đen nữ sinh, ôm ngã luôn cả trái tim khán giả.''

``Ở nội dung bóng rổ, Đội Đỏ chiến thắng.''
Ayame giả bộ ném ảo: ``Ba điểm cuối cùng như viên kẹo đắng, bỏng vào miệng Đội Trắng!''

``Tại nội dung chạy vượt chướng ngại vật, Đội Đỏ chiến thắng.''
Ayame nắm tay đưa cao: ``Tường cao, hố sâu, nhưng cao hơn vẫn là tinh thần đỏ!''

``Và cuối cùng, môn thi chạy tiếp sức, Đội Đỏ chiến thắng.'' Ayame hạ giọng, như khép lại bản hùng ca: ``Gậy đỏ trao tay, niềm tin trao nhau, và họ chạy thẳng vào lịch sử một ngày hội thao!'

Mỗi lần ``Đội Đỏ'' vang lên, khán đài phải bên như sóng thủy triều đỏ dâng cao, reo đến nỗi lấn át cả tiếng ve đầu hè, để lại sau lưng một vệt sáng loang dần trên mặt cỏ như mực đỏ khô giữa trang giấy cuối cùng của một ngày hội thao không thể quên.

Kaigou khoanh tay trước ngực, mắt híp lại, khẽ phát ra một tiếng hừ như thể vừa chứng kiến điều đã nằm lòng từ đầu:  ``Cũng may, không phí công cả tháng luyện đến nát giày.''

Yae đưa tay chống hông, nở nụ cười rạng rỡ đến mức nắng cuối chiều cũng phải chùng xuống: ``Nhưng tới mức áp đảo đội bên họ thì\dots\ tớ cũng hơi ngạc nhiên đấy.''

Koishin đứng nghiêng nửa vai, má quay đi chỗ khác, như thể tiếng reo hò còn vang vọng làm cậu khó chịu:  ``\dots Chậc. Đã đoán trước từ lúc thấy đọi bên kia run chân với vẻ mặt ỉu xìu rồi.''

Kuon ở phía sau, ống kính vẫn hướng về đám đông, chỉ thả ra một hơi thở nhẹ tênh, đủ để làn sương mỏng bám vào miếng bọc nilon: ``\textit{Không cẩn thận đội bạn còn ra về không có gì luôn. Đội lớp mình mạnh thật.}''

Ayame vỗ tay như muốn gom cả sân vào lòng bàn tay, đôi mắt cười híp lại thành hai vầng trăng thượng tuần.

``Nhưng mà\~{}!'' cô kéo dài giọng, như thể đang nhấc một tấm màn nhung lên trước mặt khán giả, ``hạng mục cực kỳ quan trọng chưa công bố đấy!''

Tiếng xì xào lan ra từng hàng ghế, từng bậc cỏ, như gió thổi qua cánh đồng lau trước cơn bão.

Ayame quay lưng, đôi mắt lướt qua hai khối màu đỏ-trắng đang căng như dây đàn chuẩn bị đứt. Cô nhếch míc-rô, để gió cuốn giọng mình đi xa: ``Giờ này, hai đội đã nếm đủ mùi thắng-thua\dots\ nhưng còn món tráng miệng tinh thần chưa dọn lên bàn.''

Wata bước sát, vai hơi cúi như chỉ huy trận chiến cuối: ``Mời đại diện Đội Trắng và Đội Đỏ lên sân khấu!''

Khán đài khẽ xao động.

Ayame cất giọng rõ ràng:
``Đội Trắng các bạn như bản thiết kế hoàn hảo: đồng đều, kỹ thuật gần như không tì vết, đội hình đẹp đến từng milimet.''

Nàng Phó gật nhẹ, bổ sung:  ``Mình đếm được đúng một lần giật nhịp, còn lại là mọi thứ khá là bình yên và không có sai sót gì quá nghiêm trọng.''

Mấy bạn của Đội Trắng ở hàng đầu khẽ nở mũi, vai thả lỏng.

Ayame giơ ngón tay, cười nhẹ: ``Nhưng mà, các bạn hơi\dots\ an toàn. Thiếu chút lửa bùng lên khiến người ta phải bật dậy khỏi ghế.''

Wata nhìn sang khối Trắng, trầm giọng: ``Cảm giác như xem bản giao hưởng chỉ toàn nốt trắng. Nó đẹp, nhưng lại không có nốt cao nào để khiến người xem chú ý cả.''

Bên dưới, một bạn nữ Trắng cắn nhẹ môi, tay siết cờ. Người bên cạnh cô thì thầm: ``Mình\dots\ có nên xin thêm bản nhạc nổ không nhỉ?''

Ayame khép míc-rô, mắt nửa híp: ``Đừng buồn, đó là lời khen có gai để lần sau các bạn có thể rút kinh nghiệm cho kỳ đại hội sau. Chúc mừng các bạn vì đã cho khán giả chiêm ngưỡng một màn trình diễn tuyệt vời.''

Hội trưởng xoay người, ánh mắt lướt qua bãi đất còn vương vết giày đỏ: ``Đội Đỏ các bạn\dots\ nói thật là\dots\ lỗi thì có đấy nha\~{}!''

Tiếng cười rộ lên như sóng nhỏ. Suzuki nóng bừng tai, nhớ lại cú nhảy có phần bạo dạn và liều lĩnh của mình.

Kaigou cúi xuống, khẽ huých:  ``\dots Còn nhớ lúc em tự ý nhảy từ trên cao xuống chứ? Làm chị sợ chết khiếp đấy!''

``\dots Em xin lỗi\dots\ Tại em nghĩ như thế sẽ khiến mọi người chú ý hơn\dots'' cô thỏ thẻ, giấu mặt sau lá cờ đỏ đang phất phơ.

Wata gõ nhẹ míc-rô, xen vào: ``Nhưng, đúng như Suzuki-sann nói đấy! Đó lại thành điểm nhấn khiến khán giả nhớ mãi.''

Ayame gật, giọng bỗng ấm: ``Các bạn có thứ mà Đội Trắng thiếu, và cũng là cái cực kỳ quan trọng: đó là sức kéo khán giả.'' Cô vung tay về phía khán đài, nơi vẫn còn những dải cờ đỏ chưa kịp gỡ.

Nàng tóc vàng tiếp lời, mắt nheo lại: ``Mình đếm được không dưới bốn lần khán đài tự nhún theo, chưa kể chỗ Suzuki-san hạ cánh, cả khối ba hét oà lên.''

Kuon hơi ngẩn ra, rõ ràng ngạc nhiên với những lời có cánh của hai người. Shion đứng sau cười khúc khích, vỗ vai Kuon: ``Ô, ai đó được khen kìa\~{}'' Aku khoanh tay, nhún vai:  ``Không ngờ luôn đấy\dots''

Phía Đội Trắng, một bạn nữ khẽ siết cờ, thì thầm với bạn bên cạnh: ``\dots Họ có ma lực thật.'' Người đội trưởng Trắng thở dài, nhưng khóe miệng hơi nhếch, như ngầm công nhận.

Ayame quay lưng, giọng bỗng sắc như dao: ``Và vì tiêu chí quan trọng nhất của ouen gassen là khán giả, vậy nên\dots\''

Cô vung tay lên, khoét một khoảng trời đỏ rực.
``Đội chiến thắng trong cuộc thi cổ vũ xin được gọi tên\dots''

Một giây im lặng, đủ để mọi nhịp tim lạc nhịp. Ngay lúc này, cả hai người trên sân khấu cùng nhau đồng thanh: ``\dots \textbf{ĐỘI ĐỎ!}''

Sân trường vỡ òa như chai sâm panh khổng lồ. Tiếng hò reo dội vào nhau, tạo thành bản nhạc không cần dàn dựng.

Suzuki bị Kaoru nhấc bổng giữa không trung, cả hai xoay một vòng như con quay.
``Chúng ta làm được rồi! Đội chúng ta thắng rồi!!''

Yae đập nắm đấm vào bầu không khí, mắt sáng quắc: ``Không chỉ chuẩn, mà còn đét!''

Phía bên kia, hàng Trắng buông cờ, thở ra một hơi dài. Họ gật nhẹ, như thể nói: ``Được thôi, năm sau ta sẽ làm tốt hơn.''

Ayame nhe răng, cười toe toét: ``Năm nay, Đội Trắng không chỉ mặc áo trắng\dots'' cô nghiêng đầu, đôi mắt nháy lấp lánh trò đùa, ``\dots mà còn về đích trắng tay luôn, nha\~{}!''

Tiếng cười bùng lên khắp khán đài, lan nhanh như sóng lớn vừa đổ bờ. Nó cuốn theo ba gương mặt đứng cuối hàng - ba kẻ côn đồ ở Đội Trắng - khiến chúng càng cúi sát, hai má như phủ lớp than đặc quánh đêm giao thừa. Mỗi tiếng bật cười là một nhát búa đóng thêm vào mũi tên hổ thẹn đang găm sau lưng chúng.

Bên kia chiến tuyến, khối Đỏ quây quần thành vòng tròn kín, vai kề vai, cánh tay quấn chặt lấy nhau. Áo đẫm mồ hôi nhưng đôi mắt vẫn rực lên thứ lửa vừa thắp từ khán đài, thứ lửa mà họ đã đốt bằng cả trái tim lẫn nhịp chân lệch tông. Họ không cần khẩu hiệu nữa, chỉ cần cảm nhận hơi thở của nhau là đủ để biết mình đang cùng nhau chiến thắng.

Khối Trắng đứng thành hàng thẳng tắp, môi mím đến trắng bạc. Lá cờ trắng trong tay họ rung nhè nhẹ, không phải vì gió mà vì những ngón tay đang cố giấy đi thứ run rẩy khó nói. Họ nhìn thấy Đội Đỏ được nâng bổng giữa vòng vây cổ vũ, và trong khoảnh khắc ấy, họ hiểu ra: thứ vũ khí đáng sợ nhất không phải là sự hoàn hảo, mà là khi khán giả quên mất bạn đang thi đấu để cùng nhảy theo nhịp của đối thủ.

\ct

Ayame vỗ tay hai cái, mắt híp lại như mèo vừa bắt được chuột: ``Nào, giờ đến món quà bất ngờ cuối ngày: những giải thưởng phụ!''

Wata bước lên, giọng trầm ấm đọc từng hạng mục như gõ trống trận.

``Cầu thủ bóng né xuất sắc nhất: Akaba Koishin-san.''

Ayame nhún vai, cười tinh quái: ``Nghe đồn cậu ấy né trái né phải còn né được cả tên trời đánh nữa!''

Koishin quay mặt đi, tai đỏ ửng: ``\dots Không cần phải công bố đâu.''

``Cầu thủ bóng rổ nổi bật: Shiromi Yae-san!''

Ayame giơ hai tay làm điệu ném ảo, miệng kêu ``vèo'': ``Ba điểm cuối cùng như viên kẹo đắng, bỏng luôn cả lưỡi Đội Trắng!''

Yae cười toe toét, mắt lấp lánh.

``Thành viên năng động nhất: Hikawa Suzuki-san!''

Ayame nháy mắt về phía khán đài: ``Cô bé này nhảy cao hơn cả điểm số toán của mình ngày xưa!''

Suzuki lúng túng, tay tự chỉ vào mình như không thể tin đựo: ``M-Mình á\dots?''

``Thành viên fair-play nhất: Kaizuka Yuki-san!''

Ayame đưa ngón tay lên môi, ra hiệu ``suỵt'' dễ thương: ``Không chỉ thân thiện, cô ấy còn là người duy nhất không chửi thề trong cả ngày thi!''

Yuki gãi má, cười nhẹ: ``Chỉ là tớ làm tròn trách nhiệm thôi mà.''

``Thành viên ấn tượng nhất: Fukami Shiina-san!''

Ayame hạ giọng, như thể đang kể chuyện ma: ``Knock-out đối thủ bằng cái nhìn và tránh né đòn bẩn bằng cách\dots\ không thèm nhìn lại!''

Shiina khoanh tay, mặt không biểu cảm: ``Bình thường.''

Kuon đứng phía sau, khẽ gật, ống kính lia qua từng gương mặt: ``Cậu cũng xứng đáng đấy chứ.''

Ayame giơ cao míc-rô, kết luận như bật nút: ``Một lần nữa, hội thao năm nay đã chứng minh rằng, đẹp trai không bằng chai mặt, mà chai mặt thì\dots\ vẫn thua con gái Đội Đỏ!''

Những tiếng cười vang lên như sóng, từng hàng ghế đỏ truyền tay nhau cười, từng hàng ghế trắng cũng phải bật cười khe khẽ, còn Hội trưởng thì nghiêng đầu, mắt híp lên như mèo vừa nhấc khúc đuôi dài khỏi bát sữa.

``Nhưng vẫn chưa hết!'' Ayame xoay míc-rô, giọng bỗng trầm xuống như múa roi trên mặt hồ. Một luồng gió lạ chạy dọc khán đài; lá cờ cuối cùng phất phơ rồi buông thõng.

``Và\~{}! Cuối cùng, trong hội thao năm nay, chúng ta có thêm một giải thưởng nữa!'' Cô mỉm cười, nụ cười không vừa: khóe môi nhếch cao như mũi kéo căng giấy, mắt lấp lánh ánh đèn sân khấu.

``\textbf{Giải thưởng vinh dự dành cho những người chơi không fair-play nhất!}''

Cả sân trường Aogami như bị bóp nghẹt. Tiếng ve ngừng kêu, tiếng giày nghiến cỏ cũng tắt. Ba tên côn đồ giật mình, vai thót lại, mắt lén liếc nhau tìm lối thoát.

``Không ai khác, đó chính là\dots'' ngón tay Ayame vẽ một vòng cung rồi dừng thẳng vào họ: ``\textbf{Ba bạn bên Đội Trắng kia\~{}!}''

Tiếng ``Cái gì!?'' vang lên như vỡ kính, nhưng lập tức bị khán đài nuốt chửng. Họ không la ó, chỉ dồn hết hơi thở vào im lặng, chờ xem lá bài úp cuối cùng.

Hội trưởng hội học sinh không đáp. Cô chỉ búng tay một cái một cách gọn lẹ, như cú nhấn đóng cảnh của đạo diễn.

Màn hình LED khổng lồ từ sau cô bật sáng. Ánh sáng trắng xanh lạnh tưới xuống mặt cỏ, biến sân thành rạp ngoài trời. Đoạn phim chạy chậm phat lại cú huých khuỷu trong tamaire khiến Kuon suýt ngã, cú giậm giày tạo vũng trơn trên đường chạy, cú lấn làn trong phần thi chạy tiếp sức vừa rồi, những trò bẩn tưởng chừng rất kín kẽ và không ai thấy, nhưng từng khung hình lại rõ đến từng gân cổ.

Không lời bình luận, không nhạc nền, chỉ tiếng tim ai đó đập loạn trong khán đài. Ba tên kia cứng họng; mồ hôi lấp lánh trên thái dương dưới ánh đèn màn hình.

Uzuki đứng bên máy quay, khẽ cười như mèo vừa bắt chuột: ``Máy quay không biết nói dối đâu nha.''

Kuon thở dài, đủ để làn sương mỏng bám kính ống: ``Tôi đã bảo rồi mà.''

Kaigou chéo tay, giọng lạnh như băng vụn rơi: ``Vải thưa không che được mắt Thánh. Quả thật linh nghiệm thật\dots''

Ba cái đầu cúi xuống vì ánh đèn quá chói, vì tiếng cười khúc khích đang lan từng hàng ghế như lửa cháy rơm. Dĩ nhiên bọn chúng chẳng hề tỏ vẻ ăn năn gì, nhưng chẳng thể nào chối cãi được.

Ayame vỗ *BỐP* một tiếng, đủ làm cả sân giật mình: ``Yên tâm! Không có hình phạt gì đâu\~{}!'' Cô nở nụ cười tươi như vừa tung kẹo cho trẻ con, tiếp lời: ``Chỉ là ghi nhận thôi!''

Khán đài bật ra những tiếng cười nhẹ, rồi to dần, cuối cùng trở thành tràng pháo tay rộn rã. Ba tên kia cắn môi, hai bàn tay siết thành nắm đấm trắng bệch, nhưng không thể làm gì hơn ngoài cúi đầu giữa trận cười vang trời của cả trường.

\ct

Loa phát thanh rè rè buông tiếng nhạc hòa tấu nhẹ, như thứ nước xả vải phủ lên cả ngày căng thẳng. Các khối lớp bắt đầu di tản, từng dòng người đi chậm hơn, bước ngắn hơn, như kẻ vừa tháo dây giày sau cuộc chạy marathon. Trên sân chỉ còn lại những mảnh giấy màu đỏ-trắng bị gió thổi vo ve, vài chai nước ngửa cổ, vài chiếc băng rôn tuột khỏi dây, phơi phới trên khán đài trống.

Ánh đèn đường bật sáng từng bóng cam lờ mờ, in hình lá phong đầu mùa rơi lả tả xuống đường chạy, lát những bậc thang bằng bóng loang lổ. Mùi cỏ nóng hôi hổi dần nhường chỗ cho mùi đất se lạnh; tiếng cười râm ran lùi xa, thay bằng tiếng dép lê lê lết trên xi-măng, lộc cộc, đều đặn như kim đồng hồ đếm ngược.

Trên bục sân khấu tạm, Ayame gập mic lại, vỗ nhẹ như ru nó ngủ; Wata kéo khóa túi đựng bảng tổng sắp, cái kéo *KẸT* một tiếng, như chấm dứt một kỳ quan. Cột cờ cao nhất giật mình rung, lá cờ đỏ thẫm buông thõng, phần phật một cái rồi im bặt, như trái tim khổng lồ của cả trường vừa thôi đập.

Wata bước lên bục, giọng trầm ấm như tiếng chuông trường tan học: ``Thưa toàn thể các bạn học sinh trờng Trung học Aogami, hội thao lần thứ 82 của trường chính thức khép lại.''

Ayame giơ cao míc-rô, nắng cuối chiều lướt qua mái tóc cô như dải lụa đỏ: ``Cảm ơn tất cả mọi người đã cùng nhau đốt cháy hết mình trong ngày hôm nay!''

Tiếng vỗ tay vang lên, không dày như lúc chiến thắng, mà rời rạc, như mưa đầu thu.

Mặt trời đã lặn phân nửa, chỉ còn một mảnh tròn cam cháy treo trên mái nhà xa tít. Ánh sáng cuối ngày đổ xuống sân cỏ, nhuộm đỏ rực, trùng màu áo đội chiến thắng, như thể thiên nhiên cũng gật đầu công nhận.

Gió thổi nhẹ, hơi lạnh len lỏi vào từng kẽ áo, cuốn theo mùi cỏ khô, mùi bụi sân, và cả mùi mồ hôi còn vương trên khán đài. Những tiếng cười râm ran tắt dần, thay vào đó là những cái thở dài mỏi mệt, nhưng nhẹ nhõm.

Ở đâu đó trong đám đông, có người chỉ đơn giản ôm lấy chiến thắng, cười tươi như trẻ con được kẹo. Có người lặng lẽ nhìn về phía đội bại, môi mím chặt, trong lòng bắt đầu nhen nhóm một điều gì đó không tên: có thể là nuối tiếc, có thể là quyết tâm, hoặc chỉ là một câu hỏi nhỏ: ``Năm sau, liệu chúng ta có làm được không?''

Nhưng dù thế nào, tất cả đều biết: ngày hôm nay đã khép lại, và ngày mai sẽ lại mở ra, với hay không với chiếc cúp trên tay.

\ct

Sau buổi tổng kết, sân trường như bãi biển thủy triều rút: học sinh các lớp dần tản ra trở về sau một ngày, chỉ còn lớp 1-C bám rễ dưới tán cây bàng. Họ tụ thành một khối vuông vức, không vì lệnh giáo viên, mà vì tiếng gọi nhỏ nhắn nhưng đầy quyền lực của Yuko.

Cô chủ nhiệm đứng trên bậc xi-măng cao hơn nửa bậc, tay không cầm míc-rô vẫn đủ khiến mọi ánh mắt hướng về. Gương mặt cô không cười, chỉ hơi nhếch đuôi mày như người thầy đang chờ lớp im lặng. Khi tiếng xì xào tắt hẳn, Yuko mới mở miệng, giọng vừa đủ nghe, không vang, không nhỏ: ``Tôi biết, hôm nay các em đã chạy, đã đổ mồ hôi, đã hét đến khản cổ. Nhưng trước khi về, tôi muốn nói ba điều.''

Cô giơ ngón tay thứ nhất, thẳng tắp như cây thước kẻ.

``Thứ nhất: chiến thắng này không tự nhiên mà có. Nó là kết quả của việc tất cả mọi người trong lớp mình đã dày công tập luyện cực kỳ chăm chỉ và nghiêm túc.''

Ngón tay thứ hai vươn lên, không nhanh, không chậy.

``Thứ hai: điểm cao nhất không nằm ở bảng tổng sắp, mà nằm ở chỗ hôm nay không ai bỏ cuộc. Dù đã bị đối thủ chơi xấu, các em đều đã làm đúng nhiệm vụ của mình.''

Ngón tay cuối cùng giơ lên, Yuko khẽ siết lại như đang cầm một hạt cát.

``Thứ ba: từ ngày mai, chúng ta lại về guồng học tập. Tôi không muốn nghe ai đó khoe thành tích rồi nghỉ học sớm. Hãy nhớ, lớp 1-C thắng hôm nay, nhưng cũng có thể thua ngày mai nếu chúng ta không giữ lửa.''

Cô dừng một nhịp, mắt lướt qua từng khuôn mặt, từ Kaigou đang gãi cổ đến Suzuki đỏ ửng vì xúc động. Rồi cô khép ba ngón tay lại thành nắm đấm nhỏ, đặt lên ngực.

``Vì vậy, trước khi tan lớp, tôi muốn cả lớp mình chụp một bức ảnh. Không phải để khoe, mà để nhắc nhau: đây là lần đầu tiên, và cũng phải là tiêu chuẩn cho mọi lần sau.''

Yuko bước xuống khỏi bậc xi-măng, hất nhẹ cằm về phía khung cửa sổ có ánh hoàng hôn đang đổ. Giọng cô vẫn cứng, nhưng đuôi câu đã mềm:

``Giờ thì xếp hàng. Ai cao đứng sau, ai thấp đứng trước. Không chen, không cười hở lợi. Tôi muốn bức ảnh này đẹp theo cách của 1-C: ngăn nắp, im lặng, và đầy kiêng hãnh.''

Cả lớp không reo hò, chỉ khẽ di chuyển, tiếng giày kêu lộp cộp nhịp nhàng như một đội quân nhỏ vừa được rèn luyện. Trong khoảnh khắc trước khi máy ảnh bấm, Yuko đứng ngoài khung hình, hai tay khoanh, mắt nhìn thẳng vào ống kính. Cô không cười, nhưng đôi môi mím lại ẩn một niềm vui được cất giấu kỹ càng.

Ngay lập tức, cả lớp nhốn nháo.

``Ơ khoan đã, tớ đứng chỗ này cơ!''

``Đẩy sang tí đi, che mất tớ rồi!''

``Đứng thấp xuống đi, cao quá rồi đấy!''

Akane kéo tay Yuka, cố chen lên phía trước. ``Yuka! Lại đây, lại đây! Chỗ này đẹp này!''

Cô bạn thân dụi mắt, giọng lười nhác: ``Ở đâu chả được mà\~{}''

``Không được! Tôi muốn đứng cạnh bà cơ!'' nàng bờm Chó kéo tay cô bạn, nhắng nhít chạy về phía đằng sau.

``Ừ rồi, rồi\dots\ Đành chiều bà vậy\~{}'' nàng bờm Mèo mệt mỏi đáp lại.

Ở một bên khác, Yuki đứng chống hông, nhìn quanh như đang đánh giá bố cục.
``Đứng lệch sang bên trái một chút đi. Nếu đứng thế kia thì ánh sáng chiếu vào không đều đâu.''

Shiina khoanh tay, đứng yên một chỗ, mặt vẫn giữ vẻ lạnh lùng như mọi ngày: ``Tôi thì đứng đâu cũng được. Xấu đẹp gì cũng chẳng quan trọng.''

Trong lúc đó, ba anh em Hikawa đứng cạnh nhau một cách tự nhiên. Kuon ở giữa, Kaigou đứng bên phải, còn Suzuki đứng bên trái.

Suzuki khẽ mỉm cười, ôm chặt lấy hai anh chị: ``Hiếm khi chúng ta đứng thế này nhỉ. Lần cuối chắc là hồi chúng ta ở thế giới Cầu Vồng ha?''

Kaigou liếc sang Kuon:``Cậu đừng có cử động linh tinh đấy.''

Kuon thở dài, giọng như nói đểu: ``Tôi đứng yên từ nãy đến giờ rồi đấy thôi. Câu đó tôi phải nói với cô mới đúng.''

``Cái cậu này!'' cô nàng mắt xanh đấm nhẹ vào vai của cậu, bật cười khe khẽ.

``Được rồi, mọi người ổn định nhé!'' Yuko giơ tay lên cao. ``3\dots\ 2\dots\ 1\dots''

*TÁCH*

Khoảnh khắc ấy được giữ lại. Một tấm ảnh, nơi tất cả cùng cười, cùng chen chúc, cùng tồn tại trong một khoảnh khắc tưởng như rất bình thường.

Không khí trở nên nhẹ nhàng hơn, không còn sự căng thẳng của thi đấu. Những câu chuyện bắt đầu râm ran trở lại: về những cú ném bóng, về pha kéo co ``bay người'', về màn cổ vũ suýt làm Kaigou đứng tim.

``Anh chị thấy không!? Em suýt rơi thật đấy!'' Suzuki vừa nói vừa cười.

Kaigou nhíu mày, tỏ vẻ không hài lòng: ``Không phải suýt đâu. Là nguy hiểm thật. Ngộ nhỡ em bị thương thì sao?''

``Nhưng em tiếp đất ổn mà.''

``\dots Lần sau không được làm vậy nữa.''

Kuon đứng bên cạnh, chỉ khẽ lắc đầu: ``Đúng là không có trong kịch bản thật. À mà này\dots'' Cậu quay về phía cô chị lớn, nhướng mày hỏi: ``Cô chạy nhanh thật đấy. Kinh nghiệm nào vậy?''

Kaigou khoanh tay, đáp tỉnh queo: ``Từ thế giới cũ.''

Cậu Tổng không mất quá nhiều thời gian để suy đoán: ``Chắc là chạy trốn khỏi đám con trai dâm dục chứ gì?''

Cô nàng liếc xéo, không chút đỏ mặt: ``Không sai.''

Suzuki nghe vậy thì bật cười, vỗ vai hai người: ``Hội thao năm nay em thích nhất là được thấy cả lớp gắn bó với nhau.''

Kuon gật gù, mắt vẫn nhìn theo dòng người đang tan dần: ``Anh chưa từng thấy hội thao nào sôi động đến thế.''

Kaigou nhếch mép: ``Ừ, đúng dịp để cậu tập thể dục. Chắc là từ ngày mai chúng ta sẽ dậy tập thể dục buổi sáng thôi nhỉ?''

Kuon đỏ ửng tai, lẩm bẩm: ``Thôi, xin kiếu.''

Ở một góc khác, Koishin khoanh tay, mặt hơi đỏ: ``Thắng thì thắng rồi đấy\dots\ nhưng đừng có mà tưởng bở.''

Kanade cười khúc khích: ``Cậu vui mà, đúng không?''

``Không có!''

Yae thì đã bắt đầu bàn sang chuyện\dots\ hội thao năm sau: ``Năm sau mình phải thử thêm môn bóng rổ nữa. Lần này chưa đủ đã.''

``Cậu vẫn còn sức à\dots'' một người nào đó trong lớp thở dài.

Ở đằng sau, Ayame và Wata xách dụng cụ, đi song hành với nhau. Nàng tóc trắng-đỏ xoa vai, nhìn cô bạn thân cười nhẹ: ``Năm nay mình đã thấy đủ kiểu người rồi, kẻ liều mạng, kẻ an toàn, kẻ lén lút. Đúng là một bản giao hưởng đầy màu sắc.''

Nàng tóc vàng gật, giọng trầm ấm: ``Và điều khiến mình nhớ nhất không phải chiến thắng, mà là cách cả trường cùng nhịp tim khi cờ đỏ vừa phất. Cảm giác ấy\dots\ thực sự khó tái hiện lắm.''

Ayame nhếch môi, đôi mắt lim dim như đang giữ lại chút nắng cuối: ``Năm sau, mình sẽ cố gắng để trở thành một phát ngôn viên bùng cháy nhất!''

Wata khép bảng tổng sắp, mỉm cười kín đáo: ``Mong rằng lúc đó, chúng ta vẫn còn đủ nhiệt để cùng nhau đỏ một lần nữa.''

Yuko nhìn cả lớp, ánh mắt dịu lại: ``Lớp mình làm tốt lắm. Không chỉ là kết quả\dots\ mà là cách các em phối hợp với nhau.''

Không ai nói gì thêm, nhưng ai cũng hiểu.

Hội thao không chỉ là thắng thua.

Mà là việc họ đã thực sự trở thành một tập thể.

\begin{center}
  \uline{Ngày hôm sau\\Phòng Câu lạc bộ Nhiếp ảnh.}
\end{center}

Thứ Hai, khi hơi nắng đầu tuần còn lạt lẻo như bàn tay vừa rời khỏi chăn sau một ngày hội thao đầy nhiệt huyết. Phòng CLB Nhiếp ảnh nằm chót về phía cuối, cửa sắt xanh đã mở hé, để lộ khung cảnh bên trong: màn sáng xám của buổi sớm lùa qua khe cửa sổ, loang lổ trên bàn chồng lens, dây cáp, và những chiếc cốc cà phê còn nguyên vệt môi đêm qua. Mùi giấy ẩm, mùi nhựa máy ảnh nóng, và chút mồ hôi khô lại trên áo đội tình nguyện hòa quyện thành ``mùi hội thao'' chưa kịp tan.

Trong góc phòng, màn hình máy tính đang bật sáng, chiếu ra thứ ánh sáng trắng xanh lạnh như băng cuối hè. Uzuki ngồi trước bàn, cằm tựa lên mu bàn tay, mắt dán vào từng khung hình đang chạy chậm. Dưới ánh đèn vi tính, gương mặt cô trông hơi tái, nhưng đôi mắt lại sáng lên mỗi khi thấy góc quay ổn định, cú pan mượt như bơ hoặc pha phóng to thu nhỏ nối tiếp không vỡ khung.

Kuon ngồi cạnh, tựa tay vào ghế mà cô nàng bờm Thỏ đang ngồi. Mái tóc cậu hơi rối, có lẽ vì cả đêm cắt ghép, cũng có thể vì lúc quay cậu đã lăn lộn trên khán đài cùng đám cổ động viên. Trên màn hình, đoạn thanh thời gian màu xanh lá hiện đầy những đoạn cắt nhỏ: tamaire, kéo co, bóng né, đấu vật\dots\ mỗi môn một màu, như bản đồ nhiệt của cảm xúc.

``Cậu quay chắc tay thật đấy,'' Uzuki nói, giọng khàn vì ngủ ít, nhưng vẫn đủ để cất lên thành lời khen chân thành. Cô tua lại một đoạn bóng né, cho chạy chậm đến mức có thể thấy từng giọt mồ hôi trên trán Koishin khi cậu lách qua quả bóng. ``Nhìn cái pan này đi, mượt như có ray dolly ấy.''

Kuon nhún vai, đáp nhỏ: ``Chỉ là quen tay thôi.'' Giọng cậu trầm, hơi khàn, như thể đang cố giấu đi vẻ mệt mỏi lẫn chút tự hào.

``Không, không phải kiểu 'quen' bình thường đâu,'' Uzuki dừng lại, chỉ vào đoạn chuyển động camera theo hình zíc-zắc giữa hai đội kéo co. ``Góc này\dots\ tốc độ này\dots\ cậu phải từng làm chuyên nghiệp mới có thể cảm được nhịp như vậy.''

Kuon im lặng. Đôi mắt cậu vẫn dán vào màn hình, nhưng đầu ngón tay khẽ gõ nhịp lên thành bàn, như đang đếm khung hình trong đầu. Mỗi lần chuyển cảnh, ánh sáng trên mặt hai người lại thay đổi, như bị kéo theo cảm xúc của chính những người trong khung hình.

``Dù vậy,'' cô nàng nghiêng đầu, kéo con trỏ về một đoạn rung nhẹ cuối cùng của bài cổ vũ, ``cũng có vài khúc bị giật. Chắc do cậu mải hò hét theo khán đài, phải không?''

Kuon khẽ cười nhạt, khóe môi hếch lên một chút: ``Có lẽ vậy'' Cậu dừng lại một giây, như đang nhớ lại cảm giác khi cả khối Đỏ đồng thanh reo tên Kaigou, khiến cậu quên mất tay đang giữ máy. ``Không tránh được mà.''

``Ừ, cũng dễ hiểu thôi,'' Uzuki gật gù, ngồi soát lại những đoạn phim do chính cô và cậu Tổng quay lại.

Ngoài cửa sổ, tiếng chuông báo vào tiết đầu tiên vang lên, như nhắc họ rằng cuộc sống học sinh vẫn tiếp diễn, nhưng ít nhất thì trong căn phòng nhỏ này, họ vừa kết thúc một bộ phim của riêng mình, bằng những thước phim rung lắc, mồ hôi, và cả trái tim nồng nhiệt của tuổi trẻ.

Uzuki kéo thanh thời gian sang phải, định xem nốt đoạn bóng rổ, nhưng ngón tay cô chợt khựng lại giữa chừng. Màn hình vẫn chạy, nhưng mắt cô không theo nữa; có thứ gì đó ở mép ngoài khung hình vừa lướt qua, mờ như bụi đất chiều, đủ khiến cô thót tim.

``\dots Khoan đã.'' Tiếng cô nhỏ, như sợ ai nghe thấy.

Kuon ngồi thẳng lưng, đôi mắt vốn buồn ngủ bỗng đăm đăm. ``Sao vậy?''

Cô nàng bờm Thỏ không đáp. Cô rút tay khỏi chuột, kéo con trỏ về phía trước, tua lại từng khung hình chậm rãi. Máy chiếu rè rè, ánh sáng xanh lạnh loang trên mặt hai người như màn đêm thủy tinh. Khung hình dừng lại ở phút thứ mười bảy của hiệp hai bóng rổ tại góc xa khán đài, nơi hàng ghế trống lác đác. Ở đó, một bóng dáng nhỏ xíu đứng lẻ loi, tựa vào lan can. Không cổ vũ, không nhảy, không reo. Chỉ đứng.

``Cái này\dots\ là gì vậy?'' cô khẽ nghiêng đầu, mắt dán vào màn hình như sợ bóng kia tan biến nếu cô chớp mắt.

Cậu con trai híp mắt, lòng bàn tay tựa lên cằm. ``Phóng to lên đi.''

Uzuki nhấp chuột hai lần. Khung hình nổ tung thành từng điểm ảnh, nhưng bóng kia vẫn rõ: một cô bé tóc đen tết đuôi sam, dài chạm gối, mảnh mai như nhánh liễu non. Đôi mắt xanh lục sâu hoắm, nhìn vào ống kính máy ảnh, khẽ lửng lơ trên không. Môi cô bé khép, tay buông thõng, đứng yên như tượng gỗ giữa dòng người đang cuồng nhiệt.

Cả hai im bặt. Trong phòng chỉ còn tiếng quạt máy rít nhẹ, và tiếng tim họ đập rộn ràng. ``Có phải\dots\ học sinh không?'' Uzuki lên tiếng trước, giọng khàn như thể đang cố nuốt một hòn đá.

Kuon không trả lời. Cậu kéo chuột, nhảy sang đoạn video tiếp theo quay lại trận chạy tiếp sức. Góc quay rộng hơn, hướng về đích. Và lại\dots\ cô bé ấy. Lần này cô đứng dưới gốc cây bàng, tay chạm vào thân cây như đếm từng vết sần. Mắt vẫn xanh lục, vẫn không nhìn máy quay, vẫn không chớp.

Một bóng hình nhỏ nhắn, lặng lẽ hiện ra ở khung cửa sau khán đài. Không gian khác, thời điểm khác, nhưng vẫn là mái tóc tết đuôi sam đen nhánh, vẫn là đôi mắt xanh lục không chớp.

Uzuki nghiêng người về phía trước, màn hình khẽ rung theo nhịp thở gấp, mặt hơi tái nhợt đi, cảm tưởng như có dòng điện chạy dọc lưng.

Kuon không rời mắt khỏi màn hình, ánh nhìn như bị dính chặt vào từng khung hình. Giọng cậu trầm xuống: ``Mở thử ảnh đi.''

Uzuki gật nhẹ, tay lần vào thư mục. Bức ảnh kỷ niệm lớp 1-C hiện lên, đầy màu sắc và tiếng cười. Cả lớp chen chúc, ai cũng háo hức. Nhưng rồi, gần chỗ ba anh em Hikawa, một bóng hình nhỏ nhắn lẫn vào như một mảnh ghép vô hình. Cô bé tóc tết đen, đôi mắt xanh lục, đứng im lìm, không ai để ý, không ai nhìn, nhưng lại hiện diện rõ ràng.Cô bé đứng thẳng, không chạm vào ai, không được ai chạm, nhưng lại vừa khít vào bố cục như thể đã được chỉnh sửa từ ngoài đời thật.

Uzuki nuốt khan, cổ họng như khô đắng. Giọng cô khẽ run: ``\dots Cậu cũng thấy rồi, đúng không?''

Kuon không nói, chỉ gật nhẹ. Cánh tay cậu tựa lên thành ghế, ngón tay thôi gõ nhịp, như thể cả thế giới vừa ngừng lại sau một dấu lặng.

Bầu không khí trong phòng như đặc quánh lại gần như ngay tức khắc. Uzuki cảm thấy gáy mình tê dại, cô vô thức thu người về phía trước, mắt vẫn dán vào màn hình như sợ bóng kia sẽ biến mất nếu cô chớp mắt. Cô bé tóc tết đen nhánh, đôi mắt xanh lục không học sinh Aogami nào có, hiện ra rõ đến mức cô có thể đếm từng sợi tóc, nhưng lại mờ đến mức cô không thể chắc đó là thật hay chỉ là bóng ma của thuật toán nén ảnh.

Kuon ngồi im như tượng, chỉ có đôi mắt vẫn theo dõi từng pixel. Cậu không run, không giật mình, chỉ khẽ siết chặt cằm như đang cố giữ một bí mật đã cất giấu từ lâu. Cậu biết cô bé ấy là ai, hoặc ít nhất, biết cô không thuộc về ngày hội thao này, nhưng cậu không nói. Chỉ có hơi thở dài nhẹ, gần như vô thức, thoát ra từ kẽ môi cậu như một lời thú tội không thành lời.

``Mình\dots\ liệu mình có nhìn nhầm không? Hay là do mình thiếu ngủ nên sinh ra hoang tưởng\dots?'' Uzuki lên tiếng, giọng khản đặc. Cô cố gắng dùng từ kỹ thuật để trấn an mình, nhưng ngón tay run rẩy khi kéo thanh thời gian qua lại, để rồi lại thấy cô bé ấy ở góc sân bóng rổ, dưới tán cây bàng, sau khán đài, như một bản sao bị dính vào mọi khung hình, không mời mà đến, không mời mà ở lại.

Kuon không trả lời. Cậu chỉ im lặng, như thể im lặng là cách duy nhất để không kéo cô bé ấy ra khỏi màn hình và vào đời thực.

Ở đâu đó, ngoài khung cửa sổ, gió khẽ đẩy cánh màn trắng rung rinh. Trong căn phòng chỉ còn tiếng quạt máy rít nhẹ, và tiếng tim ai đó đập rộn ràng như muốn thoát khỏi lồng ngực.

Cô bé tóc tết vẫn đứng đó.
Không cười.
Không chớp mắt.
Chỉ đứng đó quan sát - hoặc đang chờ đợi - một điều gì đó mà chỉ mình cô bé mới biết.

\end{document}